
\documentclass[11pt]{article}
\usepackage[utf8]{inputenc}
\usepackage[T1]{fontenc}
\usepackage[margin=0.88in]{geometry}
\usepackage{amsmath,amssymb,amsthm,mathtools}
\usepackage{enumitem}
\usepackage{booktabs,longtable,array,tabularx}
\usepackage{xcolor}
\usepackage[hidelinks,bookmarks=false]{hyperref}
\usepackage{microtype}
\usepackage{titlesec}
\usepackage{fancyhdr}
\usepackage{stmaryrd}
\usepackage{makecell}

\setlist[itemize]{leftmargin=2.0em,itemsep=0.14em,topsep=0.15em}
\setlist[enumerate]{leftmargin=2.2em,itemsep=0.14em,topsep=0.15em}
\renewcommand{\arraystretch}{1.13}
\newcolumntype{L}[1]{>{\raggedright\arraybackslash}p{#1}}
\setlength{\parskip}{0.52em}
\setlength{\parindent}{0pt}
\setlength{\headheight}{14pt}
\pagestyle{fancy}
\fancyhf{}
\lhead{RH Prime-Trace Amplitude Separator}
\rhead{Full Audit-Grade Manuscript}
\cfoot{\thepage}

\titleformat{\section}{\Large\bfseries}{\thesection}{0.7em}{}
\titleformat{\subsection}{\large\bfseries}{\thesubsection}{0.7em}{}
\titleformat{\subsubsection}{\normalsize\bfseries}{\thesubsubsection}{0.7em}{}

\newtheorem{definition}{Definition}[section]
\newtheorem{theorem}{Theorem}[section]
\newtheorem{lemma}{Lemma}[section]
\newtheorem{proposition}{Proposition}[section]
\newtheorem{corollary}{Corollary}[section]
\newtheorem{remark}{Remark}[section]
\newtheorem{audit}{Audit}[section]
\newtheorem{nonclaim}{Nonclaim}[section]

\DeclareMathOperator{\Ree}{Re}
\DeclareMathOperator{\ImPart}{Im}
\DeclareMathOperator{\supp}{supp}

\newcommand{\Az}{A_{\zeta}}
\newcommand{\Lcrit}{L_{\mathrm{crit}}}
\newcommand{\RH}{\mathrm{RH}}
\newcommand{\EF}{\mathrm{EF}}
\newcommand{\EndpointUse}{\mathrm{EndpointUse}}

\newcommand{\BareCx}{\mathrm{BareCounterexampleContent}}
\newcommand{\OrdCxThm}{\mathrm{OrdCounterexampleThm}}
\newcommand{\OfficialCxForce}{\mathrm{OfficialCounterexampleForce}}
\newcommand{\LocalCxForce}{\mathrm{LocalCounterexampleForce}}
\newcommand{\LocalEndpointUse}{\mathrm{LocalEndpointUse}}
\newcommand{\RHOffCxForce}{\mathrm{RHOffLineCounterexampleForce}}
\newcommand{\CxForce}{\mathrm{CounterexampleForce}}
\newcommand{\EndpointCounterexampleSlot}{\mathrm{EndpointCounterexampleSlot}}
\newcommand{\FixedDomain}{\mathrm{FixedDomain}}
\newcommand{\Aplus}{A^+}
\newcommand{\ATS}{\mathrm{ATS}}
\newcommand{\UEAP}{\mathrm{UEAP}}
\newcommand{\KR}{K_R}
\newcommand{\CRHEF}{C_{\mathrm{RH}\text{-}\mathrm{EF}}}
\newcommand{\ModelEF}{\mathrm{Model}_{\mathrm{EF}}}
\newcommand{\AdeqEF}{\mathrm{Adeq}_{\mathrm{EF}}}
\newcommand{\OfficialEFRoute}{\mathrm{OfficialEFRoute}}
\newcommand{\OSEvalEF}{\mathrm{OSEval}_{\mathrm{RH}\text{-}\mathrm{EF}}}
\newcommand{\CritAmpOcc}{\mathrm{CritAmpOcc}}
\newcommand{\CritAmpImageOcc}{\mathrm{CritAmpImageOcc}}
\newcommand{\OffAmpExcl}{\mathrm{OffLineAmpExcl}}
\newcommand{\AmpSepOcc}{\mathrm{AmpSep}^{\mathrm{occ}}}
\newcommand{\AmpSepLoc}{\mathrm{AmpSep}^{\mathrm{loc}}}
\newcommand{\NonCritAmpWork}{\mathrm{NonCriticalAmpStatusWork}}
\newcommand{\Damp}{D_{\mathrm{amp}}}
\newcommand{\Namp}{N_{\mathrm{amp}}}
\newcommand{\ProjRoleEF}{\mathrm{ProjectionRoleContent}_{\mathrm{EF}}}
\newcommand{\Counterforce}{\mathrm{EndpointCounterforce}}
\newcommand{\OrdOfficialRes}{\mathrm{OrdOfficialRes}}
\newcommand{\EndpointAdmRel}{\mathrm{EndpointAdmRel}}
\newcommand{\SameDomIncomp}{\mathrm{SameDomainIncomp}}
\newcommand{\EndpointForce}{\mathrm{EndpointForce}}
\newcommand{\EndpointInert}{\mathrm{EndpointInert}}
\newcommand{\CarrierShift}{\mathrm{CarrierShift}}
\newcommand{\Repair}{\mathrm{Repair}}
\newcommand{\Book}{\mathrm{Bookkeeping}}
\newcommand{\Support}{\mathrm{SupportOnly}}
\newcommand{\Surrogate}{\mathrm{Surrogate}}
\newcommand{\Skin}{\mathrm{Skin}}
\newcommand{\Tensor}{\mathrm{Tensor}}
\newcommand{\Anchor}{\mathrm{Anchor}}
\newcommand{\Amp}{\mathrm{Amp}}
\newcommand{\psierr}{E_{\psi}}
\newcommand{\psinorm}{\mathcal E_{\psi}}
\newcommand{\AmpStand}{\mathrm{EndpointStandingNonNeutralAmp}}
\newcommand{\NonNativeDisc}{\mathrm{NonNativeAmpDisc}}
\newcommand{\SameZero}{\mathrm{SameZero}}
\newcommand{\PrimeTraceAdeq}{\mathrm{PrimeTraceEndpointAdeq}}
\newcommand{\ordhalf}{\operatorname{ord}_{1/2}}
\newcommand{\Dord}{D_{\mathrm{ord}}}
\newcommand{\NativeFifth}{\mathrm{NativeFifthRole}}
\newcommand{\NativeCounterforce}{\mathrm{NativeCounterforce}}
\newcommand{\BridgeNeutralized}{\mathrm{BridgeNeutralized}}
\newcommand{\EndpointStanding}{\mathrm{EndpointStanding}}
\newcommand{\BzetaEF}{B^{\mathrm{EF}}_{\zeta}}
\newcommand{\BridgeCompleteEF}{\mathrm{BridgeComplete}_{\mathrm{EF}}}
\newcommand{\BridgeFailEF}{\mathrm{BridgeFail}_{\mathrm{EF}}}
\newcommand{\SameZeroEF}{\mathrm{SameZero}_{\mathrm{EF}}}
\newcommand{\DomainShiftEF}{\mathrm{DomainShift}_{\mathrm{EF}}}
\newcommand{\Czeta}{C_{\zeta}}
\newcommand{\Kgov}{K_{\mathrm{gov}}}
\newcommand{\BareCounterContent}{\mathrm{BareCounterexampleContent}}
\newcommand{\OrdCounterThm}{\mathrm{OrdCounterexampleThm}}
\newcommand{\OfficialCounterforce}{\mathrm{OfficialCounterexampleForce}}
\newcommand{\LocalCounterforce}{\mathrm{LocalCounterexampleForce}}
% duplicate removed: \newcommand{\LocalEndpointUse}{\mathrm{LocalEndpointUse}}
\newcommand{\RHOffCounterforce}{\mathrm{RHOffLineCounterexampleForce}}


\begin{document}

\begin{center}
{\LARGE \bfseries The Riemann Hypothesis by Exclusion of the Prime-Trace Amplitude Separator}\\[0.35em]
{\Large A Full Audit-Grade AASC Endpoint-Rigidity Proof via Explicit-Formula Scale Discrimination}\\[0.8em]
{\large Amos Jay Maley}\\[0.2em]
{\large June 2026}\\[0.8em]
\end{center}

\begin{abstract}
This manuscript gives a full audit-grade AASC endpoint-exclusion proof of the Riemann Hypothesis through the prime-trace amplitude route.  It does not introduce a new zero-free region, zero-density estimate, positivity criterion, Hilbert--Polya operator, random-matrix theorem, or numerical verification.  It uses the standard explicit-formula zero-to-prime-trace correspondence as the official endpoint route on the fixed zeta/prime-trace carrier.

Let
\[
\Az(\rho) :\Longleftrightarrow \zeta(\rho)=0\ \wedge\ 0<\Ree(\rho)<1
\]
be nontrivial zeta zerohood, and let
\[
\Lcrit(\rho) :\Longleftrightarrow \Ree(\rho)=\frac12
\]
be critical-line readout.  The positive endpoint is
\[
\RH^+ := \forall \rho\,(\Az(\rho)\Rightarrow \Lcrit(\rho)).
\]
The bare negative branch is the ordinary off-line zerohood assertion
\[
\RH^-_{\mathrm{bare}} := \exists \rho\,(\Az(\rho)\wedge \neg\Lcrit(\rho)).
\]
Descriptive off-line coordinate reports are preserved.  The contradiction below is not derived from \(\Ree(\rho)\ne 1/2\) alone.

\textbf{Front-loaded dependency lock.} This manuscript has one nonclassical dependency: official endpoint-defeating counterexample force is theorem-bearing endpoint use.  Once that is granted, the remaining proof is target instantiation plus route exhaustion.  The dependency is not hidden in the explicit formula, in the zeta carrier, in the critical-normalized amplitude interface, or in the local reductio rule.  Bare negative content may exist as syntax, hypothesis, support, or ordinary theoremhood; but when it is used to defeat the official endpoint on the same fixed carrier, that use fixes target, carrier, branch, reportability, downstream reuse, positive-branch exclusion, act-time finality, and lawful-redescription stability.  That is endpoint use.  A reader who denies this dependency should attack the no-autonomous-counterexample theorem directly; the rest of the manuscript is the RH-EF instantiation of that dependency.

The decisive transport is through the explicit formula and the structural-to-explicit-formula equivalence bridge.  The structural RH role decomposition supplies analytic zerohood, carrier preservation, kernel-governed report identity, carrier-preservation bridge, and critical-line readout.  The present manuscript instantiates that bridge as \(\BzetaEF\), the explicit-formula amplitude bridge, which does not contain line support by definition.  It binds same-zero zerohood to an explicit-formula contribution, line readout to neutral critical-normalized order, and off-line readout to non-neutral order.  For a zero \(\rho=\beta+i\gamma\), the zero contribution carries scale \(x^{\rho}/\rho\) in a standard Chebyshev prime trace.  Under critical normalization by \(x^{1/2}\), its amplitude profile is governed by \(x^{\beta-1/2}\).  Critical-line zeros occupy the neutral amplitude interface \(x^0\).  An off-line zero supplies a non-neutral scale displacement.  When this displacement is used as same-carrier endpoint counterforce against the official RH endpoint through the explicit-formula route, it is not merely a native coordinate report; it is endpoint-governing prime-trace amplitude separator occupation.  Native-counterforce exhaustion then proves that ordinary native off-line amplitude falsification is not a fifth endpoint role: endpoint-standing non-neutral order is bridge-neutralized, support-only, domain/interface shift, or independent amplitude discriminator work.  In the live off-line countercase, bridge-neutralization is unavailable, support-only use lacks endpoint counterforce, and fixed-domain evaluation excludes domain/interface shift; the remaining branch is \(\Damp(\rho)\).

The proof uses the AASC kernel, UEAP report discipline, and ATS anchor--tensor--skin decomposition as structural locking support.  UEAP prevents report laundering from coordinate or support status into endpoint status.  ATS fixes the zeta/prime-trace anchors, identifies the amplitude exponent as tensor content rather than notation skin, and blocks surrogate or normalization skins from carrying endpoint force.  The AMetric and no-independent-discriminator consequences of endpoint use exclude independent same-domain amplitude endpoint discriminators.  The live exact-complement local countercase \([\neg\Lcrit(\rho)]_i\) induces local amplitude separator occupation; that occupation induces \(D_{\mathrm{amp}}(\rho)\), while the same endpoint use yields \(N_{\mathrm{amp}}(\rho)\Rightarrow \neg D_{\mathrm{amp}}(\rho)\).  The local countercase discharges and gives \(\Lcrit(\rho)\) for arbitrary nontrivial zerohood.  Universal generalization gives the Riemann Hypothesis.
\end{abstract}

\textbf{Mathematical-status lock.} AASC is used here as a mathematical constraint formalism.  The proof is a theorem chain about endpoint use, target instantiation, explicit-formula amplitude projection, role occupation, and no-independent-discriminator closure.  It is not philosophical commentary or an optional interpretive overlay.

\textbf{Proof-class lock.} The proof class is \emph{kernel-forced prime-trace amplitude separator exclusion}.  Wrong-mode demands for a zero-free region, density estimate, spectral operator, positivity criterion, or numerical extension do not touch the proof class unless they identify a failure in target instantiation, model adequacy, amplitude projection, separator induction, or no-independent closure.

\textbf{Corpus-support boundary.} The AASC Corpus Control Matrix is used only to select structural support anchors for kernel forcing, UEAP report discipline, ATS role decomposition, AMetric no-selector discipline, role-occupancy closure, no post-selection repair, silent-redescription exclusion, and same-domain incompleteness control.  Those corpus anchors do not supply analytic zero-location facts.  They are summarized in Appendix~\ref{app:corpus}.

\textbf{Formalization boundary.} This manuscript is the mathematical proof draft.  A future Lean audit would formalize the endpoint-role surface: explicit-formula route objects, projection-role content, amplitude separator occupation, discriminator induction, no-independent closure, and reductio discharge.  No Lean formalization is claimed here.

\section*{Front-Loaded Dependency Lock}
\addcontentsline{toc}{section}{Front-Loaded Dependency Lock}

This manuscript has one nonclassical dependency: official endpoint-defeating counterexample force is theorem-bearing endpoint use.  Once that is granted, the remaining proof is target instantiation plus route exhaustion.  This sentence is the dependency lock for the whole paper.  The proof does not ask the reader to accept a new analytic estimate, a hidden positivity principle, a disguised zero-free region, or a redefinition of the zeta carrier.  It asks the reader to accept the endpoint-use thesis for official counterexample force, then verifies the RH-EF target instantiation and exhausts the possible roles of the projected off-line amplitude branch.

The lock has two sides.  First, raw counterexample content is not endpoint use: an object \(\rho\), the formula \(\Az(\rho)\wedge\neg\Lcrit(\rho)\), and the non-neutral order statement \(\ordhalf(\rho)\ne0\) may remain syntax, hypothesis, descriptive analytic content, support, or ordinary theoremhood.  Second, endpoint-defeating counterexample force is not merely an object: it is the mathematical act of offering that object or theorem as the same-carrier defeat of the official endpoint.  Such an act fixes the target, carrier, branch, exclusion of the positive endpoint, reportability, downstream reuse, act-time status, and lawful-redescription role.  Those are exactly the endpoint-use features governed by the kernel and the A+ consequence layer.

After this dependency is stated, the rest of the proof has no hidden nonclassical step.  The RH target is instantiated through the explicit-formula route; the structural RH roles are bridged to critical-normalized prime-trace order; native off-line counterforce is exhausted into bridge-neutralized, support-only, bookkeeping, domain/interface-shift, or independent-amplitude-discriminator roles; the live fixed-domain subproof eliminates the first four; and A+ excludes the remaining discriminator.  Therefore a same-mode objection should not scatter across the manuscript.  It should either deny the endpoint-use status of official counterexample force, or identify a failure in target instantiation, role projection, native-counterforce exhaustion, local discharge, or A+ closure.

\tableofcontents
\newpage

\section*{Read This First: Endpoint Proof Spine}
\addcontentsline{toc}{section}{Read This First: Endpoint Proof Spine}

Let \(\rho\in\mathbb C\) be arbitrary, and write \(\rho=\beta+i\gamma\).  The proof works locally under the zerohood assumption \(\Az(\rho)\).  The raw local split is
\[
\Az(\rho)\Rightarrow \bigl(\Lcrit(\rho)\vee \neg\Lcrit(\rho)\bigr).
\]
Open the exact-complement local countercase
\[
[\neg\Lcrit(\rho)]_i.
\]
Under the official explicit-formula endpoint route, \(\Az(\rho)\wedge\neg\Lcrit(\rho)\) projects to non-neutral critical-normalized prime-trace amplitude:
\[
\Az(\rho)\wedge \beta\ne\frac12
\Rightarrow
\ordhalf(\rho)=\beta-\frac12\ne0
\Rightarrow
x^{\beta-1/2}\ne x^0.
\]
The proof explicitly does \emph{not} use
\[
x^{\beta-1/2}\ne x^0\Rightarrow \bot.
\]
Only after endpoint-countercase use is fixed, and only after the structural-to-EF bridge classifies the branch as non-neutral order on the fixed prime-trace endpoint image, does the native-counterforce collapse apply.  The native off-line branch is not granted a fifth role:
\[
\NativeCounterforce_{\RH\text{-}\EF}(\rho)
\Rightarrow
\BridgeNeutralized(\rho)\vee\Support(\rho)\vee\DomainShiftEF(\rho)\vee\Damp(\rho).
\]
Inside the exact-complement off-line subproof, bridge-neutralization contradicts the live non-neutral order, support-only status lacks endpoint counterforce, and fixed-domain evaluation excludes domain/interface shift.  Thus the only endpoint-standing native-counterforce branch is \(\Damp(\rho)\).  Equivalently, the projected role becomes local separator occupation:
\[
\ProjRoleEF(\neg\Lcrit(\rho),\rho)
\wedge \Counterforce_{\RH\text{-}\EF}(\rho)
\Rightarrow \AmpSepLoc(\rho).
\]
Then
\[
\AmpSepLoc(\rho)
\Rightarrow \NonCritAmpWork_R(\rho)
\Rightarrow \Damp(\rho).
\]
Endpoint use of the same local separator act yields
\[
\AmpSepLoc(\rho)
\Rightarrow \EndpointUse_R(\AmpSepLoc(\rho))
\Rightarrow \Namp(\rho)
\Rightarrow \neg\Damp(\rho).
\]
Thus \(\Damp(\rho)\wedge\neg\Damp(\rho)\), contradiction.  The annotation-discharge theorem discharges the original local assumption \([\neg\Lcrit(\rho)]_i\), not a strengthened conjunction.  Therefore \(\Lcrit(\rho)\).  Since \(\rho\) was arbitrary under \(\Az(\rho)\),
\[
\forall \rho\,(\Az(\rho)\Rightarrow \Lcrit(\rho)).
\]
This is the classical Riemann Hypothesis.

\textbf{Anti-trivialization lock.} The manuscript does not license the schema \(\neg P\Rightarrow \mathrm{Separator}(P)\Rightarrow P\).  Endpoint exclusion is permitted only after target instantiation, official route fixation, projection-role content, endpoint-counterforce use, route-locus exhaustion, and the proof that the projected role performs independent same-domain endpoint-status work.

\textbf{Counterexample-force lock.} A counterexample object is not automatically endpoint use.  This manuscript distinguishes raw counterexample content, ordinary counterexample theoremhood, official counterexample force, and local exact-complement countercase force.  Only official counterexample force and local endpoint countercase force enter endpoint-use discipline.  Thus the proof does not classify the object \(\rho\) as endpoint use; it classifies the mathematical act of offering \(\rho\), or a theorem asserting its existence, as endpoint-defeating counterexample force against the official RH endpoint.

\textbf{No-autonomous-negative-endpoint lock.} The proof's final same-mode commitment is that ordinary negative endpoint force is not a free-standing fifth role.  If the negative branch has no endpoint force, it is raw content, support, bookkeeping, or ordinary theoremhood.  If it has endpoint force against the official target on the same carrier, it is endpoint use and is governed by the same kernel-forced consequence layer as any official endpoint act.

\newpage

\section{Counterexample Force as Endpoint Use}\label{sec:counterexample-force-endpoint-use}

This section supplies the final general bridge needed by the native-counterforce layer. A classical reader may agree that the off-line branch projects to non-neutral explicit-formula order and still object that a counterexample is simply a mathematical object satisfying the negation of the endpoint, not endpoint governance. The response is a role distinction. A bare object or proposition is not an endpoint act. The mathematical use of that object or proposition to defeat, displace, or settle the official endpoint on the same fixed carrier is an endpoint act.

The section is deliberately placed before target instantiation. Later RH-specific constructions rely on the following general discipline:
\[
\OfficialCxForce_R(B)\Rightarrow\EndpointUse_R(B).
\]
It is the counterexample analogue of ordinary official endpoint resolution. Ordinary theoremhood is allowed. Raw counterexample syntax is allowed. Local reductio assumptions are allowed. The controlled object is endpoint-defeating counterexample force.

\subsection{Bare content, theoremhood, official force, and local force}

\begin{definition}[Bare counterexample content]
For a fixed endpoint context \(R\) and positive endpoint branch \(P_R\), \(\BareCx_R(B)\) means that \(B\) is raw logical counterexample content whose truth would contradict or refute the positive endpoint proposition on the relevant carrier. It records proposition-level content only. It does not assert proof, reportability, endpoint standing, downstream reuse, official resolution, or endpoint use.

For RH-EF, the local bare counterexample content attached to a point \(\rho\) is
\[
\BareCx_{\RH\text{-}\EF}(\rho):=\Az(\rho)\wedge\neg\Lcrit(\rho).
\]
The global bare negative branch is
\[
\RH^-_{\mathrm{bare}}:=\exists \rho\,\BareCx_{\RH\text{-}\EF}(\rho).
\]
\end{definition}

\begin{definition}[Ordinary counterexample theoremhood]
\(\OrdCxThm_R(B)\) means ordinary theoremhood of the counterexample branch \(B\):
\[
\OrdCxThm_R(B):=\mathrm{Proof}_R(B).
\]
For RH-EF, this includes a hypothetical ordinary proof of
\[
\exists \rho\,\bigl(\Az(\rho)\wedge\neg\Lcrit(\rho)\bigr).
\]
Ordinary counterexample theoremhood is proof legitimacy. It is not, by itself, the independent amplitude discriminator excluded later.
\end{definition}

\begin{definition}[Official counterexample force]
\(\OfficialCxForce_R(B)\) means that a counterexample branch \(B\), or a theorem-bearing act asserting \(B\), is used to defeat, displace, or settle the official endpoint on the same fixed carrier. It includes the following endpoint-use features:
\begin{enumerate}
\item the official target is fixed;
\item the carrier of evaluation is fixed;
\item the branch being used as counterexample is fixed;
\item the act is theorem-bearing, or is locally assumption-bearing inside an exact-complement subproof;
\item the positive endpoint branch is excluded or blocked by the counterexample use;
\item the counterexample is reportable as endpoint-defeating content if the act succeeds;
\item the counterexample licenses downstream reuse as negative endpoint material if accepted;
\item the act is final at act-time rather than repaired into standing later as the same act;
\item lawful redescription preserves target, carrier, branch, and endpoint-defeating role.
\end{enumerate}
\end{definition}

\begin{definition}[Local counterexample force]
For a proof target \(P\) and a branch \(B\), write
\[
\LocalCxForce_R(B;P).
\]
This means that \(B\) has been opened inside a reductio or proof-by-cases subproof as the live exact-complement countercase whose survival would block \(P\) on the fixed endpoint carrier. It is local endpoint-facing countercase use. It is not global official negative endpoint resolution and not theoremhood of \(B\).

For RH-EF, the local form is
\[
\LocalCxForce_{\RH\text{-}\EF}(\neg\Lcrit(\rho);\Lcrit(\rho)).
\]
\end{definition}

\begin{definition}[Local endpoint use]
\(\LocalEndpointUse_R(B)\) means endpoint use inside a governed local subproof rather than a global endpoint outcome. It fixes the same target, carrier, branch role, local report force, act-time status, and admissible continuation inside the subproof. Local endpoint use is enough to trigger local fixed-domain consequences. It does not assert global official negative resolution.
\end{definition}

\subsection{Bare content is not endpoint use}

\begin{theorem}[Bare counterexample content is not endpoint use]
Bare counterexample content alone does not imply endpoint use:
\[
\BareCx_R(B)\not\Rightarrow \EndpointUse_R(B).
\]
In the RH-EF instance,
\[
\Az(\rho)\wedge\neg\Lcrit(\rho)\not\Rightarrow
\EndpointUse_R(\neg\Lcrit(\rho)).
\]
\end{theorem}

\begin{proof}
Bare content is a proposition-level role. It may be true, false, hypothetical, descriptive, support-level, unproved, or later proved. None of those statuses by itself fixes official endpoint target, reportability, downstream reuse, branch exclusion, act-time finality, or endpoint-role stability. Endpoint use requires an act that offers the branch as endpoint-defeating or endpoint-resolving. Therefore bare counterexample content is not endpoint use.
\end{proof}

\begin{audit}[Objects are not acts]
The manuscript does not classify the object \(\rho\) as endpoint use. It classifies a mathematical act: using \(\rho\), or a theorem asserting its existence, as counterexample force against the official endpoint. An object can be raw content. A counterexample act has endpoint role.
\end{audit}

\begin{theorem}[Ordinary counterexample theoremhood is kernel-internal, not automatically forbidden]
For any fixed endpoint context,
\[
\OrdCxThm_R(B)\Rightarrow K_R.
\]
But ordinary counterexample theoremhood, considered only as theoremhood, is not the endpoint discriminator later excluded:
\[
\OrdCxThm_R(B)\not\Rightarrow D(B).
\]
\end{theorem}

\begin{proof}
A genuine proof is a governed derivation and therefore kernel-internal. Its polarity does not make it defective. The endpoint discriminator arises only when the counterexample theorem is used as official endpoint counterforce or endpoint resolution on the fixed carrier and thereby determines branch standing, reportability, downstream reuse, or exclusion of the positive branch. Ordinary theoremhood alone does not perform that role.
\end{proof}

\begin{audit}[Negative theoremhood is protected]
A hypothetical ordinary proof of an off-line zero is not rejected merely because it is negative. The paper distinguishes proof legitimacy from endpoint-role occupation. The rejected object, if the endpoint route is fixed, is endpoint-standing amplitude separator authority after official counterexample use.
\end{audit}

\subsection{Official counterexample force is endpoint use}

\begin{theorem}[Official counterexample force is endpoint use]\label{thm:official-counterexample-force-endpoint-use}
For every fixed endpoint context \(R\) and counterexample branch \(B\),
\[
\OfficialCxForce_R(B)\Rightarrow \EndpointUse_R(B).
\]
\end{theorem}

\begin{proof}
A counterexample used only as raw syntax or background possibility is not official counterexample force. By definition, official counterexample force uses the branch to defeat or settle the fixed endpoint on the same carrier. Such an act fixes the target being defeated, fixes the carrier on which defeat is claimed, fixes the negative branch as counterexample branch, fixes the exclusion of the positive branch, and fixes reportability and downstream reuse if the act succeeds. It must be final at act-time, since a later bridge or repair would be a new endpoint act rather than the same counterexample already standing. It is stable under lawful redescription, since otherwise the alleged counterexample would not be a counterexample to the same endpoint. These are exactly the features of theorem-bearing endpoint use. Therefore \(\OfficialCxForce_R(B)\Rightarrow\EndpointUse_R(B)\).
\end{proof}


\begin{corollary}[No autonomous counterexample endpoint role]\label{cor:no-autonomous-counterexample-endpoint-role}
There is no same-carrier official counterexample force that stands outside endpoint-use discipline:
\[
\neg\exists B\,\bigl(\OfficialCxForce_R(B)\wedge \neg\EndpointUse_R(B)\bigr).
\]
Equivalently,
\[
\neg\EndpointUse_R(B)\Rightarrow \neg\OfficialCxForce_R(B).
\]
\end{corollary}

\begin{proof}
The first display is the direct negation of the possibility excluded by Theorem~\ref{thm:official-counterexample-force-endpoint-use}.  That theorem says: whenever a branch has official counterexample force, it has endpoint-use status.  Hence no branch can have official counterexample force while lacking endpoint-use status.  The second display is the contrapositive.
\end{proof}

\begin{theorem}[Autonomous ordinary-negative endpoint role collapse]\label{thm:autonomous-negative-role-collapse}
Let \(B\) be a negative or counterexample branch on a fixed endpoint carrier.  If \(B\) is not endpoint use, then it cannot defeat the official endpoint as the settled same-carrier counterexample branch.  Its remaining lawful roles are raw content, ordinary theoremhood considered only as proof legitimacy, support, bookkeeping, local syntax without endpoint force, or carrier/domain-shifted material.
\end{theorem}

\begin{proof}
Suppose \(B\) is not endpoint use.  By Corollary~\ref{cor:no-autonomous-counterexample-endpoint-role}, \(B\) has no official counterexample force.  Therefore \(B\) cannot simultaneously exclude the positive endpoint branch, fix negative endpoint reportability, license downstream reuse as the settled negative endpoint, and preserve the same carrier and act-time finality.  If the branch is present at all, it remains in a role that does not settle the endpoint: raw syntax, proof legitimacy without official endpoint use, support, bookkeeping, local assumption without endpoint force, or a shifted carrier/domain.  These cases exhaust the ordinary negative roles that do not exercise official endpoint-defeating force.
\end{proof}

\begin{audit}[No autonomous ordinary-negative endpoint role]
This is the explicit form of the manuscript's central nonclassical commitment.  The proof does not deny ordinary negative content.  It denies that ordinary negative content can retain official endpoint-defeating force while refusing endpoint-use status.  A same-mode objection must therefore exhibit a counterexample act that defeats the official endpoint while failing to fix the endpoint-use features listed in Definition~2.3.
\end{audit}

\begin{remark}[Counterexample use, not counterexample object]
The theorem does not say that every object satisfying a negative predicate is an endpoint act. It says that the act of offering such an object, or a theorem asserting its existence, as same-carrier official counterexample force is endpoint use. This is the same role distinction used throughout the manuscript between support-level content and endpoint occupation.
\end{remark}

\begin{theorem}[Official counterexample force satisfies target adequacy]\label{thm:counterexample-force-target-adequacy}
If official counterexample force is exercised on a fixed domain, then the kernel and A+ consequence layer are in force:
\[
\OfficialCxForce_R(B)\wedge\FixedDomain(R)
\Rightarrow K_R\wedge\Aplus_R(B).
\]
\end{theorem}

\begin{proof}
By Theorem~\ref{thm:official-counterexample-force-endpoint-use}, official counterexample force is endpoint use. Endpoint use is governed construction: it fixes target, carrier, branch, theorem-status or local countercase force, downstream reuse, act-time finality, and redescription-stable role. The target-adequacy-to-kernel theorem gives \(K_R\). Fixed-domain endpoint use then triggers the A+ consequence package: AMetric no-selector discipline, unique admissible interior, report preservation, no carrier transfer, no same-act repair, no third endpoint status, and no independent same-domain endpoint discriminator. Hence \(K_R\wedge\Aplus_R(B)\).
\end{proof}

\begin{audit}[Why counterexample force is governed]
A counterexample with no fixed target is not a counterexample to this endpoint. A counterexample with no fixed carrier may be a different problem. A counterexample with no reportability or downstream reuse is support or syntax. A counterexample whose standing is supplied only later is a repaired act. Thus a same-carrier official counterexample act has already accepted the endpoint-use features whose consequences the proof applies.
\end{audit}

\subsection{Local exact-complement counterexample force}

\begin{theorem}[Exact-complement local force is local endpoint use]\label{thm:local-counterexample-force-endpoint-use}
In a fixed endpoint-closure proof, if \(B\) is opened as the live exact-complement countercase to target \(P\) and is used with endpoint counterforce, then it is local endpoint use:
\[
\LocalCxForce_R(B;P)\wedge \Counterforce_R(B)
\Rightarrow \LocalEndpointUse_R(B).
\]
For the RH-EF local subproof,
\[
\begin{aligned}
&\LocalCxForce_{\RH\text{-}\EF}(\neg\Lcrit(\rho);\Lcrit(\rho))\\
&\quad\wedge\Counterforce_{\RH\text{-}\EF}(\rho)\\
&\quad\Rightarrow \LocalEndpointUse_R(\neg\Lcrit(\rho)).
\end{aligned}
\]
\end{theorem}

\begin{proof}
Local counterexample force is the local proof role in which the exact complement is introduced to block the endpoint target. Endpoint counterforce says that the branch is not inert syntax but is being used to defeat the target inside the subproof. Together they fix the local target, the local branch role, the same carrier, the exclusion relation to the positive target, and the act-time subproof status. Those are the local analogues of endpoint use. The conclusion is local endpoint use, not global negative endpoint outcome.
\end{proof}

\begin{theorem}[Local endpoint use triggers local A+ closure]\label{thm:local-endpoint-use-aplus}
On a fixed domain,
\[
\LocalEndpointUse_R(B)\wedge\FixedDomain(R)
\Rightarrow \Aplus_{R,\mathrm{loc}}(B).
\]
For the RH-EF amplitude separator case, local A+ includes the no-independent-amplitude-discriminator closure:
\[
\Aplus_{R,\mathrm{loc}}(B)
\Rightarrow N_{\mathrm{amp}}(\rho).
\]
\end{theorem}

\begin{proof}
The kernel roles do not suspend inside a reductio subproof. A local endpoint act still has target, carrier, branch role, standing-bearing local force, admissible continuation, and act-time status. Therefore fixed-domain closure supplies the local A+ consequences. In the RH-EF setting, the relevant local consequence is the no-independent-amplitude-discriminator closure \(\Namp(\rho)\).
\end{proof}

\begin{audit}[No global promotion]
The local A+ theorem does not promote the local assumption to global official negative resolution. It only states that a governed local endpoint countercase on the same fixed carrier cannot use a local independent discriminator while the kernel-forced consequence layer is in force.
\end{audit}

\subsection{RH-EF specialization}

\begin{definition}[RH off-line counterexample force]
\(\RHOffCxForce(\rho)\) means that the bare analytic content
\[
\Az(\rho)\wedge\neg\Lcrit(\rho)
\]
is used to defeat the official Riemann Hypothesis endpoint on the fixed zeta/explicit-formula carrier. It is stronger than off-line coordinate description and stronger than support-level amplitude displacement.
\end{definition}

\begin{theorem}[RH off-line counterexample force is endpoint use]
\[
\RHOffCxForce(\rho)\Rightarrow \EndpointUse_R(\RH^-_{\mathrm{bare}}).
\]
\end{theorem}

\begin{proof}
RH off-line counterexample force is the RH-EF instance of official counterexample force. It uses the off-line zerohood branch to defeat the official RH endpoint on the same fixed carrier. The general counterexample-force endpoint-use theorem therefore applies.
\end{proof}

\begin{theorem}[RH off-line counterexample force projects to EF role content]
Under the official explicit-formula route,
\[
\begin{aligned}
&\RHOffCxForce(\rho)\wedge\OfficialEFRoute(R,\zeta,\psi)\\
&\quad\Rightarrow \ProjRoleEF(\neg\Lcrit(\rho),\rho).
\end{aligned}
\]
\end{theorem}

\begin{proof}
The official EF route binds branch roles to explicit-formula amplitude roles. Off-line counterexample force is not merely coordinate description; it is the endpoint-defeating use of the off-line branch. Under projection, that branch has non-neutral EF order as its endpoint-role content. This is precisely \(\ProjRoleEF(\neg\Lcrit(\rho),\rho)\).
\end{proof}

\begin{theorem}[RH off-line counterexample force enters native-counterforce exhaustion]
\[
\begin{aligned}
&\RHOffCxForce(\rho)\wedge\OfficialEFRoute(R,\zeta,\psi)\\
&\quad\wedge\ModelEF(\zeta,\psi,R)\\
&\quad\Rightarrow \NativeCounterforce_{\RH\text{-}\EF}(\rho).
\end{aligned}
\]
\end{theorem}

\begin{proof}
By the preceding theorem, the off-line counterexample act has projected EF role content. The model adequacy theorem identifies that content with non-neutral critical-normalized order for the same zerohood instance. Since the use is endpoint-defeating counterexample force rather than descriptive or support-only use, it is native endpoint counterforce in the sense used by the native-counterforce collapse theorem.
\end{proof}

\begin{corollary}[No separate ordinary-counterexample escape]
For RH-EF, official off-line counterexample force has no endpoint role outside the native-counterforce exhaustion:
\[
\begin{aligned}
\RHOffCxForce(\rho)\Rightarrow{}&\BridgeNeutralized(\rho)\vee\Support(\rho)\vee\Book(\rho)\\
&\vee\DomainShiftEF(\rho)\vee\Damp(\rho).
\end{aligned}
\]
Inside the live off-line endpoint countercase on a fixed domain, the first four branches are unavailable, so \(\Damp(\rho)\) remains.
\end{corollary}

\begin{proof}
The first statement follows from entry into native-counterforce exhaustion. In the live local off-line subproof, bridge-neutralization gives endpoint-neutral order and hence conflicts with \(\neg\Lcrit(\rho)\); support-only and bookkeeping lack endpoint counterforce; fixed-domain endpoint evaluation excludes domain/interface shift. Therefore the remaining exhaustive endpoint role is \(\Damp(\rho)\).
\end{proof}

\begin{audit}[Where the classical objection now lands]
After this section, a classical objection cannot merely say that a counterexample is an object rather than governance. The manuscript agrees. The object is raw content. The endpoint-defeating use of that object is an endpoint act. A same-mode denial must now show a same-carrier official counterexample act that defeats the endpoint while failing to fix target, carrier, branch, reportability, downstream reuse, branch exclusion, act-time finality, or redescription-stable role.
\end{audit}

\newpage

\section{Official Endpoint, Target Instantiation, and Context Non-Smuggling}

\subsection{Official target and branch split}

\begin{definition}[Nontrivial zerohood and critical-line readout]
For a complex number \(\rho\), define
\[
\Az(\rho):=\bigl(\zeta(\rho)=0\bigr)\wedge \bigl(0<\Ree(\rho)<1\bigr),
\]
and
\[
\Lcrit(\rho):=\bigl(\Ree(\rho)=1/2\bigr).
\]
The official positive endpoint is
\[
\RH^+ := \forall \rho\, (\Az(\rho)\Rightarrow \Lcrit(\rho)).
\]
The bare negative branch is
\[
\RH^-_{\mathrm{bare}}:=\exists\rho\,\bigl(\Az(\rho)\wedge\neg\Lcrit(\rho)\bigr).
\]
\end{definition}

\begin{lemma}[Raw complementarity]
At the bare analytic proposition layer,
\[
\neg \RH^+\quad\Longleftrightarrow\quad \RH^-_{\mathrm{bare}}.
\]
Locally, under \(\Az(\rho)\), the exact complement of \(\Lcrit(\rho)\) is \(\neg\Lcrit(\rho)\), equivalently \(\Ree(\rho)\ne 1/2\).
\end{lemma}

\begin{proof}
The first equivalence is classical quantifier negation applied to \(\forall \rho(\Az(\rho)\Rightarrow \Lcrit(\rho))\).  The local equivalence follows from the definition of \(\Lcrit\).  No endpoint standing, reportability, or theorem status is inferred from this bare complementarity.
\end{proof}

\begin{audit}[Raw complementarity is not endpoint standing]
The proof uses raw complementarity only to open a local reductio branch and later discharge it.  It does not infer that the negative branch has official endpoint standing.  The standing-bearing role is supplied only by exact-complement countercase use in the fixed endpoint-closure proof, and then only locally.
\end{audit}

\subsection{The RH-EF context object}

\begin{definition}[Fixed RH explicit-formula endpoint context]
The object
\[
\CRHEF(R,\zeta,\psi)
\]
records the official RH endpoint context through the explicit-formula route.  It fixes the zeta carrier, the prime-trace carrier, nontrivial zerohood, critical-line readout, the positive endpoint branch, the bare negative branch, the critical-normalized amplitude image, the local countercase slot, the endpoint counterexample slot, and the amplitude separator occupation slot.
\end{definition}

\begin{theorem}[Context non-smuggling]
The context object \(\CRHEF(R,\zeta,\psi)\) does not assert any of the following:
\[
\RH^+,
\quad \Lcrit(\rho),
\quad \neg\RH^-_{\mathrm{bare}},
\quad \Aplus_R(B),
\quad \Namp(\rho),
\quad \neg\Damp(\rho),
\quad \neg\AmpSepLoc(\rho),
\quad \RHOffCxForce(\rho).
\]
It fixes carrier, route, branch slots, endpoint-counterexample slots, and endpoint-image roles only.
\end{theorem}

\begin{proof}
By definition the context object records what is being evaluated and by which route.  The zeta carrier does not include line support; the prime-trace carrier does not assert neutral amplitude occupation for all zeros; the route predicate does not assert the endpoint.  All status consequences are derived later from endpoint use, model adequacy, projection, and the kernel-forced consequence layer.
\end{proof}

\begin{audit}[Context smuggling objection]
A hostile objection must identify a clause inside \(\CRHEF\) that already states line membership, neutral amplitude occupation for all zeros, no-independent closure, or exclusion of the off-line branch.  Merely fixing the critical-line endpoint image is not smuggling; it is target instantiation.  The positive image is fixed, not asserted as universally occupied.
\end{audit}

\subsection{Official explicit-formula endpoint route}

\begin{definition}[Official explicit-formula route]
\(\OfficialEFRoute(R,\zeta,\psi)\) means that the official RH endpoint is being evaluated through the standard explicit-formula relation between zeta zerohood and prime-trace oscillation, while preserving the same zeta zerohood carrier, the same prime-trace carrier, and the same critical-line readout target.  It is stronger than mentioning the explicit formula as background evidence.
\end{definition}

\begin{definition}[Endpoint-closure evaluation]
\(\OSEvalEF(R,\zeta,\psi)\) means that the official RH endpoint is being evaluated in the endpoint-closure proof class on the fixed explicit-formula carrier.  It fixes target, carrier, route, branch slots, and local counterforce discipline.  It is not endpoint-complete evaluation and does not assert either branch.
\end{definition}

\begin{theorem}[Official route is target instantiation, not endpoint assertion]
\[
\CRHEF(R,\zeta,\psi)\wedge \OfficialEFRoute(R,\zeta,\psi)
\not\Rightarrow \RH^+.
\]
The route supplies the bridge by which branch roles are read; it does not supply line support.
\end{theorem}

\begin{proof}
The official route says how endpoint roles are projected onto the prime-trace amplitude image.  Projection is not occupation.  In particular, a route can map the positive branch to neutral amplitude and the negative branch to non-neutral amplitude without deciding which branch is true.  Therefore no positive endpoint follows from route fixation alone.
\end{proof}

\section{Kernel, A+, ATS, and UEAP Replay}

\subsection{Kernel-neutral target adequacy}

\begin{definition}[Kernel-neutral target adequacy]
A theorem-bearing endpoint act is target-adequate when it satisfies: target determinacy, step evaluability, act-time finality, and same-regime fidelity.  These conditions are neutral with respect to AASC vocabulary; they are the conditions under which a raw trace can be assessed as a determinate same-carrier mathematical act.
\end{definition}

\begin{theorem}[Target adequacy forces the kernel]
If an endpoint act is target-adequate, then reference, standing, admissibility, and irreversibility are already in force for that act:
\[
TargetAdequacy_R(B)\Rightarrow Ref_R\wedge St_R\wedge Adm_R\wedge Irr_R.
\]
\end{theorem}

\begin{proof}
Target determinacy supplies fixed reference.  Step evaluability supplies standing-bearing step status.  Act-time finality supplies irreversibility: later repair is a new act, not the same act made valid.  Same-regime fidelity supplies the admissibility boundary preserving those statuses under lawful redescription.  Thus the kernel roles are forced roles, not optional AASC additions.
\end{proof}

\begin{theorem}[Endpoint use is governed construction]
For any fixed endpoint context \(R\) and branch or local endpoint act \(B\),
\[
\EndpointUse_R(B)\Rightarrow GovernedConstruction_R(B)\Rightarrow K_R.
\]
\end{theorem}

\begin{proof}
Endpoint use fixes target, carrier, theorem-status or countercase status, downstream reuse, act-time finality, and redescription-stable role.  These are exactly the features of governed construction on a fixed domain.  The kernel-forcing theorem applies.
\end{proof}

\subsection{A+ consequence package}

\begin{definition}[A+ for the RH-EF endpoint]
\(\Aplus_R(B)\) is the fixed-domain consequence layer forced by kernel-instantiated endpoint use.  For the RH-EF route it includes boundary fixation, endpoint bivalence, AMetric no-selector discipline, standing as reuse-stable admissibility, conservation of standing, unique admissible interior, no raw trace sufficiency, no carrier transfer, no post-selection repair, no silent redescription, no third endpoint status, report preservation, ATS anchor/tensor/skin stability, and no independent same-domain endpoint-status discriminator.
\end{definition}

\begin{theorem}[Kernel forces A+]
\[
\EndpointUse_R(B)\wedge \FixedDomain(R)\Rightarrow \Aplus_R(B).
\]
\end{theorem}

\begin{proof}
By endpoint use the kernel is in force.  Fixed-domain closure applies the kernel consequences: bivalence and AMetricity at the boundary, uniqueness of the admissible interior, standing conservation, no raw-trace sufficiency, no carrier transfer, no same-act repair, no skin authority, and no independent same-domain discriminator.  A+ is therefore a consequence layer rather than an independent premise.
\end{proof}

\begin{audit}[A+ is not doing analytic estimation]
A+ does not assert a numerical bound for \(\psi(x)-x\), does not prove a zero-free region, and does not alter the explicit formula.  It controls whether a theorem-bearing endpoint act may introduce a second same-domain endpoint-status classifier.  The analytic transport to amplitude status is supplied separately by the explicit-formula route.
\end{audit}

\subsection{ATS local packet}

\begin{definition}[Local ATS packet]
\(\ATS_{\RH\text{-}\EF}\) is the conjunction of: fixed zeta anchor, fixed prime-trace anchor, same-zero identity, fixed critical normalization, explicit-formula route preservation, tensor status of the exponent \(\beta-1/2\), skin-inertness of presentation choices, and no second same-domain amplitude classifier.
\end{definition}

\begin{theorem}[ATS blocks skin authority]
Under \(\ATS_{\RH\text{-}\EF}\), a change of notation, smoothing presentation, truncation convention, or equivalent formula display does not create endpoint standing.  If it changes the same-zero contribution, the normalized exponent, or the prime-trace carrier, it is not skin but a carrier or tensor change requiring bridge discharge.
\end{theorem}

\begin{proof}
ATS separates anchors from tensor content and from skin.  Endpoint status can be carried only by anchor-preserving tensor content.  A skin change carries no standing.  A non-skin change alters the target or carrier and therefore cannot be used as same-domain endpoint resolution without fresh bridge certification.
\end{proof}

\subsection{UEAP local packet}

\begin{definition}[UEAP report ladder]
The RH-EF report ladder is:
\begin{enumerate}
\item zerohood report \(\Az(\rho)\);
\item coordinate report \(\beta\ne 1/2\);
\item explicit-formula contribution report \(x^\rho/\rho\);
\item support-level amplitude report \(x^{\beta-1/2}\ne x^0\);
\item endpoint-counterforce report;
\item endpoint-result report.
\end{enumerate}
\end{definition}

\begin{theorem}[UEAP no-overreporting for amplitude]
Support-level non-neutral amplitude does not by itself become endpoint counterforce:
\[
OffLineAmpSupport(\rho)\not\Rightarrow \Counterforce_{\RH\text{-}\EF}(\rho).
\]
Endpoint counterforce requires official route fixation, local exact-complement use, and role occupation.
\end{theorem}

\begin{proof}
UEAP forbids a report from exceeding its carrier, bridge, and support.  The amplitude report says what scale is contributed by a zero under the explicit formula.  It does not by itself assert endpoint resolution.  Endpoint counterforce is a higher report role: the amplitude displacement is used to defeat the critical-neutral endpoint image on the same carrier.  That use must be separately established.
\end{proof}

\subsection{Official endpoint resolution target adequacy}

\begin{theorem}[Official RH-EF endpoint resolution satisfies target adequacy]
For any branch or local countercase act \(B\),
\[
\OrdOfficialRes_R(B)\wedge \OfficialEFRoute(R,\zeta,\psi)\wedge \ModelEF(\zeta,\psi,R)\wedge \FixedDomain(R)
\Rightarrow K_R\wedge \Aplus_R(B).
\]
\end{theorem}

\begin{proof}
Target determinacy holds because the endpoint is the official RH assertion, the carrier is the fixed zeta/prime-trace carrier, the route is the explicit formula, and the endpoint image is critical-normalized amplitude.  Step evaluability holds because theorem-bearing use classifies a branch as endpoint outcome, support, bookkeeping, or counterforce.  Act-time finality holds because later analytic certification creates a new act rather than repairing an earlier unbridged endpoint report.  Same-regime fidelity holds because lawful redescription preserves the zeta carrier, prime-trace carrier, same-zero identity, normalization, and branch role.  Target adequacy forces the kernel, and the kernel on a fixed domain forces A+.
\end{proof}

\section{Explicit-Formula Prime-Trace Model Adequacy}

\subsection{Prime-trace carrier and critical normalization}

Let
\[
\psi(x)=\sum_{n\le x}\Lambda(n),\qquad \psierr(x)=\psi(x)-x,
\]
and
\[
\psinorm(x)=\frac{\psi(x)-x}{x^{1/2}}.
\]
The normalization by \(x^{1/2}\) is not an assertion of RH.  It is the critical-line endpoint normalization: on the positive endpoint image, a zero contribution has neutral magnitude exponent.

\begin{definition}[Explicit-formula contribution role]
For a nontrivial zero \(\rho=\beta+i\gamma\), the explicit-formula contribution role is represented, suppressing smoothing and harmless presentation choices, by a term of scale
\[
\frac{x^\rho}{\rho}.
\]
After critical normalization, the amplitude profile is governed by
\[
\frac{x^{\rho-1/2}}{\rho},
\]
and its magnitude exponent is \(\beta-1/2\).
\end{definition}

\begin{definition}[Context-bound explicit-formula model adequacy]
\(\ModelEF(\zeta,\psi,R)\) is the adequate explicit-formula endpoint model when the following clauses hold:
\begin{enumerate}[label=(EF\arabic*)]
\item zeta-zero coverage: every nontrivial zerohood instance in the official RH carrier appears in the zero contribution class;
\item contribution soundness: every contribution used by the model is attached to the same zeta carrier;
\item same-zero identity preservation: the contribution attributed to \(\rho\) remains tied to the same zerohood instance;
\item symmetry-pair awareness: paired and conjugate contributions are lawful but do not erase individual scale-class typing unless a bridge proves endpoint-neutral cancellation;
\item smoothing/truncation discipline: test functions and truncations are skin or support unless they change endpoint role;
\item amplitude exactness: \(\rho=\beta+i\gamma\) has critical-normalized amplitude scale \(x^{\beta-1/2}\);
\item critical-neutral exactness: \(\beta=1/2\) iff the normalized amplitude exponent is zero;
\item off-line displacement exactness: \(\beta\ne 1/2\) iff the normalized amplitude exponent is nonzero;
\item no RH smuggling: no clause asserts that all zero contributions are neutral;
\item support/endpoint separation: amplitude data are lawful support unless endpoint counterforce use is supplied.
\end{enumerate}
\end{definition}

\begin{theorem}[Model adequacy]
\[
\CRHEF(R,\zeta,\psi)\wedge \FixedDomain(R)\wedge \OfficialEFRoute(R,\zeta,\psi)
\Rightarrow \ModelEF(\zeta,\psi,R).
\]
\end{theorem}

\begin{proof}
The context fixes the official zeta and prime-trace carriers.  The route fixes the explicit formula as the zero-to-prime-trace projection.  The zero contribution of \(\rho\) has the standard scale \(x^\rho/\rho\); critical normalization divides by \(x^{1/2}\), leaving exponent \(\beta-1/2\).  Symmetry, smoothing, truncation, and equivalent formula presentations are included as lawful carrier data but do not assert universal neutral occupation.  Therefore the adequacy clauses hold.
\end{proof}

\begin{audit}[Model adequacy burden]
A valid model-adequacy objection must identify a nontrivial zero lost by the contribution model, a contribution not tied to the zeta carrier, a failure of same-zero identity, a normalization mismatch, a hidden RH assumption, or a support/endpoint collapse.  Merely objecting that the explicit formula contains error terms or smoothing choices does not meet the burden unless those choices alter endpoint role.
\end{audit}

\subsection{Term isolation and pairing closure}

\begin{definition}[Term-isolation discipline]
A zero contribution may be treated as endpoint-relevant amplitude content only when same-zero identity, explicit-formula route, normalization, and contribution class are fixed.  Term isolation is not a claim of analytic dominance, absence of cancellation, or global asymptotic lower bound.
\end{definition}

\begin{theorem}[Term-isolation discipline]
If \(\ModelEF(\zeta,\psi,R)\) holds and \(\Az(\rho)\), then the contribution associated with \(\rho=\beta+i\gamma\) has a well-defined critical-normalized scale class \(\beta-1/2\).  This scale class may be used in endpoint-role analysis without asserting that the single term dominates the entire prime trace.
\end{theorem}

\begin{proof}
Model adequacy fixes same-zero identity and the explicit-formula contribution class.  The scale class is an algebraic feature of the contribution after normalization.  Dominance and cancellation are separate analytic questions.  The endpoint analysis requires only that off-line zerohood projects to non-neutral scale class, not that this class dominates all other terms.
\end{proof}

\begin{theorem}[Pairing closure]
Suppose \(\rho=\beta+i\gamma\) is off line.  The symmetry-related family of zero contributions is lawful analytic content.  However, pairing does not change the individual scale-class fact \(\beta-1/2\ne 0\) unless a bridge proves that the paired family is endpoint-neutral in the same critical-normalized image.  Therefore symmetry pairing is support or bridge material, not automatic neutral occupation.
\end{theorem}

\begin{proof}
The functional equation and conjugation symmetries generate related zeros and paired terms.  Those terms may combine in oscillatory or real-valued ways.  But the endpoint image under audit is critical-neutral amplitude occupation for the same zerohood contribution class.  Pairing can participate in a proof that the endpoint image is neutral only if it provides a bridge from the paired contribution family to neutral occupation.  Without that bridge, it is lawful analytic structure, not erasure of non-neutral scale-class typing.
\end{proof}

\begin{audit}[Pairing/cancellation non-overclaim]
The theorem does not say off-line contributions dominate the prime trace.  It says cancellation does not license endpoint-neutral occupation unless proved as endpoint-neutral bridge content.  This blocks the hostile objection that the paper ignores paired sums while avoiding the stronger analytic claim that no cancellation can occur.
\end{audit}

\section{Branch Projection and the Single Amplitude Endpoint Interface}

\subsection{Positive and negative amplitude projections}

\begin{definition}[Critical-neutral amplitude occupation]
For a fixed zerohood instance \(\rho=\beta+i\gamma\), define \(\CritAmpOcc(\rho)\) by
\[
\CritAmpOcc(\rho) :\Longleftrightarrow x^{\beta-1/2}=x^0
\]
as critical-normalized amplitude role occupation.  Equivalently, under model adequacy, \(\CritAmpOcc(\rho)\) is the amplitude image of \(\Lcrit(\rho)\).
\end{definition}

\begin{definition}[Off-line amplitude exclusion]
\[
\OffAmpExcl(\rho):\Longleftrightarrow x^{\beta-1/2}\ne x^0.
\]
This is non-occupation of the critical-neutral amplitude image.  It is support-level analytic content unless used as endpoint counterforce.
\end{definition}

\begin{theorem}[Explicit-formula branch projection]
Under \(\OfficialEFRoute(R,\zeta,\psi)\wedge \ModelEF(\zeta,\psi,R)\), for each zerohood instance \(\rho\),
\[
\Lcrit(\rho)\Longleftrightarrow \CritAmpOcc(\rho),
\]
and
\[
\neg\Lcrit(\rho)\Longleftrightarrow \OffAmpExcl(\rho).
\]
\end{theorem}

\begin{proof}
Write \(\rho=\beta+i\gamma\).  By definition \(\Lcrit(\rho)\) is \(\beta=1/2\).  By amplitude exactness, critical-normalized scale has exponent \(\beta-1/2\).  Thus \(\beta=1/2\) exactly when the scale is \(x^0\), and \(\beta\ne 1/2\) exactly when the scale is non-neutral.  This proves both equivalences.
\end{proof}

\begin{audit}[Projection is not truth selection]
The projection theorem does not decide which branch holds.  It tells us what the branch means on the explicit-formula amplitude image if that branch is used in the official route.  This mirrors target instantiation: branch roles are fixed before branch truth is decided.
\end{audit}

\subsection{Projection-role content}

\begin{definition}[Projection-role content]
\(\ProjRoleEF(B,\rho)\) means that, under the official explicit-formula route, the endpoint-use act \(B\) has the explicit-formula amplitude role associated with \(\rho\) as its endpoint-role content.  For the local negative branch, this content is \(\OffAmpExcl(\rho)\).
\end{definition}

\begin{theorem}[Local off-line countercase binds to projection-role content]
\[
\Az(\rho)\wedge \OfficialEFRoute(R,\zeta,\psi)\wedge \ModelEF(\zeta,\psi,R)
\wedge \mathrm{ReductioCountercase}_{\RH\text{-}\EF}(\neg\Lcrit(\rho);\Lcrit(\rho))
\Rightarrow \ProjRoleEF(\neg\Lcrit(\rho),\rho).
\]
\end{theorem}

\begin{proof}
The local assumption is the exact complement of critical-line readout for the fixed zerohood instance.  The official route requires that branch content be read through the explicit-formula amplitude image.  By branch projection, the complement of line readout is non-neutral amplitude exclusion.  The projection-role predicate records that this is the endpoint-role content of the local countercase, not a free-standing descriptive assertion.
\end{proof}

\begin{audit}[Projection-role content is not yet contradiction]
Projection-role content is still not forbidden.  The act becomes separator occupation only when the projected role is used as endpoint counterforce and is not support, bookkeeping, carrier shift, repair, surrogate substitution, or coequal occupation.
\end{audit}

\subsection{Single amplitude interface}

\begin{theorem}[Single amplitude endpoint interface]
On the fixed RH-EF carrier, critical-neutral amplitude occupation has exactly one occupation interface:
\[
\CritAmpImageOcc(\rho)\Longleftrightarrow \CritAmpOcc(\rho)\Longleftrightarrow \Lcrit(\rho).
\]
Thus the positive endpoint image is neutral critical-normalized amplitude occupation, not a family of coequal neutral and non-neutral occupants.
\end{theorem}

\begin{proof}
The endpoint image is the explicit-formula projection of critical-line readout.  Critical-line readout is \(\beta=1/2\), which is exactly normalized exponent zero.  A nonzero normalized exponent is non-occupation of this image.  Treating non-neutral scale as coequal occupation changes from the critical-neutral endpoint image to a broader symmetric outcome carrier.  It is not the same endpoint image.
\end{proof}

\begin{corollary}[Off-line amplitude is non-occupation]
\[
\OffAmpExcl(\rho)\Rightarrow \neg \CritAmpImageOcc(\rho).
\]
\end{corollary}

\begin{proof}
Off-line amplitude exclusion means non-neutral exponent; the single-interface theorem says occupation requires neutral exponent.  Hence non-occupation follows.
\end{proof}

\begin{theorem}[Single-interface non-explosiveness]
The manuscript does not infer
\[
\OffAmpExcl(\rho)\Rightarrow\bot.
\]
It permits non-neutral amplitude as descriptive analytic consequence, support-level hypothetical content, and ordinary counterexample syntax.  Contradiction arises only from endpoint-standing use of that non-neutral amplitude as same-carrier separator occupation.
\end{theorem}

\begin{proof}
The proof later adds endpoint-counterforce use, non-support status, no carrier shift, no repair, no surrogate substitution, and no coequal occupation.  Without those additional hypotheses, non-neutral amplitude is a lawful analytic description of what an off-line zero contribution would look like.
\end{proof}


\section{Structural-to-Explicit-Formula Equivalence Bridge}\label{sec:structural-ef-bridge}

This section is the new load-bearing bridge from the structural RH role-compression paper into the explicit-formula amplitude carrier.  Its purpose is to prevent the amplitude proof from relying on a loose analogy.  The structural RH paper separates analytic carrier, zerohood standing, critical-line readout, and carrier-preservation bridge.  The present manuscript instantiates that bridge through the explicit-formula amplitude route.  The result is not the assertion that off-line coordinates are illegal.  The result is a typed equivalence between structural line readout and critical-normalized amplitude order.

\subsection{Structural roles and explicit-formula roles}

\begin{definition}[Structural RH role packet]
For a fixed nontrivial zerohood instance \(\rho\), the structural RH packet consists of
\[
\Az(\rho),\qquad \Czeta(\rho),\qquad \Kgov(\rho),\qquad B_\zeta(\rho),\qquad \Lcrit(\rho).
\]
Here \(\Az(\rho)\) is analytic zerohood in the critical strip, \(\Czeta(\rho)\) is carrier preservation in the fixed zeta carrier, \(\Kgov(\rho)\) is kernel-governed same-scope report standing, \(B_\zeta(\rho)\) is the carrier-preservation bridge, and \(\Lcrit(\rho)\) is critical-line readout.
\end{definition}

\begin{definition}[Explicit-formula bridge object]
\(\BzetaEF(\rho)\) is the explicit-formula instantiation of the structural carrier-preservation bridge.  It records a route from zeta zerohood standing to critical-normalized amplitude readout on the same carrier.  It requires:
\begin{enumerate}[label=(B\arabic*)]
\item same-zero preservation: the contribution is attached to the same \(\rho\), not a symmetry partner, averaged cloud, finite height class, or model analogue;
\item route preservation: the branch is read through \(\OfficialEFRoute(R,\zeta,\psi)\);
\item contribution preservation: the explicit-formula term class associated with \(\rho\) is the carrier of the amplitude report;
\item normalization preservation: the critical normalization by \(x^{1/2}\) is fixed;
\item bridge certification: any pairing, smoothing, truncation, or cancellation used to neutralize a non-neutral contribution is itself certified as endpoint-neutral on the same carrier;
\item UEAP status preservation: the report does not exceed zerohood, contribution, amplitude, and bridge support;
\item ATS trajectory preservation: the same anchor, tensor, and skin decomposition is preserved through the route.
\end{enumerate}
\end{definition}

\begin{audit}[Bridge object anti-circularity]
\(\BzetaEF(\rho)\) does not contain \(\Lcrit(\rho)\) by definition.  It is a route- and report-preservation bridge.  A proof may instantiate it in bridge-complete form by showing endpoint-neutral amplitude order for the same zerohood contribution.  Once endpoint-neutral order is obtained, the projection theorem yields line readout.  The line readout is downstream of bridge completion, not built into the name of the bridge.
\end{audit}

\subsection{Critical-normalized order invariant}

\begin{definition}[Critical-normalized amplitude order]
For \(\rho=\beta+i\gamma\), define the critical-normalized amplitude order of the associated explicit-formula contribution by
\[
\ordhalf(\rho):=\lim_{x\to\infty}\frac{\log |x^{\beta-1/2}|}{\log x},
\]
whenever the contribution class is fixed by \(\ModelEF(\zeta,\psi,R)\).  The corresponding order discriminator is
\[
\Dord(\rho):\Longleftrightarrow \ordhalf(\rho)\ne 0.
\]
\end{definition}

\begin{theorem}[Normalized amplitude order invariant]
Under \(\ModelEF(\zeta,\psi,R)\), if \(\rho=\beta+i\gamma\), then
\[
\ordhalf(\rho)=\beta-\frac12.
\]
Consequently,
\[
\Lcrit(\rho)\Longleftrightarrow \ordhalf(\rho)=0,
\qquad
\neg\Lcrit(\rho)\Longleftrightarrow \Dord(\rho).
\]
\end{theorem}

\begin{proof}
For \(x>1\), \(|x^{\beta-1/2}|=x^{\beta-1/2}\).  Therefore
\[
\frac{\log |x^{\beta-1/2}|}{\log x}=\frac{(\beta-1/2)\log x}{\log x}=\beta-\frac12.
\]
The limit is constant in \(x\), so it equals \(\beta-1/2\).  The two equivalences follow from \(\Lcrit(\rho)\) being \(\beta=1/2\).
\end{proof}

\begin{audit}[Order invariant is not a hidden estimate]
The theorem does not estimate \(\psi(x)-x\), does not assert dominance of a zero contribution, and does not deny cancellation in the full explicit formula.  It identifies the order class of the contribution associated with a fixed zero under the fixed route.  Pairing and cancellation remain bridge issues, not reasons to erase the order invariant at support level.
\end{audit}

\subsection{Zerohood-to-contribution equivalence}

\begin{definition}[Same-zero explicit-formula contribution]
\(\SameZeroEF(\rho)\) means that the explicit-formula endpoint model contains the contribution class associated with the same nontrivial zeta zero \(\rho\), with the same carrier, same route, and same critical normalization.
\end{definition}

\begin{theorem}[Zerohood-to-contribution equivalence]
Under \(\CRHEF(R,\zeta,\psi)\wedge \OfficialEFRoute(R,\zeta,\psi)\wedge \ModelEF(\zeta,\psi,R)\),
\[
\Az(\rho)\Longleftrightarrow \SameZeroEF(\rho)
\]
for the endpoint-role purposes of this proof.
\end{theorem}

\begin{proof}
The forward direction is the zero-coverage clause of \(\ModelEF\): every nontrivial zerohood instance in the official carrier has an associated explicit-formula contribution class.  The reverse direction is contribution soundness and same-zero identity: every contribution used by the endpoint model is attached to the same zeta zerohood carrier and the same zero instance.  The equivalence is role-equivalence for the endpoint route; it does not add an analytic zero-location theorem.
\end{proof}

\begin{audit}[No surrogate contribution]
A symmetry partner, random-matrix eigenvalue, spectral model point, finite verification entry, or smoothed auxiliary residue may support a bridge, but it is not \(\SameZeroEF(\rho)\) unless same-zero identity and carrier preservation are certified.  This is the ATS anchor condition applied to the EF route.
\end{audit}

\subsection{Line-readout-to-neutral-order bridge}

\begin{theorem}[Line readout equals neutral order on the EF route]
Under \(\ModelEF(\zeta,\psi,R)\), the structural critical-line readout role and the EF neutral-order role are equivalent:
\[
\Lcrit(\rho)\Longleftrightarrow \bigl(\ordhalf(\rho)=0\bigr)
\Longleftrightarrow \CritAmpOcc(\rho).
\]
The off-line branch is likewise equivalent to non-neutral order:
\[
\neg\Lcrit(\rho)\Longleftrightarrow \Dord(\rho)\Longleftrightarrow \OffAmpExcl(\rho).
\]
\end{theorem}

\begin{proof}
The first equivalence is the normalized order theorem.  The second equivalence is the definition of critical-neutral amplitude occupation.  The off-line equivalence is the exact complement: non-critical readout is nonzero normalized order, which is precisely off-line amplitude exclusion in the fixed EF model.
\end{proof}

\begin{audit}[Equivalence is projection, not decision]
The theorem fixes the role-content of the branches under \(\OfficialEFRoute\).  It does not decide whether \(\ordhalf(\rho)\) is zero for every zero.  That decision is obtained only after the local endpoint-countercase route is closed.
\end{audit}

\subsection{Bridge completeness and bridge failure}

\begin{definition}[EF bridge completion]
\(\BridgeCompleteEF(\rho)\) means that the EF route supplies a same-zero, carrier-preserving, UEAP-adequate, ATS-stable bridge that neutralizes the endpoint-facing order contribution for \(\rho\) in the critical-normalized amplitude image.  Formally, it entails
\[
\BridgeCompleteEF(\rho)\Rightarrow \ordhalf(\rho)=0.
\]
\(\BridgeFailEF(\rho)\) means that no such endpoint-neutral bridge is supplied for the endpoint-facing use under audit.
\end{definition}

\begin{theorem}[Bridge-complete EF route yields critical-line readout]
Under \(\ModelEF(\zeta,\psi,R)\),
\[
\BridgeCompleteEF(\rho)\Rightarrow \Lcrit(\rho).
\]
\end{theorem}

\begin{proof}
Bridge completion yields endpoint-neutral normalized order by definition of the EF bridge-completion role.  The line-readout-to-neutral-order theorem identifies neutral order with \(\Lcrit(\rho)\) on the fixed EF model.  Hence critical-line readout follows.  The proof does not define the bridge as line readout; it uses the EF order interface as the bridge output.
\end{proof}

\begin{corollary}[Off-line assumption excludes bridge-complete neutralization]
Under \(\ModelEF(\zeta,\psi,R)\),
\[
\neg\Lcrit(\rho)\Rightarrow \neg\BridgeCompleteEF(\rho).
\]
\end{corollary}

\begin{proof}
If \(\BridgeCompleteEF(\rho)\), then the previous theorem gives \(\Lcrit(\rho)\), contradicting \(\neg\Lcrit(\rho)\).  Therefore bridge-complete neutralization is unavailable inside the exact-complement off-line subproof.
\end{proof}

\begin{audit}[Why bridge failure is not assumed globally]
The corollary is local to the off-line reductio branch.  The manuscript does not globally assert that EF bridge completion fails.  The final proof instead shows that the off-line branch cannot survive: if a bridge completes, line readout follows; if it does not and endpoint counterforce remains on the fixed carrier, the non-neutral order becomes discriminator work.
\end{audit}

\subsection{Bridge-discriminator-domain fork}

\begin{theorem}[Structural-to-EF bridge fork]
Assume \(\Az(\rho)\), \(\OfficialEFRoute(R,\zeta,\psi)\), \(\ModelEF(\zeta,\psi,R)\), and \(\ProjRoleEF(\neg\Lcrit(\rho),\rho)\).  If the projected off-line branch is used with endpoint counterforce, then exactly one of the following endpoint-route dispositions applies:
\[
\BridgeCompleteEF(\rho)\quad\vee\quad \Damp(\rho)\quad\vee\quad \DomainShiftEF(\rho).
\]
\end{theorem}

\begin{proof}
Projection-role content binds the local negative branch to non-neutral order content on the EF amplitude carrier.  Endpoint counterforce says that this content is used to defeat critical-neutral endpoint closure.  There are three possible dispositions.  First, a bridge may neutralize the order content in the same endpoint image; this is \(\BridgeCompleteEF(\rho)\).  Second, the non-neutral order may remain endpoint-standing on the same fixed image.  Then it is not support, bookkeeping, or skin; it changes branch exclusion and endpoint standing.  By the endpoint-discriminator trichotomy, it is independent same-domain amplitude discriminator work, i.e. \(\Damp(\rho)\).  Third, the act may avoid same-domain discrimination only by changing the endpoint image, carrier, route, normalization, or same-zero identity; that is \(\DomainShiftEF(\rho)\).  UEAP excludes overreporting a support-level non-neutral order as endpoint result, and ATS excludes hiding a carrier or tensor change as notation skin.  These cases exhaust the route.
\end{proof}

\begin{audit}[The fork is the equivalence bridge]
This is the bridge missing from the compact version.  It does not say non-neutral order is impossible.  It says endpoint-facing non-neutral order has only three lawful dispositions: bridge-complete neutralization, domain shift, or independent discriminator work.  The local off-line reductio eliminates bridge completion; fixed-domain endpoint use eliminates domain shift; the remaining branch is the amplitude discriminator.
\end{audit}

\begin{corollary}[Fixed-domain off-line counterforce yields the amplitude discriminator]
Under \(\Az(\rho)\), \(\OfficialEFRoute(R,\zeta,\psi)\), \(\ModelEF(\zeta,\psi,R)\), \(\FixedDomain(R)\), \(\ProjRoleEF(\neg\Lcrit(\rho),\rho)\), \(\Counterforce_{\RH\text{-}\EF}(\rho)\), and the local off-line assumption \(\neg\Lcrit(\rho)\), one has
\[
\Damp(\rho).
\]
\end{corollary}

\begin{proof}
By the structural-to-EF bridge fork, the endpoint-counterforce use yields \(\BridgeCompleteEF(\rho)\), \(\Damp(\rho)\), or \(\DomainShiftEF(\rho)\).  The local off-line assumption excludes \(\BridgeCompleteEF(\rho)\), since bridge completion yields \(\Lcrit(\rho)\).  Fixed-domain preservation excludes \(\DomainShiftEF(\rho)\).  Therefore \(\Damp(\rho)\) remains.
\end{proof}

\begin{audit}[Hostile-referee target after the bridge]
A successful objection must now attack a specific branch of the fork: show a bridge-complete neutralization compatible with \(\neg\Lcrit(\rho)\), show a same-domain endpoint use that is a domain shift and yet not a domain shift, or exhibit a fourth endpoint-standing role for non-neutral order that is not support, bookkeeping, bridge completion, domain shift, or independent discriminator work.
\end{audit}


\section{Native Counterforce Collapse: No Fifth Endpoint Role}\label{sec:native-counterforce-collapse}

This section closes the remaining hostile-referee target left after the structural-to-explicit-formula bridge.  A critic may grant that the off-line branch projects to non-neutral explicit-formula order, but still insist that this is simply native same-carrier falsification rather than illicit endpoint discrimination.  The present section proves that ordinary native falsification is not a fifth endpoint role.  Native analytic origin is preserved; endpoint-standing authority is classified.

\subsection{Native analytic content versus native endpoint authority}

\begin{definition}[Native analytic content]
A datum is native analytic content for the RH-EF route when it is produced inside the fixed zeta/prime-trace carrier by the official explicit-formula model.  In particular, if \(\rho=\beta+i\gamma\) is a nontrivial zero, then the order class
\[
\ordhalf(\rho)=\beta-\frac12
\]
is native analytic content of the associated contribution class.
\end{definition}

\begin{definition}[Native endpoint counterforce]
\(\NativeCounterforce_{\RH\text{-}\EF}(\rho)\) means that native explicit-formula non-neutral order associated with the same zerohood instance \(\rho\) is used with endpoint-standing force to defeat the critical-neutral RH-EF endpoint while preserving the same official route, same zero, same prime-trace carrier, and same act-time endpoint evaluation.  In this manuscript, that endpoint-standing force is supplied only by \(\OfficialCxForce_R(\RH^-_{\mathrm{bare}})\), \(\RHOffCxForce(\rho)\), or the local exact-complement force \(\LocalCxForce_{\RH\text{-}\EF}(\neg\Lcrit(\rho);\Lcrit(\rho))\) together with \(\Counterforce_{\RH\text{-}\EF}(\rho)\).  Native counterforce is therefore not a raw property of \(\ordhalf(\rho)\); it is a use role.
\end{definition}

\begin{audit}[The point of the definition]
The definition does not deny that \(\ordhalf(\rho)\ne0\) is native analytic content.  It isolates the stronger use: the native datum is no longer only descriptive or support-level; it is being used as endpoint-standing counterforce against the critical-neutral endpoint image.  The proof below classifies that use.  It does not outlaw the datum.
\end{audit}

\begin{lemma}[Native content is not native authority]
Native analytic origin does not by itself supply endpoint-standing authority.  In symbols,
\[
\Dord(\rho)\not\Rightarrow \NativeCounterforce_{\RH\text{-}\EF}(\rho).
\]
\end{lemma}

\begin{proof}
\(\Dord(\rho)\) states that the associated contribution has non-neutral critical-normalized order.  It does not state that the datum resolves the endpoint, defeats the endpoint, excludes the positive branch, licenses downstream reuse, or stands as official endpoint result.  UEAP separates support-level report from endpoint-standing report, and ATS separates tensor content from the act that uses that tensor content as endpoint authority.  Therefore native analytic content alone is not native endpoint authority.
\end{proof}

\subsection{The native-counterforce fork}

\begin{definition}[Bridge-neutralized native counterforce]
\(\BridgeNeutralized(\rho)\) means that the apparent non-neutral endpoint-facing contribution has a same-zero, same-route, carrier-preserving bridge that neutralizes its endpoint order in the critical-normalized image.  Bridge-neutralization is stronger than cancellation as a numerical or formal possibility: it is endpoint-neutralization certified at the same standing-bearing role level.
\end{definition}

\begin{theorem}[Native-counterforce role exhaustion]
For the fixed RH-EF endpoint carrier, native endpoint counterforce by non-neutral order has no fifth role.  If
\[
\Az(\rho)\wedge \OfficialEFRoute(R,\zeta,\psi)\wedge \ModelEF(\zeta,\psi,R)
\wedge \NativeCounterforce_{\RH\text{-}\EF}(\rho),
\]
then exactly one of the following endpoint-role dispositions applies:
\[
\BridgeNeutralized(\rho)\vee \Support(\rho)\vee \Book(\rho)\vee \DomainShiftEF(\rho)\vee \Damp(\rho).
\]
\end{theorem}

\begin{proof}
The datum under audit is native analytic content because it is produced by the fixed explicit-formula route.  The question is what role it plays when used with endpoint-standing counterforce.

If a same-zero, same-route bridge proves endpoint-neutralization of the non-neutral contribution at the critical-normalized endpoint image, then \(\BridgeNeutralized(\rho)\) holds.  If the datum is used only as part of an argument, obstruction, heuristic, or local analytic description, then it is \(\Support(\rho)\), and it lacks endpoint-standing counterforce.  If the datum changes no endpoint standing, reportability, downstream reuse, branch exclusion, or admissible continuation, then it is bookkeeping.  If the datum obtains endpoint force only by changing the zeta carrier, the prime-trace route, the normalization, same-zero identity, the smoothing/truncation regime as tensor authority, or the endpoint image, then \(\DomainShiftEF(\rho)\) holds.

After these cases are removed, the remaining use is endpoint-standing, same-domain, non-neutral, non-support, non-bookkeeping, non-bridge-neutralized, and non-domain-shifting.  It changes whether the live off-line branch can defeat critical-neutral endpoint closure.  By the endpoint-discriminator trichotomy and the ATS no-second-classifier lock, that remaining use is independent same-domain amplitude endpoint-status discrimination.  This is exactly \(\Damp(\rho)\).  The cases are exhaustive because every endpoint-facing use either changes endpoint governance or it does not; if it changes governance, it either preserves the invariant bundle or changes it.
\end{proof}

\begin{audit}[Ordinary native falsification is not a fifth role]
The theorem answers the strongest hostile objection.  An off-line zero's non-neutral order is native analytic content.  But when that content is used as endpoint-standing counterforce, ordinary native falsification is not a role outside the endpoint taxonomy.  It is bridge-neutralized, support-only, bookkeeping, domain/interface shift, or independent endpoint-status discrimination.  The proof does not deny counterexample syntax; it classifies endpoint-standing counterexample force.
\end{audit}

\subsection{Collapse under the live off-line countercase}

\begin{lemma}[Support-only native content has no endpoint counterforce]
\[
\Support(\rho)\Rightarrow \neg\Counterforce_{\RH\text{-}\EF}(\rho).
\]
\end{lemma}

\begin{proof}
Support-only material may contribute to a proof, bridge program, heuristic, obstruction, or local analytic diagnosis, but it does not itself defeat the official endpoint.  Endpoint counterforce is precisely the role in which a branch blocks endpoint closure.  Therefore support-only use and endpoint counterforce are incompatible.
\end{proof}

\begin{lemma}[Bookkeeping native content has no endpoint counterforce]
\[
\Book(\rho)\Rightarrow \neg\Counterforce_{\RH\text{-}\EF}(\rho).
\]
\end{lemma}

\begin{proof}
Bookkeeping changes no endpoint standing, reportability, downstream reuse, branch exclusion, or admissible continuation.  Endpoint counterforce changes exactly the endpoint-facing status of the live branch.  Hence bookkeeping cannot carry endpoint counterforce.
\end{proof}

\begin{lemma}[Bridge-neutralization is incompatible with the live off-line endpoint countercase]
Under \(\ModelEF(\zeta,\psi,R)\), the exact-complement off-line subproof satisfies
\[
\neg\Lcrit(\rho)\Rightarrow \neg\BridgeNeutralized(\rho).
\]
\end{lemma}

\begin{proof}
Bridge-neutralization is endpoint-level neutralization of the same-zero non-neutral order in the critical-normalized image.  In the fixed EF model, neutral endpoint order is equivalent to \(\Lcrit(\rho)\).  Thus \(\BridgeNeutralized(\rho)\) gives the same endpoint effect as bridge-complete neutralization: \(\Lcrit(\rho)\).  This contradicts the live off-line assumption \(\neg\Lcrit(\rho)\).  The argument is local to the reductio subproof; it does not globally deny that bridge-neutralization may be proved for a zero.
\end{proof}

\begin{theorem}[Native same-carrier counterforce collapse]
Under the fixed RH-EF endpoint context,
\[
\begin{aligned}
&\Az(\rho)\wedge \OfficialEFRoute(R,\zeta,\psi)\wedge \ModelEF(\zeta,\psi,R)\wedge \FixedDomain(R)\\
&\qquad\wedge\ProjRoleEF(\neg\Lcrit(\rho),\rho)\wedge \Counterforce_{\RH\text{-}\EF}(\rho)
\wedge \neg\Lcrit(\rho)
\Rightarrow \Damp(\rho).
\end{aligned}
\]
\end{theorem}

\begin{proof}
Projection-role content and endpoint counterforce make the non-neutral EF order into native endpoint counterforce, so \(\NativeCounterforce_{\RH\text{-}\EF}(\rho)\) holds.  By native-counterforce role exhaustion, the role is bridge-neutralized, support-only, bookkeeping, domain shift, or \(\Damp(\rho)\).  Bridge-neutralization is excluded by the live off-line assumption.  Support-only and bookkeeping are incompatible with endpoint counterforce.  Domain shift is excluded by \(\FixedDomain(R)\).  The remaining exhaustive branch is \(\Damp(\rho)\).
\end{proof}

\begin{audit}[Final hardening of the hinge]
The theorem is the sharpened hinge.  It does not say \(\neg\Lcrit(\rho)\Rightarrow\Damp(\rho)\).  It says that once the off-line branch is projected through the official EF route and used with endpoint counterforce on the fixed carrier, the ordinary-native-falsification reading collapses into the same endpoint taxonomy: bridge-neutralized, support-only, bookkeeping, domain shift, or independent discriminator.  The first four branches are incompatible with the live local countercase and fixed-domain endpoint use, so the discriminator remains.
\end{audit}

\begin{corollary}[No native-falsification escape from A+ closure]
There is no same-carrier endpoint-standing role in which non-neutral EF order defeats RH while remaining neither bridge-neutralized, support-only, bookkeeping, domain-shifting, nor an independent amplitude discriminator.
\end{corollary}

\begin{proof}
This is the contrapositive of native-counterforce role exhaustion together with the collapse theorem.  A purported native-falsification escape that changes endpoint standing must occupy one of the enumerated roles.  If it avoids all enumerated roles, it changes no endpoint standing and is endpoint-inert.
\end{proof}

\section{Prime-Trace Neutrality and Non-Neutral Scale Standing}

\subsection{Descriptive amplitude versus endpoint-standing amplitude}

\begin{definition}[Amplitude status levels]
For a zerohood instance \(\rho\), distinguish:
\begin{enumerate}
\item \emph{descriptive amplitude}: the algebraic scale class \(x^{\beta-1/2}\);
\item \emph{support-level amplitude}: use of the scale class in an analytic argument or bridge program;
\item \emph{endpoint-standing amplitude}: use of the scale class to defeat or settle the official RH endpoint on the fixed explicit-formula carrier.
\end{enumerate}
\end{definition}

\begin{definition}[Endpoint-standing non-neutral amplitude]
\(\AmpStand(\rho)\) means that the non-neutral scale \(x^{\beta-1/2}\ne x^0\) is used as endpoint-standing content against the critical-neutral amplitude image while preserving the same zeta/prime-trace carrier.
\end{definition}

\begin{theorem}[Descriptive non-neutral amplitude is lawful]
\[
\OffAmpExcl(\rho)\not\Rightarrow \AmpStand(\rho).
\]
Descriptive or support-level non-neutral amplitude remains lawful analytic content.
\end{theorem}

\begin{proof}
The existence of a scale class is not the same as endpoint-standing use.  UEAP separates the support report from the endpoint report.  ATS separates tensor content from the act that uses that tensor content as endpoint authority.  Therefore non-neutral amplitude by itself does not become endpoint-standing amplitude.
\end{proof}

\subsection{Prime-trace endpoint adequacy}

\begin{definition}[Prime-trace endpoint adequacy]
\(\PrimeTraceAdeq(R,\zeta,\psi)\) means that the explicit-formula route and critical-normalized prime-trace image are adequate for the official RH endpoint: critical-line readout projects to neutral amplitude occupation, off-line readout projects to non-neutral amplitude non-occupation, and the projection preserves the same zeta zerohood carrier.  It does not assert universal neutral occupation.
\end{definition}

\begin{theorem}[Neutral amplitude is endpoint image, not assumption]
\(\PrimeTraceAdeq(R,\zeta,\psi)\) fixes neutral amplitude as the positive endpoint image, but does not assert that every zero contribution occupies that image.  In symbols,
\[
\PrimeTraceAdeq(R,\zeta,\psi)\not\Rightarrow \forall \rho(\Az(\rho)\Rightarrow \CritAmpOcc(\rho)).
\]
\end{theorem}

\begin{proof}
Endpoint adequacy specifies what occupation would mean if the positive branch is occupied.  It is a typing condition, not an existence or universal occupation condition.  The distinction is the same as fixing a candidate endpoint image without asserting a non-failing candidate in the SAT setting.
\end{proof}

\subsection{Non-neutral scale standing}

\begin{theorem}[Non-neutral scale standing theorem]
If non-neutral amplitude is used with endpoint-standing force on the fixed RH-EF carrier, then it performs same-domain noncritical amplitude status work:
\[
\begin{aligned}
&\AmpStand(\rho)\wedge \FixedDomain(R)\\
&\qquad\wedge \OfficialEFRoute(R,\zeta,\psi)
\Rightarrow \NonCritAmpWork_R(\rho).
\end{aligned}
\]
\end{theorem}

\begin{proof}
Endpoint-standing non-neutral amplitude is not merely descriptive because it is used to decide endpoint standing.  It is not support-only because it defeats or displaces the endpoint.  It is not bookkeeping because removing it changes endpoint counterforce.  It is not carrier shift because the fixed-domain premise preserves the zeta/prime-trace carrier.  It is not repair because the standing is asserted at the local act time.  It is not coequal occupation because the single amplitude interface says non-neutral scale is non-occupation of the critical-neutral image.  Therefore it is same-domain noncritical amplitude status work.
\end{proof}

\begin{audit}[No RH-equivalent bound smuggled]
The theorem does not assert a bound such as \(\psi(x)-x=O(x^{1/2+\epsilon})\), nor does it deny that a hypothetical off-line zero would affect the prime trace.  It classifies only the endpoint-standing use of the non-neutral scale.  Descriptive and support uses remain lawful.
\end{audit}

\begin{theorem}[Non-native amplitude theorem]
On the fixed critical-neutral amplitude image, same-domain noncritical amplitude status work is non-native as endpoint authority unless it is support, bookkeeping, carrier shift, repair, surrogate substitution, or lawful coequal occupation.  Consequently,
\[
\NonCritAmpWork_R(\rho)
\Rightarrow \NonNativeDisc(\rho).
\]
\end{theorem}

\begin{proof}
The adjective ``native'' is role-relative.  The scale exponent is native analytic data in the explicit formula.  It is not native as endpoint authority over the critical-neutral image when it claims to defeat that image while preserving the same carrier.  If the non-neutral scale is only descriptive, it is not endpoint-status work.  If it is support, it does not resolve the endpoint.  If it changes the carrier or route, it is not same-domain.  If it is repair, it violates act-time finality.  If it is coequal occupation, it has changed to the broad outcome carrier.  After these lawful roles are removed, the remaining role is an independent amplitude endpoint-status discriminator.
\end{proof}

\begin{audit}[Native-data objection]
A classical analyst may correctly say that \(x^{\beta-1/2}\) is native analytic data.  The manuscript agrees.  The excluded object is not the data but its endpoint-standing use as same-carrier authority over the critical-neutral image.  This distinction is what prevents the proof from banning analytic content.
\end{audit}

\section{Amplitude Separator Occupation and Discriminator Closure}

\subsection{Separator occupation}

\begin{definition}[Local amplitude separator occupation]
\(\AmpSepLoc(\rho)\) means that off-line amplitude non-occupation is used inside the live exact-complement reductio subproof as the endpoint-defeating negative branch on the fixed RH-EF carrier.  It is local separator occupation, not global endpoint outcome.
\end{definition}

\begin{definition}[Global amplitude separator occupation]
\(\AmpSepOcc(\rho)\) means that off-line amplitude non-occupation is used as the official endpoint-resolving negative branch for RH on the fixed RH-EF carrier.
\end{definition}

\begin{theorem}[Separator occupation is endpoint-standing non-neutral amplitude]
\[
\begin{aligned}
\AmpSepLoc(\rho)&\Rightarrow \AmpStand(\rho),\\
\AmpSepOcc(\rho)&\Rightarrow \AmpStand(\rho).
\end{aligned}
\]
\end{theorem}

\begin{proof}
By definition, separator occupation is the endpoint-facing use of non-neutral amplitude non-occupation.  Local separator occupation has this role inside the reductio subproof; global separator occupation has it as official negative endpoint resolution.  Both are endpoint-standing uses rather than descriptive or support reports.
\end{proof}

\subsection{Endpoint-discriminator trichotomy}

\begin{definition}[Amplitude endpoint-status discriminator]
\(C\) is an amplitude endpoint-status discriminator for \(\rho\) when it changes endpoint standing, branch exclusion, reportability, downstream reuse, or admissible continuation for the fixed RH-EF endpoint act involving \(\rho\).
\end{definition}

\begin{theorem}[Endpoint-discriminator trichotomy]
Every amplitude endpoint-status classifier attached to a fixed RH-EF endpoint act is exactly one of: bookkeeping, governance-equivalent to the admitted interface, independent same-domain endpoint-status work, or domain shift.
\end{theorem}

\begin{proof}
If the classifier changes no endpoint verdict, it is bookkeeping.  If it reproduces the existing admitted interface, it is governance-equivalent.  If it changes endpoint-status verdicts while preserving the same invariant bundle, it is independent same-domain work.  If it changes the invariant bundle, it is domain shift.  The cases are exhaustive by whether governance work changes and whether the change remains on the same fixed domain.
\end{proof}

\begin{theorem}[Hinge theorem: noncritical amplitude work induces the amplitude discriminator]
\[
\NonCritAmpWork_R(\rho)\Rightarrow \Damp(\rho).
\]
\end{theorem}

\begin{proof}
Noncritical amplitude work is endpoint-facing, same-carrier, non-neutral amplitude status determination.  It is not positive critical-neutral occupation.  It is not support-only, bookkeeping, repair, carrier shift, surrogate substitution, or coequal occupation by the hypotheses built into \(\NonCritAmpWork\).  By the endpoint-discriminator trichotomy, the only remaining role is independent same-domain endpoint-status work.  In the explicit-formula amplitude carrier, that role is exactly \(\Damp(\rho)\).
\end{proof}

\begin{audit}[This is the central hostile-referee target]
The theorem does not say that non-neutral amplitude is impossible.  It says that noncritical amplitude \emph{status work}, after support/bookkeeping/carrier-shift/repair/coequal roles are excluded, is independent endpoint-status discrimination.  A successful objection must produce a remaining lawful role for endpoint-standing non-neutral amplitude that is not independent discrimination.
\end{audit}

\begin{theorem}[Endpoint use triggers no-independent amplitude closure]
\[
\EndpointUse_R(\AmpSepLoc(\rho))\wedge \FixedDomain(R)
\Rightarrow \Namp(\rho).
\]
\end{theorem}

\begin{proof}
Local separator occupation is endpoint use inside a governed reductio subproof.  Endpoint use on the fixed domain instantiates the kernel and the A+ consequence layer.  The local A+ consequence includes no independent same-domain endpoint-status discriminator on the same RH-EF carrier.  This is the predicate \(\Namp(\rho)\).
\end{proof}

\begin{theorem}[No-independent amplitude closure excludes the discriminator]
\[
\Namp(\rho)\Rightarrow \neg\Damp(\rho).
\]
\end{theorem}

\begin{proof}
\(\Namp\) is the no-independent-discriminator closure specialized to the amplitude endpoint-status discriminator.  Therefore it denies \(\Damp\).
\end{proof}

\begin{corollary}[Local amplitude separator occupation is impossible]
\[
\AmpSepLoc(\rho)\Rightarrow \bot.
\]
\end{corollary}

\begin{proof}
Assume \(\AmpSepLoc(\rho)\).  Separator occupation gives endpoint-standing non-neutral amplitude.  By the non-neutral scale standing theorem and the hinge theorem, \(\Damp(\rho)\).  Endpoint use of the local separator gives \(\Namp(\rho)\), hence \(\neg\Damp(\rho)\).  Contradiction.
\end{proof}

\section{Anti-Trivialization Discipline}

\begin{theorem}[Exact-complement assumptions are not automatically separators]
For an arbitrary proposition \(P\), a local reductio assumption \([\neg P]_i\) does not become separator occupation merely because it is the exact complement of \(P\).
\end{theorem}

\begin{proof}
Separator occupation requires a fixed endpoint-image carrier, an official endpoint route, projection-role content on that carrier, endpoint-counterforce use, and exclusion of support, bookkeeping, carrier shift, repair, surrogate substitution, and coequal occupation.  A bare local assumption has none of those features by itself.  Therefore arbitrary negation is not separator occupation.
\end{proof}

\begin{corollary}[RH-specific separator conditions]
For RH-EF, \([\neg\Lcrit(\rho)]_i\) becomes local amplitude separator occupation only after the official explicit-formula route binds it to non-neutral amplitude projection and endpoint counterforce uses that projection as endpoint-defeating content.
\end{corollary}

\begin{audit}[No \(\neg P\Rightarrow P\) schema]
The proof cannot be instantiated for arbitrary \(P\).  It uses the explicit-formula amplitude projection of zeta zerohood.  Without such a projection and a proof that endpoint-standing non-neutral amplitude is independent discriminator work, the A+ contradiction cannot be triggered.
\end{audit}

\section{Local Reductio Engine and Universal Endpoint}

\subsection{Exact-complement annotation}

\begin{definition}[Exact-complement reductio annotation]
\[
\mathrm{ReductioCountercase}_{\RH\text{-}\EF}(\neg\Lcrit(\rho);\Lcrit(\rho))
\]
is the proof-act annotation generated when the sole local assumption \([\neg\Lcrit(\rho)]_i\) is opened as the live exact complement of \(\Lcrit(\rho)\) under \(\Az(\rho)\) in the official RH-EF endpoint-closure proof.  It is not an additional object-level premise.
\end{definition}

\begin{theorem}[Annotation is not a premise]
The annotation does not transform the discharged assumption into
\[
\neg\Lcrit(\rho)\wedge \mathrm{ReductioCountercase}_{\RH\text{-}\EF}(\neg\Lcrit(\rho);\Lcrit(\rho)).
\]
It records the local proof role of the sole assumption \([\neg\Lcrit(\rho)]_i\).
\end{theorem}

\begin{proof}
A reductio subproof has an object-level assumption and a proof role for that assumption.  When the assumption is the exact complement of the target in a fixed endpoint proof, the proof role is live countercase use.  The role annotation is metaproof bookkeeping for the subproof; it is not a formula added to the object language.
\end{proof}

\begin{theorem}[Local countercase to endpoint counterforce]
Under \(\Az(\rho)\wedge\OSEvalEF(R,\zeta,\psi)\), the exact-complement local annotation yields endpoint counterforce:
\[
\mathrm{ReductioCountercase}_{\RH\text{-}\EF}(\neg\Lcrit(\rho);\Lcrit(\rho))
\Rightarrow \Counterforce_{\RH\text{-}\EF}(\rho).
\]
\end{theorem}

\begin{proof}
The live exact-complement assumption is introduced precisely to block the target \(\Lcrit(\rho)\) under the official endpoint split.  That is local endpoint-facing force, not global theoremhood and not official negative endpoint resolution.  The endpoint-closure evaluation supplies the fixed target and branch slots required for this local role.
\end{proof}

\begin{theorem}[Local exact-complement force is local endpoint use]
Under \(\Az(\rho)\wedge\OSEvalEF(R,\zeta,\psi)\), the exact-complement reductio role and endpoint counterforce yield local endpoint use:
\[
\mathrm{ReductioCountercase}_{\RH\text{-}\EF}(\neg\Lcrit(\rho);\Lcrit(\rho))
\wedge \Counterforce_{\RH\text{-}\EF}(\rho)
\Rightarrow \LocalEndpointUse_R(\neg\Lcrit(\rho)).
\]
\end{theorem}

\begin{proof}
The reductio annotation supplies \(\LocalCxForce_{\RH\text{-}\EF}(\neg\Lcrit(\rho);\Lcrit(\rho))\): the off-line branch is the live exact-complement countercase to the local target.  Endpoint counterforce says that this branch is being used to block endpoint closure.  The general local counterexample-force theorem, Theorem~\ref{thm:local-counterexample-force-endpoint-use}, therefore gives local endpoint use.  This is still local; it does not assert global official negative endpoint resolution.
\end{proof}

\begin{corollary}[Local A+ closure for the exact-complement countercase]
Under fixed-domain RH-EF evaluation,
\[
\LocalEndpointUse_R(\neg\Lcrit(\rho))\wedge\FixedDomain(R)
\Rightarrow \Aplus_{R,\mathrm{loc}}(\neg\Lcrit(\rho))\Rightarrow \Namp(\rho).
\]
\end{corollary}

\begin{proof}
Apply Theorem~\ref{thm:local-endpoint-use-aplus}.  In the RH-EF route, the local A+ consequence relevant to amplitude separator use is the no-independent-amplitude-discriminator closure \(\Namp(\rho)\).
\end{proof}


\begin{theorem}[Endpoint counterforce induces local amplitude separator occupation]
Under the fixed RH-EF context,
\[
\Az(\rho)\wedge \ProjRoleEF(\neg\Lcrit(\rho),\rho)
\wedge \Counterforce_{\RH\text{-}\EF}(\rho)
\Rightarrow \AmpSepLoc(\rho).
\]
\end{theorem}

\begin{proof}
Projection-role content says that the local negative branch has non-neutral amplitude non-occupation as its endpoint-role content.  Endpoint counterforce says that this projected non-occupation is being used to defeat the critical-neutral endpoint image inside the local subproof.  It is therefore local amplitude separator occupation.  This is not a promotion of raw negation; it is the endpoint-use classification of the projected amplitude role under counterforce.
\end{proof}

\begin{theorem}[Local countercase contradiction]
Assume \(\Az(\rho)\), \(\CRHEF(R,\zeta,\psi)\), \(\OfficialEFRoute(R,\zeta,\psi)\), \(\ModelEF(\zeta,\psi,R)\), \(\FixedDomain(R)\), and \(\OSEvalEF(R,\zeta,\psi)\).  Then the local exact-complement assumption \([\neg\Lcrit(\rho)]_i\) leads to contradiction.
\end{theorem}

\begin{proof}
Open \([\neg\Lcrit(\rho)]_i\) as the live exact complement of \(\Lcrit(\rho)\).  The annotation theorem supplies the reductio-countercase role but no additional object-level premise.  The local countercase has endpoint counterforce.  By the local off-line projection theorem, it has projection-role content \(\ProjRoleEF(\neg\Lcrit(\rho),\rho)\).  Endpoint counterforce plus projection-role content gives \(\AmpSepLoc(\rho)\).  The same local exact-complement use is local endpoint use, so local A+ closure supplies \(\Namp(\rho)\).  Local amplitude separator occupation is impossible by the discriminator contradiction.  Hence the subproof derives \(\bot\).
\end{proof}

\begin{theorem}[Annotation discharge]
If the annotated exact-complement subproof derives contradiction from the sole object-level assumption \([\neg\Lcrit(\rho)]_i\), then \(\Lcrit(\rho)\) follows.  The discharged formula is \(\neg\Lcrit(\rho)\), not a strengthened conjunction.
\end{theorem}

\begin{proof}
The reductio annotation is a proof-act role, not an object-level assumption.  Therefore ordinary reductio discharges the sole opened formula.  Since the target is \(\Lcrit(\rho)\), contradiction under \(\neg\Lcrit(\rho)\) gives \(\Lcrit(\rho)\).
\end{proof}

\begin{theorem}[Arbitrary zerohood gives critical-line readout]
Under the fixed official RH-EF endpoint context,
\[
\Az(\rho)\Rightarrow \Lcrit(\rho).
\]
\end{theorem}

\begin{proof}
Assume \(\Az(\rho)\).  If \(\Lcrit(\rho)\), done.  If \(\neg\Lcrit(\rho)\), open the exact-complement local subproof.  The local countercase contradiction discharges the assumption and yields \(\Lcrit(\rho)\).  Thus \(\Az(\rho)\Rightarrow\Lcrit(\rho)\).
\end{proof}

\begin{corollary}[Riemann Hypothesis]
\[
\forall \rho\,\bigl(\Az(\rho)\Rightarrow \Lcrit(\rho)\bigr).
\]
\end{corollary}

\begin{proof}
The preceding theorem holds for arbitrary \(\rho\).  Universal generalization yields the displayed formula, which is the Riemann Hypothesis in the standard critical-strip formulation.
\end{proof}

\section{Ordinary Negative Theoremhood, Incompleteness, and Open Status}

\begin{theorem}[Ordinary negative theoremhood is not banned]
A hypothetical ordinary theorem of \(\RH^-_{\mathrm{bare}}\) is not rejected merely because it is negative.  Ordinary theoremhood is proof legitimacy.  It becomes endpoint-role disciplined only when used as official endpoint resolution or endpoint counterforce on the fixed carrier.
\end{theorem}

\begin{proof}
The kernel applies to genuine theoremhood regardless of polarity.  The excluded object is not negative theoremhood as such, but endpoint-standing non-neutral amplitude separator work.  If a negative theorem is not used as official endpoint resolution, it is not the endpoint act under audit.  If it is so used, ordinary official resolution supplies endpoint use and triggers the endpoint-role discipline.
\end{proof}

\begin{theorem}[Official negative endpoint resolution becomes endpoint use]
\[
\OrdOfficialRes_R(\RH^-_{\mathrm{bare}})\wedge \OfficialEFRoute(R,\zeta,\psi)
\Rightarrow \EndpointUse_R(\RH^-_{\mathrm{bare}}).
\]
\end{theorem}

\begin{proof}
Ordinary official resolution fixes target, carrier, branch, theorem-status, downstream reuse, act-time finality, and redescription-stable endpoint role.  Those are exactly endpoint-use conditions.  The official explicit-formula route fixes how the negative branch is projected into amplitude content.
\end{proof}

\begin{definition}[Same-domain incompleteness or independence objection]
\(\SameDomIncomp_{\RH}(C)\) means that a Godel-style, independence, true-but-unprovable, model-theoretic, stronger-theory, or proof-system distinction is asserted to re-enter the fixed RH-EF endpoint carrier as a standing-bearing distinction.
\end{definition}

\begin{theorem}[Incompleteness firewall]
If a same-domain incompleteness or independence distinction does not affect endpoint standing, reportability, downstream reuse, branch exclusion, or admissible continuation, it is endpoint-inert.  If it does affect one of those features, it is endpoint-admissibility-relevant gate work and is subject to the same A+ no-independent-discriminator closure.
\end{theorem}

\begin{proof}
Endpoint force means altering the standing-bearing endpoint act.  A distinction that does not alter the act is inert with respect to the endpoint.  A distinction that does alter the act is a classifier of endpoint standing, reportability, reuse, or continuation.  It therefore belongs to the endpoint-admissibility-relevant invariant bundle controlled by the kernel and A+ consequence layer.
\end{proof}

\begin{corollary}[Open status is not a third endpoint branch]
Open, unresolved, independent, unproved, or review status is epistemic status.  It is not a third endpoint truth branch of the RH-EF endpoint.
\end{corollary}

\section{Route-Locus Exhaustion and Weakening Resistance}

\subsection{Route-locus exhaustion}

\begin{longtable}{L{0.21\textwidth}L{0.43\textwidth}L{0.27\textwidth}}
\toprule
\textbf{Locus} & \textbf{RH-EF form} & \textbf{Disposition}\\
\midrule
Coordinate-only report & \(\beta\ne 1/2\) as coordinate data & Lawful descriptive content; not endpoint result.\\
Explicit-formula support & \(x^{\beta-1/2}\ne x^0\) as contribution data & Lawful support; not contradiction.\\
Endpoint amplitude separator & Non-neutral amplitude used as endpoint counterforce & Routes to \(\Damp\) and closure.\\
Official counterexample force & Off-line branch used to defeat the official RH endpoint & Endpoint use by the counterexample-force theorem; enters native-counterforce exhaustion.\\
Symmetry orbit & \(\rho,1-\rho,\bar\rho,1-\bar\rho\) & Carrier structure; not same-zero neutral occupation without bridge.\\
Pairing/cancellation & Combined terms in explicit formula & Bridge material; not endpoint-neutral unless proved.\\
Spectral surrogate & Hilbert--Polya-style carrier & Support or carrier shift unless bridged.\\
Positivity surrogate & Positivity criterion or equivalent form & Bridge if proved; otherwise support.\\
Numerical finite range & Verified zeros to height \(T\) & Finite support; not universal endpoint.\\
Density estimate & Scarcity or proportion results & Support; not absence of all off-line zeros.\\
Zero-free region & Region exclusion & Support unless it covers all possible off-line zeros.\\
Repair & Later proof validates earlier unbridged report & New act; no same-act repair.\\
Hidden fifth role & Endpoint-resolving, non-governing, same-carrier non-neutral amplitude & Impossible by trichotomy.\\
\bottomrule
\end{longtable}

\begin{theorem}[Route-locus exhaustion]
Every same-carrier endpoint-resolving off-line RH-EF route enters one of the route loci above.  If it enters none, it does no endpoint work and cannot resolve or defeat the endpoint.
\end{theorem}

\begin{proof}
A negative route must use coordinate displacement, amplitude projection, endpoint counterforce, symmetry, pairing, a surrogate, finite or asymptotic support, repair, or a hidden endpoint status.  The table covers whether the route changes coordinate, amplitude, carrier, time, locality, selection, or endpoint role.  A route changing none of these has no endpoint-relevant content.
\end{proof}

\subsection{Weakening resistance}

\begin{theorem}[No strict same-carrier weakening admits amplitude separator governance]
Any weakening of the RH-EF endpoint packet that preserves the same carrier, same zerohood predicate, same explicit-formula route, same critical-normalized endpoint image, same target adequacy, and same act-time finality either preserves the endpoint behavior extensionally, lowers non-neutral amplitude to support/bookkeeping, exits the carrier, or reintroduces independent amplitude endpoint governance.
\end{theorem}

\begin{proof}
If the weakening changes no endpoint-bearing verdict, it is extensionally inert.  If it lets non-neutral amplitude resolve the endpoint without critical-neutral occupation, it performs endpoint-status work.  If that work stays on the same carrier and is not support or bookkeeping, the trichotomy classifies it as independent.  If the weakening avoids independence by changing carrier, route, or endpoint image, it exits the same endpoint.
\end{proof}

\section{Denial Burden and Faithful Counterexample Worksheet}

The structural-to-EF bridge sharpens the central denial burden.  A critic may no longer attack a bare step from off-line coordinate to separator occupation; the manuscript now passes through order-class exactness, \(\BzetaEF\), and the bridge-discriminator-domain fork.  A successful objection must break one of those links or show that bridge-failing same-carrier non-neutral order can retain endpoint counterforce without either domain shift or independent discriminator work.


\subsection{Right-mode denial targets}

A standing objection must identify at least one of the following failures:
\begin{enumerate}
\item no-autonomous-counterexample-role failure: an official counterexample act defeats the endpoint while retaining official endpoint force and nevertheless not being endpoint use;
\item counterexample-force bridge failure: an official counterexample act defeats the endpoint while not fixing target, carrier, branch, reportability, downstream reuse, branch exclusion, act-time finality, and redescription-stable role;
\item official RH endpoint misidentified;
\item zeta/prime-trace carrier inadequacy;
\item context smuggling of RH or line support;
\item explicit-formula route does not project zerohood to amplitude content;
\item amplitude exactness failure for \(x^{\beta-1/2}\);
\item term-isolation or pairing/cancellation failure;
\item projection-role content failure;
\item single amplitude interface failure;
\item descriptive/support versus endpoint-standing distinction failure;
\item non-neutral scale standing theorem failure;
\item non-native amplitude theorem failure;
\item endpoint-discriminator trichotomy failure;
\item A+ no-independent closure failure;
\item local reductio annotation or discharge failure;
\item universalization failure.
\end{enumerate}

\subsection{Faithful counterexample worksheet}

\begin{longtable}{L{0.29\textwidth}L{0.64\textwidth}}
\toprule
\textbf{Target} & \textbf{Successful objection must produce}\\
\midrule
No autonomous counterexample role & A same-carrier official counterexample act that defeats the endpoint, retains reportable negative endpoint force, licenses downstream reuse, excludes the positive endpoint branch, and nevertheless is not endpoint use.\\
Counterexample-force bridge & A same-carrier official counterexample act that defeats the endpoint while not fixing target, carrier, branch, standing, reportability, downstream reuse, branch exclusion, act-time finality, and redescription-stable endpoint role.\\
Context non-smuggling & A clause where \(\CRHEF\) asserts RH, line readout, neutral amplitude occupation, or no-independent closure.\\
EF adequacy & A nontrivial zeta zero contribution omitted by the model, or a contribution not tied to the zeta carrier.\\
Amplitude exactness & A zero \(\rho=\beta+i\gamma\) whose critical-normalized scale is not governed by \(\beta-1/2\).\\
Pairing/cancellation & A proof that symmetry or cancellation makes off-line scale endpoint-neutral without a bridge.\\
Single interface & A lawful occupant of the critical-neutral amplitude endpoint image that is non-neutral.\\
Non-explosiveness & A place where the manuscript derives contradiction from \(\OffAmpExcl\) alone.\\
Non-neutral standing & Endpoint-standing non-neutral amplitude that is support, bookkeeping, carrier shift, repair, surrogate, coequal occupation, and same-carrier endpoint result all at once.\\
Hinge theorem & Noncritical amplitude status work that is not independent endpoint-status discrimination and not any lawful alternative role.\\
A+ closure & \(\Damp(\rho)\) compatible with \(\Namp(\rho)\).\\
Reductio discharge & A proof that the annotation is an added premise or that contradiction discharges only a strengthened assumption.\\
Final endpoint & Failure of arbitrary zerohood closure to imply the universal RH statement.\\
\bottomrule
\end{longtable}

\section{Hostile-Referee Objections and Replies}

\begin{audit}[This does not prove a zero-free region]
Correct.  The proof class is not zero-free-region construction.  A zero-free region could be bridge support; it is not required for endpoint-exclusion proof if amplitude separator exclusion succeeds.
\end{audit}

\begin{audit}[Off-line zeros are analytically meaningful]
Correct.  The proof preserves off-line zerohood as raw analytic syntax and preserves the explicit-formula contribution it would have.  The excluded role is endpoint-standing non-neutral amplitude separator occupation.
\end{audit}

\begin{audit}[The real-part coordinate is native]
Correct.  The proof does not attack the real-part coordinate.  It routes endpoint counterforce through critical-normalized prime-trace amplitude and classifies endpoint-standing non-neutral amplitude as independent discriminator work.
\end{audit}

\begin{audit}[Amplitude exponents are native analytic data]
Correct as descriptive and support content.  They become AASC-relevant only when used as endpoint-standing authority over the fixed critical-neutral amplitude image.
\end{audit}

\begin{audit}[The neutral amplitude image assumes RH]
No.  It defines the positive endpoint image under the official route.  It does not assert universal occupation.  Occupation is derived only after the local countercase is excluded.
\end{audit}

\begin{audit}[A+ is doing analytic work]
No.  The explicit-formula route supplies amplitude projection.  A+ excludes independent endpoint-status discrimination once the role has been identified.  It does not supply the explicit-formula scale exponent or a zero-free estimate.
\end{audit}

\begin{audit}[The proof proves arbitrary propositions]
No.  The anti-trivialization theorem blocks arbitrary \(\neg P\).  The proof requires a fixed endpoint image, official explicit-formula route, projection-role content, and the non-neutral scale standing theorem.
\end{audit}

\section{Load-Bearing Theorem Audit Packets}

This section repeats every load-bearing transition using the full hostile-referee audit pattern: statement, dependency, proof function, anti-misreading, denial burden, and worksheet placement. The repetition is deliberate; no load-bearing theorem is allowed to appear only once.

\newpage

\subsection{Audit Packet 1: Official endpoint fixation}

\textbf{Purpose.} Ensures the manuscript solves the official RH statement rather than an AASC surrogate.\par

\textbf{Statement.} The theorem named in this packet is used only for its declared role. It does not assert the final RH endpoint unless the final universalization theorem explicitly invokes it.\par

\textbf{Dependency.} The theorem depends only on prior target fixation, carrier adequacy, ATS/UEAP locking, explicit-formula projection, projection-role content, or the kernel-forced A+ consequence layer. The final conclusion is not available upstream.\par

\textbf{Proof function.} The theorem fixes a role, prevents report laundering, classifies projected content, or excludes a same-domain independent discriminator. Analytic content is restricted to the explicit-formula scale projection declared by model adequacy.\par

\textbf{Anti-misreading.} This packet must not be read as banning analytic syntax, coordinate data, explicit-formula support, symmetry-pair information, numerical evidence, density estimates, or ordinary negative theoremhood. The excluded object is the endpoint-standing misuse named in the theorem.\par

\textbf{Hostile denial burden.} A successful objection must produce a concrete failure: carrier mismatch, hidden endpoint assumption, projection failure, lawful fifth role, annotation-discharge error, or a discriminator compatible with no-independent closure. Genre objections do not meet this burden.\par

\textbf{Worksheet entry.} The theorem appears in the theorem ladder, anti-circularity audit, route-locus disposition, and faithful counterexample worksheet. If this packet is removed, the manuscript loses a declared dependency and must not use the theorem downstream.\par

\textbf{RH-EF specialization.} In the prime-trace amplitude route, the fixed endpoint image is critical-neutral normalized amplitude. Descriptive non-neutral amplitude remains lawful. Endpoint-standing non-neutral amplitude is the only object classified as independent amplitude discriminator work.\par

\newpage

\subsection{Audit Packet 2: Context non-smuggling}

\textbf{Purpose.} Shows that fixing the zeta/prime-trace carrier does not assert line support.\par

\textbf{Statement.} The theorem named in this packet is used only for its declared role. It does not assert the final RH endpoint unless the final universalization theorem explicitly invokes it.\par

\textbf{Dependency.} The theorem depends only on prior target fixation, carrier adequacy, ATS/UEAP locking, explicit-formula projection, projection-role content, or the kernel-forced A+ consequence layer. The final conclusion is not available upstream.\par

\textbf{Proof function.} The theorem fixes a role, prevents report laundering, classifies projected content, or excludes a same-domain independent discriminator. Analytic content is restricted to the explicit-formula scale projection declared by model adequacy.\par

\textbf{Anti-misreading.} This packet must not be read as banning analytic syntax, coordinate data, explicit-formula support, symmetry-pair information, numerical evidence, density estimates, or ordinary negative theoremhood. The excluded object is the endpoint-standing misuse named in the theorem.\par

\textbf{Hostile denial burden.} A successful objection must produce a concrete failure: carrier mismatch, hidden endpoint assumption, projection failure, lawful fifth role, annotation-discharge error, or a discriminator compatible with no-independent closure. Genre objections do not meet this burden.\par

\textbf{Worksheet entry.} The theorem appears in the theorem ladder, anti-circularity audit, route-locus disposition, and faithful counterexample worksheet. If this packet is removed, the manuscript loses a declared dependency and must not use the theorem downstream.\par

\textbf{RH-EF specialization.} In the prime-trace amplitude route, the fixed endpoint image is critical-neutral normalized amplitude. Descriptive non-neutral amplitude remains lawful. Endpoint-standing non-neutral amplitude is the only object classified as independent amplitude discriminator work.\par

\newpage

\subsection{Audit Packet 3: Official explicit-formula route}

\textbf{Purpose.} Makes the explicit formula the endpoint route rather than a background analogy.\par

\textbf{Statement.} The theorem named in this packet is used only for its declared role. It does not assert the final RH endpoint unless the final universalization theorem explicitly invokes it.\par

\textbf{Dependency.} The theorem depends only on prior target fixation, carrier adequacy, ATS/UEAP locking, explicit-formula projection, projection-role content, or the kernel-forced A+ consequence layer. The final conclusion is not available upstream.\par

\textbf{Proof function.} The theorem fixes a role, prevents report laundering, classifies projected content, or excludes a same-domain independent discriminator. Analytic content is restricted to the explicit-formula scale projection declared by model adequacy.\par

\textbf{Anti-misreading.} This packet must not be read as banning analytic syntax, coordinate data, explicit-formula support, symmetry-pair information, numerical evidence, density estimates, or ordinary negative theoremhood. The excluded object is the endpoint-standing misuse named in the theorem.\par

\textbf{Hostile denial burden.} A successful objection must produce a concrete failure: carrier mismatch, hidden endpoint assumption, projection failure, lawful fifth role, annotation-discharge error, or a discriminator compatible with no-independent closure. Genre objections do not meet this burden.\par

\textbf{Worksheet entry.} The theorem appears in the theorem ladder, anti-circularity audit, route-locus disposition, and faithful counterexample worksheet. If this packet is removed, the manuscript loses a declared dependency and must not use the theorem downstream.\par

\textbf{RH-EF specialization.} In the prime-trace amplitude route, the fixed endpoint image is critical-neutral normalized amplitude. Descriptive non-neutral amplitude remains lawful. Endpoint-standing non-neutral amplitude is the only object classified as independent amplitude discriminator work.\par

\newpage

\subsection{Audit Packet 4: Endpoint-closure evaluation}

\textbf{Purpose.} Separates endpoint-closure proof discipline from endpoint-complete outcome assumption.\par

\textbf{Statement.} The theorem named in this packet is used only for its declared role. It does not assert the final RH endpoint unless the final universalization theorem explicitly invokes it.\par

\textbf{Dependency.} The theorem depends only on prior target fixation, carrier adequacy, ATS/UEAP locking, explicit-formula projection, projection-role content, or the kernel-forced A+ consequence layer. The final conclusion is not available upstream.\par

\textbf{Proof function.} The theorem fixes a role, prevents report laundering, classifies projected content, or excludes a same-domain independent discriminator. Analytic content is restricted to the explicit-formula scale projection declared by model adequacy.\par

\textbf{Anti-misreading.} This packet must not be read as banning analytic syntax, coordinate data, explicit-formula support, symmetry-pair information, numerical evidence, density estimates, or ordinary negative theoremhood. The excluded object is the endpoint-standing misuse named in the theorem.\par

\textbf{Hostile denial burden.} A successful objection must produce a concrete failure: carrier mismatch, hidden endpoint assumption, projection failure, lawful fifth role, annotation-discharge error, or a discriminator compatible with no-independent closure. Genre objections do not meet this burden.\par

\textbf{Worksheet entry.} The theorem appears in the theorem ladder, anti-circularity audit, route-locus disposition, and faithful counterexample worksheet. If this packet is removed, the manuscript loses a declared dependency and must not use the theorem downstream.\par

\textbf{RH-EF specialization.} In the prime-trace amplitude route, the fixed endpoint image is critical-neutral normalized amplitude. Descriptive non-neutral amplitude remains lawful. Endpoint-standing non-neutral amplitude is the only object classified as independent amplitude discriminator work.\par

\newpage

\subsection{Audit Packet 5: Kernel forcing}

\textbf{Purpose.} Derives reference, standing, admissibility, and irreversibility from target adequacy.\par

\textbf{Statement.} The theorem named in this packet is used only for its declared role. It does not assert the final RH endpoint unless the final universalization theorem explicitly invokes it.\par

\textbf{Dependency.} The theorem depends only on prior target fixation, carrier adequacy, ATS/UEAP locking, explicit-formula projection, projection-role content, or the kernel-forced A+ consequence layer. The final conclusion is not available upstream.\par

\textbf{Proof function.} The theorem fixes a role, prevents report laundering, classifies projected content, or excludes a same-domain independent discriminator. Analytic content is restricted to the explicit-formula scale projection declared by model adequacy.\par

\textbf{Anti-misreading.} This packet must not be read as banning analytic syntax, coordinate data, explicit-formula support, symmetry-pair information, numerical evidence, density estimates, or ordinary negative theoremhood. The excluded object is the endpoint-standing misuse named in the theorem.\par

\textbf{Hostile denial burden.} A successful objection must produce a concrete failure: carrier mismatch, hidden endpoint assumption, projection failure, lawful fifth role, annotation-discharge error, or a discriminator compatible with no-independent closure. Genre objections do not meet this burden.\par

\textbf{Worksheet entry.} The theorem appears in the theorem ladder, anti-circularity audit, route-locus disposition, and faithful counterexample worksheet. If this packet is removed, the manuscript loses a declared dependency and must not use the theorem downstream.\par

\textbf{RH-EF specialization.} In the prime-trace amplitude route, the fixed endpoint image is critical-neutral normalized amplitude. Descriptive non-neutral amplitude remains lawful. Endpoint-standing non-neutral amplitude is the only object classified as independent amplitude discriminator work.\par

\newpage

\subsection{Audit Packet 6: A+ consequence layer}

\textbf{Purpose.} Treats AMetricity, no-selector discipline, and no-independent-discriminator closure as consequences.\par

\textbf{Statement.} The theorem named in this packet is used only for its declared role. It does not assert the final RH endpoint unless the final universalization theorem explicitly invokes it.\par

\textbf{Dependency.} The theorem depends only on prior target fixation, carrier adequacy, ATS/UEAP locking, explicit-formula projection, projection-role content, or the kernel-forced A+ consequence layer. The final conclusion is not available upstream.\par

\textbf{Proof function.} The theorem fixes a role, prevents report laundering, classifies projected content, or excludes a same-domain independent discriminator. Analytic content is restricted to the explicit-formula scale projection declared by model adequacy.\par

\textbf{Anti-misreading.} This packet must not be read as banning analytic syntax, coordinate data, explicit-formula support, symmetry-pair information, numerical evidence, density estimates, or ordinary negative theoremhood. The excluded object is the endpoint-standing misuse named in the theorem.\par

\textbf{Hostile denial burden.} A successful objection must produce a concrete failure: carrier mismatch, hidden endpoint assumption, projection failure, lawful fifth role, annotation-discharge error, or a discriminator compatible with no-independent closure. Genre objections do not meet this burden.\par

\textbf{Worksheet entry.} The theorem appears in the theorem ladder, anti-circularity audit, route-locus disposition, and faithful counterexample worksheet. If this packet is removed, the manuscript loses a declared dependency and must not use the theorem downstream.\par

\textbf{RH-EF specialization.} In the prime-trace amplitude route, the fixed endpoint image is critical-neutral normalized amplitude. Descriptive non-neutral amplitude remains lawful. Endpoint-standing non-neutral amplitude is the only object classified as independent amplitude discriminator work.\par

\newpage

\subsection{Audit Packet 7: ATS anchor locking}

\textbf{Purpose.} Fixes the zeta carrier, prime-trace carrier, same-zero identity, and critical normalization.\par

\textbf{Statement.} The theorem named in this packet is used only for its declared role. It does not assert the final RH endpoint unless the final universalization theorem explicitly invokes it.\par

\textbf{Dependency.} The theorem depends only on prior target fixation, carrier adequacy, ATS/UEAP locking, explicit-formula projection, projection-role content, or the kernel-forced A+ consequence layer. The final conclusion is not available upstream.\par

\textbf{Proof function.} The theorem fixes a role, prevents report laundering, classifies projected content, or excludes a same-domain independent discriminator. Analytic content is restricted to the explicit-formula scale projection declared by model adequacy.\par

\textbf{Anti-misreading.} This packet must not be read as banning analytic syntax, coordinate data, explicit-formula support, symmetry-pair information, numerical evidence, density estimates, or ordinary negative theoremhood. The excluded object is the endpoint-standing misuse named in the theorem.\par

\textbf{Hostile denial burden.} A successful objection must produce a concrete failure: carrier mismatch, hidden endpoint assumption, projection failure, lawful fifth role, annotation-discharge error, or a discriminator compatible with no-independent closure. Genre objections do not meet this burden.\par

\textbf{Worksheet entry.} The theorem appears in the theorem ladder, anti-circularity audit, route-locus disposition, and faithful counterexample worksheet. If this packet is removed, the manuscript loses a declared dependency and must not use the theorem downstream.\par

\textbf{RH-EF specialization.} In the prime-trace amplitude route, the fixed endpoint image is critical-neutral normalized amplitude. Descriptive non-neutral amplitude remains lawful. Endpoint-standing non-neutral amplitude is the only object classified as independent amplitude discriminator work.\par

\newpage

\subsection{Audit Packet 8: ATS tensor locking}

\textbf{Purpose.} Treats the normalized exponent beta minus one-half as tensor content, not notation skin.\par

\textbf{Statement.} The theorem named in this packet is used only for its declared role. It does not assert the final RH endpoint unless the final universalization theorem explicitly invokes it.\par

\textbf{Dependency.} The theorem depends only on prior target fixation, carrier adequacy, ATS/UEAP locking, explicit-formula projection, projection-role content, or the kernel-forced A+ consequence layer. The final conclusion is not available upstream.\par

\textbf{Proof function.} The theorem fixes a role, prevents report laundering, classifies projected content, or excludes a same-domain independent discriminator. Analytic content is restricted to the explicit-formula scale projection declared by model adequacy.\par

\textbf{Anti-misreading.} This packet must not be read as banning analytic syntax, coordinate data, explicit-formula support, symmetry-pair information, numerical evidence, density estimates, or ordinary negative theoremhood. The excluded object is the endpoint-standing misuse named in the theorem.\par

\textbf{Hostile denial burden.} A successful objection must produce a concrete failure: carrier mismatch, hidden endpoint assumption, projection failure, lawful fifth role, annotation-discharge error, or a discriminator compatible with no-independent closure. Genre objections do not meet this burden.\par

\textbf{Worksheet entry.} The theorem appears in the theorem ladder, anti-circularity audit, route-locus disposition, and faithful counterexample worksheet. If this packet is removed, the manuscript loses a declared dependency and must not use the theorem downstream.\par

\textbf{RH-EF specialization.} In the prime-trace amplitude route, the fixed endpoint image is critical-neutral normalized amplitude. Descriptive non-neutral amplitude remains lawful. Endpoint-standing non-neutral amplitude is the only object classified as independent amplitude discriminator work.\par

\newpage

\subsection{Audit Packet 9: UEAP report ladder}

\textbf{Purpose.} Prevents laundering from zerohood or support-level amplitude into endpoint status.\par

\textbf{Statement.} The theorem named in this packet is used only for its declared role. It does not assert the final RH endpoint unless the final universalization theorem explicitly invokes it.\par

\textbf{Dependency.} The theorem depends only on prior target fixation, carrier adequacy, ATS/UEAP locking, explicit-formula projection, projection-role content, or the kernel-forced A+ consequence layer. The final conclusion is not available upstream.\par

\textbf{Proof function.} The theorem fixes a role, prevents report laundering, classifies projected content, or excludes a same-domain independent discriminator. Analytic content is restricted to the explicit-formula scale projection declared by model adequacy.\par

\textbf{Anti-misreading.} This packet must not be read as banning analytic syntax, coordinate data, explicit-formula support, symmetry-pair information, numerical evidence, density estimates, or ordinary negative theoremhood. The excluded object is the endpoint-standing misuse named in the theorem.\par

\textbf{Hostile denial burden.} A successful objection must produce a concrete failure: carrier mismatch, hidden endpoint assumption, projection failure, lawful fifth role, annotation-discharge error, or a discriminator compatible with no-independent closure. Genre objections do not meet this burden.\par

\textbf{Worksheet entry.} The theorem appears in the theorem ladder, anti-circularity audit, route-locus disposition, and faithful counterexample worksheet. If this packet is removed, the manuscript loses a declared dependency and must not use the theorem downstream.\par

\textbf{RH-EF specialization.} In the prime-trace amplitude route, the fixed endpoint image is critical-neutral normalized amplitude. Descriptive non-neutral amplitude remains lawful. Endpoint-standing non-neutral amplitude is the only object classified as independent amplitude discriminator work.\par

\newpage

\subsection{Audit Packet 10: Model adequacy}

\textbf{Purpose.} Pins the explicit-formula carrier to the official zeta zerohood carrier.\par

\textbf{Statement.} The theorem named in this packet is used only for its declared role. It does not assert the final RH endpoint unless the final universalization theorem explicitly invokes it.\par

\textbf{Dependency.} The theorem depends only on prior target fixation, carrier adequacy, ATS/UEAP locking, explicit-formula projection, projection-role content, or the kernel-forced A+ consequence layer. The final conclusion is not available upstream.\par

\textbf{Proof function.} The theorem fixes a role, prevents report laundering, classifies projected content, or excludes a same-domain independent discriminator. Analytic content is restricted to the explicit-formula scale projection declared by model adequacy.\par

\textbf{Anti-misreading.} This packet must not be read as banning analytic syntax, coordinate data, explicit-formula support, symmetry-pair information, numerical evidence, density estimates, or ordinary negative theoremhood. The excluded object is the endpoint-standing misuse named in the theorem.\par

\textbf{Hostile denial burden.} A successful objection must produce a concrete failure: carrier mismatch, hidden endpoint assumption, projection failure, lawful fifth role, annotation-discharge error, or a discriminator compatible with no-independent closure. Genre objections do not meet this burden.\par

\textbf{Worksheet entry.} The theorem appears in the theorem ladder, anti-circularity audit, route-locus disposition, and faithful counterexample worksheet. If this packet is removed, the manuscript loses a declared dependency and must not use the theorem downstream.\par

\textbf{RH-EF specialization.} In the prime-trace amplitude route, the fixed endpoint image is critical-neutral normalized amplitude. Descriptive non-neutral amplitude remains lawful. Endpoint-standing non-neutral amplitude is the only object classified as independent amplitude discriminator work.\par

\newpage

\subsection{Audit Packet 11: Term isolation}

\textbf{Purpose.} Allows individual zero contribution scale classification without claiming dominance.\par

\textbf{Statement.} The theorem named in this packet is used only for its declared role. It does not assert the final RH endpoint unless the final universalization theorem explicitly invokes it.\par

\textbf{Dependency.} The theorem depends only on prior target fixation, carrier adequacy, ATS/UEAP locking, explicit-formula projection, projection-role content, or the kernel-forced A+ consequence layer. The final conclusion is not available upstream.\par

\textbf{Proof function.} The theorem fixes a role, prevents report laundering, classifies projected content, or excludes a same-domain independent discriminator. Analytic content is restricted to the explicit-formula scale projection declared by model adequacy.\par

\textbf{Anti-misreading.} This packet must not be read as banning analytic syntax, coordinate data, explicit-formula support, symmetry-pair information, numerical evidence, density estimates, or ordinary negative theoremhood. The excluded object is the endpoint-standing misuse named in the theorem.\par

\textbf{Hostile denial burden.} A successful objection must produce a concrete failure: carrier mismatch, hidden endpoint assumption, projection failure, lawful fifth role, annotation-discharge error, or a discriminator compatible with no-independent closure. Genre objections do not meet this burden.\par

\textbf{Worksheet entry.} The theorem appears in the theorem ladder, anti-circularity audit, route-locus disposition, and faithful counterexample worksheet. If this packet is removed, the manuscript loses a declared dependency and must not use the theorem downstream.\par

\textbf{RH-EF specialization.} In the prime-trace amplitude route, the fixed endpoint image is critical-neutral normalized amplitude. Descriptive non-neutral amplitude remains lawful. Endpoint-standing non-neutral amplitude is the only object classified as independent amplitude discriminator work.\par

\newpage

\subsection{Audit Packet 12: Pairing and cancellation}

\textbf{Purpose.} Allows symmetry-paired terms while blocking unproved endpoint-neutral cancellation.\par

\textbf{Statement.} The theorem named in this packet is used only for its declared role. It does not assert the final RH endpoint unless the final universalization theorem explicitly invokes it.\par

\textbf{Dependency.} The theorem depends only on prior target fixation, carrier adequacy, ATS/UEAP locking, explicit-formula projection, projection-role content, or the kernel-forced A+ consequence layer. The final conclusion is not available upstream.\par

\textbf{Proof function.} The theorem fixes a role, prevents report laundering, classifies projected content, or excludes a same-domain independent discriminator. Analytic content is restricted to the explicit-formula scale projection declared by model adequacy.\par

\textbf{Anti-misreading.} This packet must not be read as banning analytic syntax, coordinate data, explicit-formula support, symmetry-pair information, numerical evidence, density estimates, or ordinary negative theoremhood. The excluded object is the endpoint-standing misuse named in the theorem.\par

\textbf{Hostile denial burden.} A successful objection must produce a concrete failure: carrier mismatch, hidden endpoint assumption, projection failure, lawful fifth role, annotation-discharge error, or a discriminator compatible with no-independent closure. Genre objections do not meet this burden.\par

\textbf{Worksheet entry.} The theorem appears in the theorem ladder, anti-circularity audit, route-locus disposition, and faithful counterexample worksheet. If this packet is removed, the manuscript loses a declared dependency and must not use the theorem downstream.\par

\textbf{RH-EF specialization.} In the prime-trace amplitude route, the fixed endpoint image is critical-neutral normalized amplitude. Descriptive non-neutral amplitude remains lawful. Endpoint-standing non-neutral amplitude is the only object classified as independent amplitude discriminator work.\par

\newpage

\subsection{Audit Packet 13: Positive projection}

\textbf{Purpose.} Identifies critical-line readout with neutral normalized amplitude occupation.\par

\textbf{Statement.} The theorem named in this packet is used only for its declared role. It does not assert the final RH endpoint unless the final universalization theorem explicitly invokes it.\par

\textbf{Dependency.} The theorem depends only on prior target fixation, carrier adequacy, ATS/UEAP locking, explicit-formula projection, projection-role content, or the kernel-forced A+ consequence layer. The final conclusion is not available upstream.\par

\textbf{Proof function.} The theorem fixes a role, prevents report laundering, classifies projected content, or excludes a same-domain independent discriminator. Analytic content is restricted to the explicit-formula scale projection declared by model adequacy.\par

\textbf{Anti-misreading.} This packet must not be read as banning analytic syntax, coordinate data, explicit-formula support, symmetry-pair information, numerical evidence, density estimates, or ordinary negative theoremhood. The excluded object is the endpoint-standing misuse named in the theorem.\par

\textbf{Hostile denial burden.} A successful objection must produce a concrete failure: carrier mismatch, hidden endpoint assumption, projection failure, lawful fifth role, annotation-discharge error, or a discriminator compatible with no-independent closure. Genre objections do not meet this burden.\par

\textbf{Worksheet entry.} The theorem appears in the theorem ladder, anti-circularity audit, route-locus disposition, and faithful counterexample worksheet. If this packet is removed, the manuscript loses a declared dependency and must not use the theorem downstream.\par

\textbf{RH-EF specialization.} In the prime-trace amplitude route, the fixed endpoint image is critical-neutral normalized amplitude. Descriptive non-neutral amplitude remains lawful. Endpoint-standing non-neutral amplitude is the only object classified as independent amplitude discriminator work.\par

\newpage

\subsection{Audit Packet 14: Negative projection}

\textbf{Purpose.} Identifies off-line readout with non-neutral amplitude non-occupation.\par

\textbf{Statement.} The theorem named in this packet is used only for its declared role. It does not assert the final RH endpoint unless the final universalization theorem explicitly invokes it.\par

\textbf{Dependency.} The theorem depends only on prior target fixation, carrier adequacy, ATS/UEAP locking, explicit-formula projection, projection-role content, or the kernel-forced A+ consequence layer. The final conclusion is not available upstream.\par

\textbf{Proof function.} The theorem fixes a role, prevents report laundering, classifies projected content, or excludes a same-domain independent discriminator. Analytic content is restricted to the explicit-formula scale projection declared by model adequacy.\par

\textbf{Anti-misreading.} This packet must not be read as banning analytic syntax, coordinate data, explicit-formula support, symmetry-pair information, numerical evidence, density estimates, or ordinary negative theoremhood. The excluded object is the endpoint-standing misuse named in the theorem.\par

\textbf{Hostile denial burden.} A successful objection must produce a concrete failure: carrier mismatch, hidden endpoint assumption, projection failure, lawful fifth role, annotation-discharge error, or a discriminator compatible with no-independent closure. Genre objections do not meet this burden.\par

\textbf{Worksheet entry.} The theorem appears in the theorem ladder, anti-circularity audit, route-locus disposition, and faithful counterexample worksheet. If this packet is removed, the manuscript loses a declared dependency and must not use the theorem downstream.\par

\textbf{RH-EF specialization.} In the prime-trace amplitude route, the fixed endpoint image is critical-neutral normalized amplitude. Descriptive non-neutral amplitude remains lawful. Endpoint-standing non-neutral amplitude is the only object classified as independent amplitude discriminator work.\par

\newpage

\subsection{Audit Packet 15: Projection-role content}

\textbf{Purpose.} Prevents jumping from coordinate negation directly to separator occupation.\par

\textbf{Statement.} The theorem named in this packet is used only for its declared role. It does not assert the final RH endpoint unless the final universalization theorem explicitly invokes it.\par

\textbf{Dependency.} The theorem depends only on prior target fixation, carrier adequacy, ATS/UEAP locking, explicit-formula projection, projection-role content, or the kernel-forced A+ consequence layer. The final conclusion is not available upstream.\par

\textbf{Proof function.} The theorem fixes a role, prevents report laundering, classifies projected content, or excludes a same-domain independent discriminator. Analytic content is restricted to the explicit-formula scale projection declared by model adequacy.\par

\textbf{Anti-misreading.} This packet must not be read as banning analytic syntax, coordinate data, explicit-formula support, symmetry-pair information, numerical evidence, density estimates, or ordinary negative theoremhood. The excluded object is the endpoint-standing misuse named in the theorem.\par

\textbf{Hostile denial burden.} A successful objection must produce a concrete failure: carrier mismatch, hidden endpoint assumption, projection failure, lawful fifth role, annotation-discharge error, or a discriminator compatible with no-independent closure. Genre objections do not meet this burden.\par

\textbf{Worksheet entry.} The theorem appears in the theorem ladder, anti-circularity audit, route-locus disposition, and faithful counterexample worksheet. If this packet is removed, the manuscript loses a declared dependency and must not use the theorem downstream.\par

\textbf{RH-EF specialization.} In the prime-trace amplitude route, the fixed endpoint image is critical-neutral normalized amplitude. Descriptive non-neutral amplitude remains lawful. Endpoint-standing non-neutral amplitude is the only object classified as independent amplitude discriminator work.\par

\newpage

\subsection{Audit Packet 16: Single amplitude interface}

\textbf{Purpose.} Distinguishes the broad RH/neg-RH outcome carrier from the critical-neutral endpoint image.\par

\textbf{Statement.} The theorem named in this packet is used only for its declared role. It does not assert the final RH endpoint unless the final universalization theorem explicitly invokes it.\par

\textbf{Dependency.} The theorem depends only on prior target fixation, carrier adequacy, ATS/UEAP locking, explicit-formula projection, projection-role content, or the kernel-forced A+ consequence layer. The final conclusion is not available upstream.\par

\textbf{Proof function.} The theorem fixes a role, prevents report laundering, classifies projected content, or excludes a same-domain independent discriminator. Analytic content is restricted to the explicit-formula scale projection declared by model adequacy.\par

\textbf{Anti-misreading.} This packet must not be read as banning analytic syntax, coordinate data, explicit-formula support, symmetry-pair information, numerical evidence, density estimates, or ordinary negative theoremhood. The excluded object is the endpoint-standing misuse named in the theorem.\par

\textbf{Hostile denial burden.} A successful objection must produce a concrete failure: carrier mismatch, hidden endpoint assumption, projection failure, lawful fifth role, annotation-discharge error, or a discriminator compatible with no-independent closure. Genre objections do not meet this burden.\par

\textbf{Worksheet entry.} The theorem appears in the theorem ladder, anti-circularity audit, route-locus disposition, and faithful counterexample worksheet. If this packet is removed, the manuscript loses a declared dependency and must not use the theorem downstream.\par

\textbf{RH-EF specialization.} In the prime-trace amplitude route, the fixed endpoint image is critical-neutral normalized amplitude. Descriptive non-neutral amplitude remains lawful. Endpoint-standing non-neutral amplitude is the only object classified as independent amplitude discriminator work.\par

\newpage

\subsection{Audit Packet 17: Non-explosiveness}

\textbf{Purpose.} Proves that off-line amplitude displacement alone is not contradiction.\par

\textbf{Statement.} The theorem named in this packet is used only for its declared role. It does not assert the final RH endpoint unless the final universalization theorem explicitly invokes it.\par

\textbf{Dependency.} The theorem depends only on prior target fixation, carrier adequacy, ATS/UEAP locking, explicit-formula projection, projection-role content, or the kernel-forced A+ consequence layer. The final conclusion is not available upstream.\par

\textbf{Proof function.} The theorem fixes a role, prevents report laundering, classifies projected content, or excludes a same-domain independent discriminator. Analytic content is restricted to the explicit-formula scale projection declared by model adequacy.\par

\textbf{Anti-misreading.} This packet must not be read as banning analytic syntax, coordinate data, explicit-formula support, symmetry-pair information, numerical evidence, density estimates, or ordinary negative theoremhood. The excluded object is the endpoint-standing misuse named in the theorem.\par

\textbf{Hostile denial burden.} A successful objection must produce a concrete failure: carrier mismatch, hidden endpoint assumption, projection failure, lawful fifth role, annotation-discharge error, or a discriminator compatible with no-independent closure. Genre objections do not meet this burden.\par

\textbf{Worksheet entry.} The theorem appears in the theorem ladder, anti-circularity audit, route-locus disposition, and faithful counterexample worksheet. If this packet is removed, the manuscript loses a declared dependency and must not use the theorem downstream.\par

\textbf{RH-EF specialization.} In the prime-trace amplitude route, the fixed endpoint image is critical-neutral normalized amplitude. Descriptive non-neutral amplitude remains lawful. Endpoint-standing non-neutral amplitude is the only object classified as independent amplitude discriminator work.\par

\newpage

\subsection{Audit Packet 18: Descriptive amplitude preservation}

\textbf{Purpose.} Allows the scale \(x^{\beta-1/2}\) as ordinary explicit-formula content.\par

\textbf{Statement.} The theorem named in this packet is used only for its declared role. It does not assert the final RH endpoint unless the final universalization theorem explicitly invokes it.\par

\textbf{Dependency.} The theorem depends only on prior target fixation, carrier adequacy, ATS/UEAP locking, explicit-formula projection, projection-role content, or the kernel-forced A+ consequence layer. The final conclusion is not available upstream.\par

\textbf{Proof function.} The theorem fixes a role, prevents report laundering, classifies projected content, or excludes a same-domain independent discriminator. Analytic content is restricted to the explicit-formula scale projection declared by model adequacy.\par

\textbf{Anti-misreading.} This packet must not be read as banning analytic syntax, coordinate data, explicit-formula support, symmetry-pair information, numerical evidence, density estimates, or ordinary negative theoremhood. The excluded object is the endpoint-standing misuse named in the theorem.\par

\textbf{Hostile denial burden.} A successful objection must produce a concrete failure: carrier mismatch, hidden endpoint assumption, projection failure, lawful fifth role, annotation-discharge error, or a discriminator compatible with no-independent closure. Genre objections do not meet this burden.\par

\textbf{Worksheet entry.} The theorem appears in the theorem ladder, anti-circularity audit, route-locus disposition, and faithful counterexample worksheet. If this packet is removed, the manuscript loses a declared dependency and must not use the theorem downstream.\par

\textbf{RH-EF specialization.} In the prime-trace amplitude route, the fixed endpoint image is critical-neutral normalized amplitude. Descriptive non-neutral amplitude remains lawful. Endpoint-standing non-neutral amplitude is the only object classified as independent amplitude discriminator work.\par

\newpage

\subsection{Audit Packet 19: Endpoint-standing amplitude}

\textbf{Purpose.} Marks the higher role in which non-neutral amplitude is used to defeat endpoint closure.\par

\textbf{Statement.} The theorem named in this packet is used only for its declared role. It does not assert the final RH endpoint unless the final universalization theorem explicitly invokes it.\par

\textbf{Dependency.} The theorem depends only on prior target fixation, carrier adequacy, ATS/UEAP locking, explicit-formula projection, projection-role content, or the kernel-forced A+ consequence layer. The final conclusion is not available upstream.\par

\textbf{Proof function.} The theorem fixes a role, prevents report laundering, classifies projected content, or excludes a same-domain independent discriminator. Analytic content is restricted to the explicit-formula scale projection declared by model adequacy.\par

\textbf{Anti-misreading.} This packet must not be read as banning analytic syntax, coordinate data, explicit-formula support, symmetry-pair information, numerical evidence, density estimates, or ordinary negative theoremhood. The excluded object is the endpoint-standing misuse named in the theorem.\par

\textbf{Hostile denial burden.} A successful objection must produce a concrete failure: carrier mismatch, hidden endpoint assumption, projection failure, lawful fifth role, annotation-discharge error, or a discriminator compatible with no-independent closure. Genre objections do not meet this burden.\par

\textbf{Worksheet entry.} The theorem appears in the theorem ladder, anti-circularity audit, route-locus disposition, and faithful counterexample worksheet. If this packet is removed, the manuscript loses a declared dependency and must not use the theorem downstream.\par

\textbf{RH-EF specialization.} In the prime-trace amplitude route, the fixed endpoint image is critical-neutral normalized amplitude. Descriptive non-neutral amplitude remains lawful. Endpoint-standing non-neutral amplitude is the only object classified as independent amplitude discriminator work.\par

\newpage

\subsection{Audit Packet 20: Non-neutral scale standing}

\textbf{Purpose.} Shows endpoint-standing non-neutral amplitude is noncritical amplitude status work.\par

\textbf{Statement.} The theorem named in this packet is used only for its declared role. It does not assert the final RH endpoint unless the final universalization theorem explicitly invokes it.\par

\textbf{Dependency.} The theorem depends only on prior target fixation, carrier adequacy, ATS/UEAP locking, explicit-formula projection, projection-role content, or the kernel-forced A+ consequence layer. The final conclusion is not available upstream.\par

\textbf{Proof function.} The theorem fixes a role, prevents report laundering, classifies projected content, or excludes a same-domain independent discriminator. Analytic content is restricted to the explicit-formula scale projection declared by model adequacy.\par

\textbf{Anti-misreading.} This packet must not be read as banning analytic syntax, coordinate data, explicit-formula support, symmetry-pair information, numerical evidence, density estimates, or ordinary negative theoremhood. The excluded object is the endpoint-standing misuse named in the theorem.\par

\textbf{Hostile denial burden.} A successful objection must produce a concrete failure: carrier mismatch, hidden endpoint assumption, projection failure, lawful fifth role, annotation-discharge error, or a discriminator compatible with no-independent closure. Genre objections do not meet this burden.\par

\textbf{Worksheet entry.} The theorem appears in the theorem ladder, anti-circularity audit, route-locus disposition, and faithful counterexample worksheet. If this packet is removed, the manuscript loses a declared dependency and must not use the theorem downstream.\par

\textbf{RH-EF specialization.} In the prime-trace amplitude route, the fixed endpoint image is critical-neutral normalized amplitude. Descriptive non-neutral amplitude remains lawful. Endpoint-standing non-neutral amplitude is the only object classified as independent amplitude discriminator work.\par

\newpage

\subsection{Audit Packet 21: Non-native amplitude theorem}

\textbf{Purpose.} Classifies endpoint-standing non-neutral amplitude as non-native endpoint authority.\par

\textbf{Statement.} The theorem named in this packet is used only for its declared role. It does not assert the final RH endpoint unless the final universalization theorem explicitly invokes it.\par

\textbf{Dependency.} The theorem depends only on prior target fixation, carrier adequacy, ATS/UEAP locking, explicit-formula projection, projection-role content, or the kernel-forced A+ consequence layer. The final conclusion is not available upstream.\par

\textbf{Proof function.} The theorem fixes a role, prevents report laundering, classifies projected content, or excludes a same-domain independent discriminator. Analytic content is restricted to the explicit-formula scale projection declared by model adequacy.\par

\textbf{Anti-misreading.} This packet must not be read as banning analytic syntax, coordinate data, explicit-formula support, symmetry-pair information, numerical evidence, density estimates, or ordinary negative theoremhood. The excluded object is the endpoint-standing misuse named in the theorem.\par

\textbf{Hostile denial burden.} A successful objection must produce a concrete failure: carrier mismatch, hidden endpoint assumption, projection failure, lawful fifth role, annotation-discharge error, or a discriminator compatible with no-independent closure. Genre objections do not meet this burden.\par

\textbf{Worksheet entry.} The theorem appears in the theorem ladder, anti-circularity audit, route-locus disposition, and faithful counterexample worksheet. If this packet is removed, the manuscript loses a declared dependency and must not use the theorem downstream.\par

\textbf{RH-EF specialization.} In the prime-trace amplitude route, the fixed endpoint image is critical-neutral normalized amplitude. Descriptive non-neutral amplitude remains lawful. Endpoint-standing non-neutral amplitude is the only object classified as independent amplitude discriminator work.\par

\newpage

\subsection{Audit Packet 22: Endpoint-discriminator trichotomy}

\textbf{Purpose.} Exhausts classifiers into bookkeeping, equivalent, independent, or domain-shifting roles.\par

\textbf{Statement.} The theorem named in this packet is used only for its declared role. It does not assert the final RH endpoint unless the final universalization theorem explicitly invokes it.\par

\textbf{Dependency.} The theorem depends only on prior target fixation, carrier adequacy, ATS/UEAP locking, explicit-formula projection, projection-role content, or the kernel-forced A+ consequence layer. The final conclusion is not available upstream.\par

\textbf{Proof function.} The theorem fixes a role, prevents report laundering, classifies projected content, or excludes a same-domain independent discriminator. Analytic content is restricted to the explicit-formula scale projection declared by model adequacy.\par

\textbf{Anti-misreading.} This packet must not be read as banning analytic syntax, coordinate data, explicit-formula support, symmetry-pair information, numerical evidence, density estimates, or ordinary negative theoremhood. The excluded object is the endpoint-standing misuse named in the theorem.\par

\textbf{Hostile denial burden.} A successful objection must produce a concrete failure: carrier mismatch, hidden endpoint assumption, projection failure, lawful fifth role, annotation-discharge error, or a discriminator compatible with no-independent closure. Genre objections do not meet this burden.\par

\textbf{Worksheet entry.} The theorem appears in the theorem ladder, anti-circularity audit, route-locus disposition, and faithful counterexample worksheet. If this packet is removed, the manuscript loses a declared dependency and must not use the theorem downstream.\par

\textbf{RH-EF specialization.} In the prime-trace amplitude route, the fixed endpoint image is critical-neutral normalized amplitude. Descriptive non-neutral amplitude remains lawful. Endpoint-standing non-neutral amplitude is the only object classified as independent amplitude discriminator work.\par

\newpage

\subsection{Audit Packet 23: Hinge theorem}

\textbf{Purpose.} Identifies noncritical amplitude status work with \(D_{\mathrm{amp}}\).\par

\textbf{Statement.} The theorem named in this packet is used only for its declared role. It does not assert the final RH endpoint unless the final universalization theorem explicitly invokes it.\par

\textbf{Dependency.} The theorem depends only on prior target fixation, carrier adequacy, ATS/UEAP locking, explicit-formula projection, projection-role content, or the kernel-forced A+ consequence layer. The final conclusion is not available upstream.\par

\textbf{Proof function.} The theorem fixes a role, prevents report laundering, classifies projected content, or excludes a same-domain independent discriminator. Analytic content is restricted to the explicit-formula scale projection declared by model adequacy.\par

\textbf{Anti-misreading.} This packet must not be read as banning analytic syntax, coordinate data, explicit-formula support, symmetry-pair information, numerical evidence, density estimates, or ordinary negative theoremhood. The excluded object is the endpoint-standing misuse named in the theorem.\par

\textbf{Hostile denial burden.} A successful objection must produce a concrete failure: carrier mismatch, hidden endpoint assumption, projection failure, lawful fifth role, annotation-discharge error, or a discriminator compatible with no-independent closure. Genre objections do not meet this burden.\par

\textbf{Worksheet entry.} The theorem appears in the theorem ladder, anti-circularity audit, route-locus disposition, and faithful counterexample worksheet. If this packet is removed, the manuscript loses a declared dependency and must not use the theorem downstream.\par

\textbf{RH-EF specialization.} In the prime-trace amplitude route, the fixed endpoint image is critical-neutral normalized amplitude. Descriptive non-neutral amplitude remains lawful. Endpoint-standing non-neutral amplitude is the only object classified as independent amplitude discriminator work.\par

\newpage

\subsection{Audit Packet 24: No-independent amplitude closure}

\textbf{Purpose.} Derives \(N_{\mathrm{amp}}\) from endpoint use and fixed-domain A+.\par

\textbf{Statement.} The theorem named in this packet is used only for its declared role. It does not assert the final RH endpoint unless the final universalization theorem explicitly invokes it.\par

\textbf{Dependency.} The theorem depends only on prior target fixation, carrier adequacy, ATS/UEAP locking, explicit-formula projection, projection-role content, or the kernel-forced A+ consequence layer. The final conclusion is not available upstream.\par

\textbf{Proof function.} The theorem fixes a role, prevents report laundering, classifies projected content, or excludes a same-domain independent discriminator. Analytic content is restricted to the explicit-formula scale projection declared by model adequacy.\par

\textbf{Anti-misreading.} This packet must not be read as banning analytic syntax, coordinate data, explicit-formula support, symmetry-pair information, numerical evidence, density estimates, or ordinary negative theoremhood. The excluded object is the endpoint-standing misuse named in the theorem.\par

\textbf{Hostile denial burden.} A successful objection must produce a concrete failure: carrier mismatch, hidden endpoint assumption, projection failure, lawful fifth role, annotation-discharge error, or a discriminator compatible with no-independent closure. Genre objections do not meet this burden.\par

\textbf{Worksheet entry.} The theorem appears in the theorem ladder, anti-circularity audit, route-locus disposition, and faithful counterexample worksheet. If this packet is removed, the manuscript loses a declared dependency and must not use the theorem downstream.\par

\textbf{RH-EF specialization.} In the prime-trace amplitude route, the fixed endpoint image is critical-neutral normalized amplitude. Descriptive non-neutral amplitude remains lawful. Endpoint-standing non-neutral amplitude is the only object classified as independent amplitude discriminator work.\par

\newpage

\subsection{Audit Packet 25: Local amplitude separator exclusion}

\textbf{Purpose.} Combines \(D_{\mathrm{amp}}\) and not-\(D_{\mathrm{amp}}\) to close the separator route.\par

\textbf{Statement.} The theorem named in this packet is used only for its declared role. It does not assert the final RH endpoint unless the final universalization theorem explicitly invokes it.\par

\textbf{Dependency.} The theorem depends only on prior target fixation, carrier adequacy, ATS/UEAP locking, explicit-formula projection, projection-role content, or the kernel-forced A+ consequence layer. The final conclusion is not available upstream.\par

\textbf{Proof function.} The theorem fixes a role, prevents report laundering, classifies projected content, or excludes a same-domain independent discriminator. Analytic content is restricted to the explicit-formula scale projection declared by model adequacy.\par

\textbf{Anti-misreading.} This packet must not be read as banning analytic syntax, coordinate data, explicit-formula support, symmetry-pair information, numerical evidence, density estimates, or ordinary negative theoremhood. The excluded object is the endpoint-standing misuse named in the theorem.\par

\textbf{Hostile denial burden.} A successful objection must produce a concrete failure: carrier mismatch, hidden endpoint assumption, projection failure, lawful fifth role, annotation-discharge error, or a discriminator compatible with no-independent closure. Genre objections do not meet this burden.\par

\textbf{Worksheet entry.} The theorem appears in the theorem ladder, anti-circularity audit, route-locus disposition, and faithful counterexample worksheet. If this packet is removed, the manuscript loses a declared dependency and must not use the theorem downstream.\par

\textbf{RH-EF specialization.} In the prime-trace amplitude route, the fixed endpoint image is critical-neutral normalized amplitude. Descriptive non-neutral amplitude remains lawful. Endpoint-standing non-neutral amplitude is the only object classified as independent amplitude discriminator work.\par

\newpage

\subsection{Audit Packet 26: Anti-trivialization}

\textbf{Purpose.} Blocks arbitrary negation from becoming separator occupation.\par

\textbf{Statement.} The theorem named in this packet is used only for its declared role. It does not assert the final RH endpoint unless the final universalization theorem explicitly invokes it.\par

\textbf{Dependency.} The theorem depends only on prior target fixation, carrier adequacy, ATS/UEAP locking, explicit-formula projection, projection-role content, or the kernel-forced A+ consequence layer. The final conclusion is not available upstream.\par

\textbf{Proof function.} The theorem fixes a role, prevents report laundering, classifies projected content, or excludes a same-domain independent discriminator. Analytic content is restricted to the explicit-formula scale projection declared by model adequacy.\par

\textbf{Anti-misreading.} This packet must not be read as banning analytic syntax, coordinate data, explicit-formula support, symmetry-pair information, numerical evidence, density estimates, or ordinary negative theoremhood. The excluded object is the endpoint-standing misuse named in the theorem.\par

\textbf{Hostile denial burden.} A successful objection must produce a concrete failure: carrier mismatch, hidden endpoint assumption, projection failure, lawful fifth role, annotation-discharge error, or a discriminator compatible with no-independent closure. Genre objections do not meet this burden.\par

\textbf{Worksheet entry.} The theorem appears in the theorem ladder, anti-circularity audit, route-locus disposition, and faithful counterexample worksheet. If this packet is removed, the manuscript loses a declared dependency and must not use the theorem downstream.\par

\textbf{RH-EF specialization.} In the prime-trace amplitude route, the fixed endpoint image is critical-neutral normalized amplitude. Descriptive non-neutral amplitude remains lawful. Endpoint-standing non-neutral amplitude is the only object classified as independent amplitude discriminator work.\par

\newpage

\subsection{Audit Packet 27: Exact-complement annotation}

\textbf{Purpose.} Keeps the reductio annotation outside the object-level premise set.\par

\textbf{Statement.} The theorem named in this packet is used only for its declared role. It does not assert the final RH endpoint unless the final universalization theorem explicitly invokes it.\par

\textbf{Dependency.} The theorem depends only on prior target fixation, carrier adequacy, ATS/UEAP locking, explicit-formula projection, projection-role content, or the kernel-forced A+ consequence layer. The final conclusion is not available upstream.\par

\textbf{Proof function.} The theorem fixes a role, prevents report laundering, classifies projected content, or excludes a same-domain independent discriminator. Analytic content is restricted to the explicit-formula scale projection declared by model adequacy.\par

\textbf{Anti-misreading.} This packet must not be read as banning analytic syntax, coordinate data, explicit-formula support, symmetry-pair information, numerical evidence, density estimates, or ordinary negative theoremhood. The excluded object is the endpoint-standing misuse named in the theorem.\par

\textbf{Hostile denial burden.} A successful objection must produce a concrete failure: carrier mismatch, hidden endpoint assumption, projection failure, lawful fifth role, annotation-discharge error, or a discriminator compatible with no-independent closure. Genre objections do not meet this burden.\par

\textbf{Worksheet entry.} The theorem appears in the theorem ladder, anti-circularity audit, route-locus disposition, and faithful counterexample worksheet. If this packet is removed, the manuscript loses a declared dependency and must not use the theorem downstream.\par

\textbf{RH-EF specialization.} In the prime-trace amplitude route, the fixed endpoint image is critical-neutral normalized amplitude. Descriptive non-neutral amplitude remains lawful. Endpoint-standing non-neutral amplitude is the only object classified as independent amplitude discriminator work.\par

\newpage

\subsection{Audit Packet 28: Local counterforce}

\textbf{Purpose.} Explains why the live local off-line branch has endpoint-facing force.\par

\textbf{Statement.} The theorem named in this packet is used only for its declared role. It does not assert the final RH endpoint unless the final universalization theorem explicitly invokes it.\par

\textbf{Dependency.} The theorem depends only on prior target fixation, carrier adequacy, ATS/UEAP locking, explicit-formula projection, projection-role content, or the kernel-forced A+ consequence layer. The final conclusion is not available upstream.\par

\textbf{Proof function.} The theorem fixes a role, prevents report laundering, classifies projected content, or excludes a same-domain independent discriminator. Analytic content is restricted to the explicit-formula scale projection declared by model adequacy.\par

\textbf{Anti-misreading.} This packet must not be read as banning analytic syntax, coordinate data, explicit-formula support, symmetry-pair information, numerical evidence, density estimates, or ordinary negative theoremhood. The excluded object is the endpoint-standing misuse named in the theorem.\par

\textbf{Hostile denial burden.} A successful objection must produce a concrete failure: carrier mismatch, hidden endpoint assumption, projection failure, lawful fifth role, annotation-discharge error, or a discriminator compatible with no-independent closure. Genre objections do not meet this burden.\par

\textbf{Worksheet entry.} The theorem appears in the theorem ladder, anti-circularity audit, route-locus disposition, and faithful counterexample worksheet. If this packet is removed, the manuscript loses a declared dependency and must not use the theorem downstream.\par

\textbf{RH-EF specialization.} In the prime-trace amplitude route, the fixed endpoint image is critical-neutral normalized amplitude. Descriptive non-neutral amplitude remains lawful. Endpoint-standing non-neutral amplitude is the only object classified as independent amplitude discriminator work.\par

\newpage

\subsection{Audit Packet 29: Annotation discharge}

\textbf{Purpose.} Discharges the original local assumption and not a strengthened conjunction.\par

\textbf{Statement.} The theorem named in this packet is used only for its declared role. It does not assert the final RH endpoint unless the final universalization theorem explicitly invokes it.\par

\textbf{Dependency.} The theorem depends only on prior target fixation, carrier adequacy, ATS/UEAP locking, explicit-formula projection, projection-role content, or the kernel-forced A+ consequence layer. The final conclusion is not available upstream.\par

\textbf{Proof function.} The theorem fixes a role, prevents report laundering, classifies projected content, or excludes a same-domain independent discriminator. Analytic content is restricted to the explicit-formula scale projection declared by model adequacy.\par

\textbf{Anti-misreading.} This packet must not be read as banning analytic syntax, coordinate data, explicit-formula support, symmetry-pair information, numerical evidence, density estimates, or ordinary negative theoremhood. The excluded object is the endpoint-standing misuse named in the theorem.\par

\textbf{Hostile denial burden.} A successful objection must produce a concrete failure: carrier mismatch, hidden endpoint assumption, projection failure, lawful fifth role, annotation-discharge error, or a discriminator compatible with no-independent closure. Genre objections do not meet this burden.\par

\textbf{Worksheet entry.} The theorem appears in the theorem ladder, anti-circularity audit, route-locus disposition, and faithful counterexample worksheet. If this packet is removed, the manuscript loses a declared dependency and must not use the theorem downstream.\par

\textbf{RH-EF specialization.} In the prime-trace amplitude route, the fixed endpoint image is critical-neutral normalized amplitude. Descriptive non-neutral amplitude remains lawful. Endpoint-standing non-neutral amplitude is the only object classified as independent amplitude discriminator work.\par

\newpage

\subsection{Audit Packet 30: Ordinary negative theoremhood}

\textbf{Purpose.} Permits negative theoremhood as theoremhood but role-classifies official use.\par

\textbf{Statement.} The theorem named in this packet is used only for its declared role. It does not assert the final RH endpoint unless the final universalization theorem explicitly invokes it.\par

\textbf{Dependency.} The theorem depends only on prior target fixation, carrier adequacy, ATS/UEAP locking, explicit-formula projection, projection-role content, or the kernel-forced A+ consequence layer. The final conclusion is not available upstream.\par

\textbf{Proof function.} The theorem fixes a role, prevents report laundering, classifies projected content, or excludes a same-domain independent discriminator. Analytic content is restricted to the explicit-formula scale projection declared by model adequacy.\par

\textbf{Anti-misreading.} This packet must not be read as banning analytic syntax, coordinate data, explicit-formula support, symmetry-pair information, numerical evidence, density estimates, or ordinary negative theoremhood. The excluded object is the endpoint-standing misuse named in the theorem.\par

\textbf{Hostile denial burden.} A successful objection must produce a concrete failure: carrier mismatch, hidden endpoint assumption, projection failure, lawful fifth role, annotation-discharge error, or a discriminator compatible with no-independent closure. Genre objections do not meet this burden.\par

\textbf{Worksheet entry.} The theorem appears in the theorem ladder, anti-circularity audit, route-locus disposition, and faithful counterexample worksheet. If this packet is removed, the manuscript loses a declared dependency and must not use the theorem downstream.\par

\textbf{RH-EF specialization.} In the prime-trace amplitude route, the fixed endpoint image is critical-neutral normalized amplitude. Descriptive non-neutral amplitude remains lawful. Endpoint-standing non-neutral amplitude is the only object classified as independent amplitude discriminator work.\par

\newpage

\subsection{Audit Packet 31: Incompleteness firewall}

\textbf{Purpose.} Classifies independence or true-but-unprovable statuses as inert or gate work.\par

\textbf{Statement.} The theorem named in this packet is used only for its declared role. It does not assert the final RH endpoint unless the final universalization theorem explicitly invokes it.\par

\textbf{Dependency.} The theorem depends only on prior target fixation, carrier adequacy, ATS/UEAP locking, explicit-formula projection, projection-role content, or the kernel-forced A+ consequence layer. The final conclusion is not available upstream.\par

\textbf{Proof function.} The theorem fixes a role, prevents report laundering, classifies projected content, or excludes a same-domain independent discriminator. Analytic content is restricted to the explicit-formula scale projection declared by model adequacy.\par

\textbf{Anti-misreading.} This packet must not be read as banning analytic syntax, coordinate data, explicit-formula support, symmetry-pair information, numerical evidence, density estimates, or ordinary negative theoremhood. The excluded object is the endpoint-standing misuse named in the theorem.\par

\textbf{Hostile denial burden.} A successful objection must produce a concrete failure: carrier mismatch, hidden endpoint assumption, projection failure, lawful fifth role, annotation-discharge error, or a discriminator compatible with no-independent closure. Genre objections do not meet this burden.\par

\textbf{Worksheet entry.} The theorem appears in the theorem ladder, anti-circularity audit, route-locus disposition, and faithful counterexample worksheet. If this packet is removed, the manuscript loses a declared dependency and must not use the theorem downstream.\par

\textbf{RH-EF specialization.} In the prime-trace amplitude route, the fixed endpoint image is critical-neutral normalized amplitude. Descriptive non-neutral amplitude remains lawful. Endpoint-standing non-neutral amplitude is the only object classified as independent amplitude discriminator work.\par

\newpage

\subsection{Audit Packet 32: Route-locus exhaustion}

\textbf{Purpose.} Forces each hostile alternative into a disposition table.\par

\textbf{Statement.} The theorem named in this packet is used only for its declared role. It does not assert the final RH endpoint unless the final universalization theorem explicitly invokes it.\par

\textbf{Dependency.} The theorem depends only on prior target fixation, carrier adequacy, ATS/UEAP locking, explicit-formula projection, projection-role content, or the kernel-forced A+ consequence layer. The final conclusion is not available upstream.\par

\textbf{Proof function.} The theorem fixes a role, prevents report laundering, classifies projected content, or excludes a same-domain independent discriminator. Analytic content is restricted to the explicit-formula scale projection declared by model adequacy.\par

\textbf{Anti-misreading.} This packet must not be read as banning analytic syntax, coordinate data, explicit-formula support, symmetry-pair information, numerical evidence, density estimates, or ordinary negative theoremhood. The excluded object is the endpoint-standing misuse named in the theorem.\par

\textbf{Hostile denial burden.} A successful objection must produce a concrete failure: carrier mismatch, hidden endpoint assumption, projection failure, lawful fifth role, annotation-discharge error, or a discriminator compatible with no-independent closure. Genre objections do not meet this burden.\par

\textbf{Worksheet entry.} The theorem appears in the theorem ladder, anti-circularity audit, route-locus disposition, and faithful counterexample worksheet. If this packet is removed, the manuscript loses a declared dependency and must not use the theorem downstream.\par

\textbf{RH-EF specialization.} In the prime-trace amplitude route, the fixed endpoint image is critical-neutral normalized amplitude. Descriptive non-neutral amplitude remains lawful. Endpoint-standing non-neutral amplitude is the only object classified as independent amplitude discriminator work.\par

\newpage

\subsection{Audit Packet 33: Weakening resistance}

\textbf{Purpose.} Prevents a weaker same-carrier packet from admitting separator governance.\par

\textbf{Statement.} The theorem named in this packet is used only for its declared role. It does not assert the final RH endpoint unless the final universalization theorem explicitly invokes it.\par

\textbf{Dependency.} The theorem depends only on prior target fixation, carrier adequacy, ATS/UEAP locking, explicit-formula projection, projection-role content, or the kernel-forced A+ consequence layer. The final conclusion is not available upstream.\par

\textbf{Proof function.} The theorem fixes a role, prevents report laundering, classifies projected content, or excludes a same-domain independent discriminator. Analytic content is restricted to the explicit-formula scale projection declared by model adequacy.\par

\textbf{Anti-misreading.} This packet must not be read as banning analytic syntax, coordinate data, explicit-formula support, symmetry-pair information, numerical evidence, density estimates, or ordinary negative theoremhood. The excluded object is the endpoint-standing misuse named in the theorem.\par

\textbf{Hostile denial burden.} A successful objection must produce a concrete failure: carrier mismatch, hidden endpoint assumption, projection failure, lawful fifth role, annotation-discharge error, or a discriminator compatible with no-independent closure. Genre objections do not meet this burden.\par

\textbf{Worksheet entry.} The theorem appears in the theorem ladder, anti-circularity audit, route-locus disposition, and faithful counterexample worksheet. If this packet is removed, the manuscript loses a declared dependency and must not use the theorem downstream.\par

\textbf{RH-EF specialization.} In the prime-trace amplitude route, the fixed endpoint image is critical-neutral normalized amplitude. Descriptive non-neutral amplitude remains lawful. Endpoint-standing non-neutral amplitude is the only object classified as independent amplitude discriminator work.\par

\newpage

\subsection{Audit Packet 34: Final universalization}

\textbf{Purpose.} Moves from arbitrary zerohood closure to the official universal RH endpoint.\par

\textbf{Statement.} The theorem named in this packet is used only for its declared role. It does not assert the final RH endpoint unless the final universalization theorem explicitly invokes it.\par

\textbf{Dependency.} The theorem depends only on prior target fixation, carrier adequacy, ATS/UEAP locking, explicit-formula projection, projection-role content, or the kernel-forced A+ consequence layer. The final conclusion is not available upstream.\par

\textbf{Proof function.} The theorem fixes a role, prevents report laundering, classifies projected content, or excludes a same-domain independent discriminator. Analytic content is restricted to the explicit-formula scale projection declared by model adequacy.\par

\textbf{Anti-misreading.} This packet must not be read as banning analytic syntax, coordinate data, explicit-formula support, symmetry-pair information, numerical evidence, density estimates, or ordinary negative theoremhood. The excluded object is the endpoint-standing misuse named in the theorem.\par

\textbf{Hostile denial burden.} A successful objection must produce a concrete failure: carrier mismatch, hidden endpoint assumption, projection failure, lawful fifth role, annotation-discharge error, or a discriminator compatible with no-independent closure. Genre objections do not meet this burden.\par

\textbf{Worksheet entry.} The theorem appears in the theorem ladder, anti-circularity audit, route-locus disposition, and faithful counterexample worksheet. If this packet is removed, the manuscript loses a declared dependency and must not use the theorem downstream.\par

\textbf{RH-EF specialization.} In the prime-trace amplitude route, the fixed endpoint image is critical-neutral normalized amplitude. Descriptive non-neutral amplitude remains lawful. Endpoint-standing non-neutral amplitude is the only object classified as independent amplitude discriminator work.\par

\newpage
\section{Expanded Hostile-Referee Matrix}

\begin{longtable}{L{0.32\textwidth}L{0.58\textwidth}}
\toprule
\textbf{Hostile objection} & \textbf{Manuscript response}\\
\midrule
This is not a zero-free-region proof. & Correct. The proof class is endpoint-exclusion via explicit-formula amplitude projection.\\
The explicit formula allows off-line terms. & Correct. They are lawful descriptive/support terms until used as endpoint-standing separator work.\\
The real part is native. & Correct. The real part is not excluded. Endpoint-standing prime-trace scale discrimination is the excluded role.\\
The amplitude exponent is native. & Native as analytic data, non-native as endpoint-standing authority over the critical-neutral image.\\
Critical normalization assumes RH. & No. It fixes the positive endpoint image; it does not assert universal occupation.\\
Pairing can cancel. & Pairing is allowed. Endpoint-neutral cancellation requires a bridge; it is not automatic.\\
A single term need not dominate. & The proof uses scale-class typing, not dominance.\\
Finite verification supports RH. & Yes, as finite support. It is not universal endpoint occupation.\\
Density estimates support RH. & Yes, as support. Scarcity is not absence unless bridged.\\
Hilbert--Polya could solve RH. & Yes, if it supplies a bridge to the same zeta/prime-trace carrier.\\
Positivity criteria can be equivalent to RH. & Yes, if the criterion is proved; otherwise it is bridge program or support.\\
The proof bans counterexamples. & It bans only endpoint-standing amplitude separator use; raw counterexample syntax is preserved until local contradiction.\\
The proof proves arbitrary propositions. & No. Anti-trivialization requires official route, projection-role content, and non-neutral scale standing.\\
The local reductio adds an assumption. & No. The annotation is a proof-act role, not object-level content.\\
Ordinary negative theoremhood is forbidden. & No. Official endpoint use of it is role-classified.\\
Open status is a third branch. & No. It is epistemic unless it changes endpoint standing; then it is gate work.\\
A+ does analytic work. & No. A+ excludes a discriminator after the explicit formula supplies the amplitude role.\\
The support matrix is not analytic evidence. & Correct. It supplies structural AASC locking support only.\\
The explicit-formula route is only a reformulation. & It is the declared endpoint route; reformulation becomes role-bearing only under official use.\\
The endpoint image is one-sided. & The broad outcome carrier is symmetric; the endpoint image carrier has critical-neutral occupation and non-occupation.\\
\bottomrule
\end{longtable}

\newpage
\section{Expanded Faithful Counterexample Tasks}
A critic must produce an object, not a complaint. The following task list is cumulative.
\begin{enumerate}
\item Exhibit a clause of the context object that asserts RH.
\item Exhibit a clause of the explicit-formula model that assumes all zeros have beta equal to one-half.
\item Exhibit a nontrivial zero contribution whose normalized exponent is not beta minus one-half.
\item Exhibit a same-zero symmetry-pair bridge that makes non-neutral exponent endpoint-neutral without proving the positive endpoint.
\item Exhibit a lawful endpoint-standing non-neutral amplitude role that is not support, bookkeeping, repair, carrier shift, surrogate substitution, coequal occupation, or independent discriminator work.
\item Exhibit an independent amplitude discriminator compatible with \(\Namp\).
\item Exhibit a proof-role annotation that functions as object-level premise in the local reductio.
\item Exhibit a failure of universal generalization from arbitrary \(\rho\).
\item Exhibit an incompleteness or independence distinction that changes endpoint standing without becoming endpoint-admissibility-relevant gate work.
\item Exhibit a strict weakening preserving the same endpoint carrier while allowing non-neutral amplitude endpoint governance without independent classifier work.
\end{enumerate}

\newpage
\section{Build-Integrity Repetition Ledger}
The defensive principle of this manuscript is that no theorem may be used only once if it is load-bearing. Every decisive bridge is restated in four registers: formal statement, proof, anti-misreading audit, and denial-burden entry. The result is intentionally redundant. The redundancy is not padding; it prevents a hostile reader from sliding between coordinate description, support content, endpoint counterforce, and endpoint outcome.

\begin{longtable}{L{0.28\textwidth}L{0.22\textwidth}L{0.36\textwidth}}
\toprule
\textbf{Load-bearing link} & \textbf{Primary section} & \textbf{Additional audit locations}\\
\midrule
Context non-smuggling & Section 1 & Theorem ladder, anti-circularity audit, counterexample worksheet.\\
Official route fixation & Section 1 & Proof spine, projection-role content, denial burden.\\
Kernel forcing & Section 2 & A+ theorem, incompleteness firewall, route-locus exhaustion.\\
ATS locking & Section 2 & Corpus ledger, skin authority audit, formalization map.\\
UEAP report ladder & Section 2 & Non-explosiveness, no-overreporting audit, hostile matrix.\\
Model adequacy & Section 3 & Term isolation, pairing closure, route-locus table.\\
Branch projection & Section 4 & Single-interface theorem, endpoint-role content, theorem ladder.\\
Non-neutral scale standing & Section 5 & Hinge theorem, hostile matrix, faithful counterexample task.\\
Discriminator induction & Section 6 & A+ closure, local contradiction, denial burden.\\
Annotation discharge & Section 8 & Anti-circularity audit, proof spine, local reductio theorem.\\
Universalization & Section 8 & Final endpoint corollary, reference integrity, theorem ladder.\\
\bottomrule
\end{longtable}

\section{Adversarial Same-Mode Denial Catalogue}

The following catalogue repeats the right-mode denial burden in adversarial form. Each entry is designed to prevent the proof from relying on a compressed or implicit defense.\par

\newpage

\subsection{Denial Packet 1: Projection failure}

\textbf{Required objection.} The critic must show that the explicit-formula route does not carry off-line zerohood into non-neutral critical-normalized amplitude.\par

\textbf{Why a generic objection fails.} A demand for an analytic estimate, a different proof class, or an expression of unfamiliarity with AASC vocabulary does not engage this denial packet. The packet requires a same-mode failure in target, carrier, role, projection, or endpoint-use governance.\par

\textbf{Manuscript lock.} The relevant lock is present in the formal theorem statement, repeated in the audit capsule, and indexed in the theorem ladder. If the objection attacks a theorem, it must identify the exact premise whose proof fails, not merely assert that the conclusion is surprising.\par

\textbf{Allowed mathematical content.} The manuscript continues to allow raw analytic counterexample syntax, descriptive coordinates, explicit-formula support, finite evidence, density estimates, surrogate bridge programs, and ordinary theoremhood. The packet concerns only endpoint-standing use on the fixed RH-EF carrier.\par

\textbf{Closure route.} If the objection changes endpoint standing, reportability, branch exclusion, or admissible continuation, it is endpoint-admissibility-relevant. It must then survive AMetricity, ATS locking, UEAP report preservation, and no-independent-discriminator closure.\par

\newpage

\subsection{Denial Packet 2: Amplitude adequacy failure}

\textbf{Required objection.} The critic must show that the normalized exponent beta minus one-half is not the correct endpoint-relevant scale class.\par

\textbf{Why a generic objection fails.} A demand for an analytic estimate, a different proof class, or an expression of unfamiliarity with AASC vocabulary does not engage this denial packet. The packet requires a same-mode failure in target, carrier, role, projection, or endpoint-use governance.\par

\textbf{Manuscript lock.} The relevant lock is present in the formal theorem statement, repeated in the audit capsule, and indexed in the theorem ladder. If the objection attacks a theorem, it must identify the exact premise whose proof fails, not merely assert that the conclusion is surprising.\par

\textbf{Allowed mathematical content.} The manuscript continues to allow raw analytic counterexample syntax, descriptive coordinates, explicit-formula support, finite evidence, density estimates, surrogate bridge programs, and ordinary theoremhood. The packet concerns only endpoint-standing use on the fixed RH-EF carrier.\par

\textbf{Closure route.} If the objection changes endpoint standing, reportability, branch exclusion, or admissible continuation, it is endpoint-admissibility-relevant. It must then survive AMetricity, ATS locking, UEAP report preservation, and no-independent-discriminator closure.\par

\newpage

\subsection{Denial Packet 3: Term isolation failure}

\textbf{Required objection.} The critic must show that an individual zero contribution cannot be typed without asserting global dominance.\par

\textbf{Why a generic objection fails.} A demand for an analytic estimate, a different proof class, or an expression of unfamiliarity with AASC vocabulary does not engage this denial packet. The packet requires a same-mode failure in target, carrier, role, projection, or endpoint-use governance.\par

\textbf{Manuscript lock.} The relevant lock is present in the formal theorem statement, repeated in the audit capsule, and indexed in the theorem ladder. If the objection attacks a theorem, it must identify the exact premise whose proof fails, not merely assert that the conclusion is surprising.\par

\textbf{Allowed mathematical content.} The manuscript continues to allow raw analytic counterexample syntax, descriptive coordinates, explicit-formula support, finite evidence, density estimates, surrogate bridge programs, and ordinary theoremhood. The packet concerns only endpoint-standing use on the fixed RH-EF carrier.\par

\textbf{Closure route.} If the objection changes endpoint standing, reportability, branch exclusion, or admissible continuation, it is endpoint-admissibility-relevant. It must then survive AMetricity, ATS locking, UEAP report preservation, and no-independent-discriminator closure.\par

\newpage

\subsection{Denial Packet 4: Pairing neutralization failure}

\textbf{Required objection.} The critic must show that pairing alone produces endpoint-neutral occupation without proving the critical-line endpoint.\par

\textbf{Why a generic objection fails.} A demand for an analytic estimate, a different proof class, or an expression of unfamiliarity with AASC vocabulary does not engage this denial packet. The packet requires a same-mode failure in target, carrier, role, projection, or endpoint-use governance.\par

\textbf{Manuscript lock.} The relevant lock is present in the formal theorem statement, repeated in the audit capsule, and indexed in the theorem ladder. If the objection attacks a theorem, it must identify the exact premise whose proof fails, not merely assert that the conclusion is surprising.\par

\textbf{Allowed mathematical content.} The manuscript continues to allow raw analytic counterexample syntax, descriptive coordinates, explicit-formula support, finite evidence, density estimates, surrogate bridge programs, and ordinary theoremhood. The packet concerns only endpoint-standing use on the fixed RH-EF carrier.\par

\textbf{Closure route.} If the objection changes endpoint standing, reportability, branch exclusion, or admissible continuation, it is endpoint-admissibility-relevant. It must then survive AMetricity, ATS locking, UEAP report preservation, and no-independent-discriminator closure.\par

\newpage

\subsection{Denial Packet 5: Support endpoint collapse}

\textbf{Required objection.} The critic must show that support-level amplitude content can be endpoint-standing without a bridge.\par

\textbf{Why a generic objection fails.} A demand for an analytic estimate, a different proof class, or an expression of unfamiliarity with AASC vocabulary does not engage this denial packet. The packet requires a same-mode failure in target, carrier, role, projection, or endpoint-use governance.\par

\textbf{Manuscript lock.} The relevant lock is present in the formal theorem statement, repeated in the audit capsule, and indexed in the theorem ladder. If the objection attacks a theorem, it must identify the exact premise whose proof fails, not merely assert that the conclusion is surprising.\par

\textbf{Allowed mathematical content.} The manuscript continues to allow raw analytic counterexample syntax, descriptive coordinates, explicit-formula support, finite evidence, density estimates, surrogate bridge programs, and ordinary theoremhood. The packet concerns only endpoint-standing use on the fixed RH-EF carrier.\par

\textbf{Closure route.} If the objection changes endpoint standing, reportability, branch exclusion, or admissible continuation, it is endpoint-admissibility-relevant. It must then survive AMetricity, ATS locking, UEAP report preservation, and no-independent-discriminator closure.\par

\newpage

\subsection{Denial Packet 6: ATS skin authority}

\textbf{Required objection.} The critic must show that a presentation skin can carry endpoint force without changing anchor or tensor content.\par

\textbf{Why a generic objection fails.} A demand for an analytic estimate, a different proof class, or an expression of unfamiliarity with AASC vocabulary does not engage this denial packet. The packet requires a same-mode failure in target, carrier, role, projection, or endpoint-use governance.\par

\textbf{Manuscript lock.} The relevant lock is present in the formal theorem statement, repeated in the audit capsule, and indexed in the theorem ladder. If the objection attacks a theorem, it must identify the exact premise whose proof fails, not merely assert that the conclusion is surprising.\par

\textbf{Allowed mathematical content.} The manuscript continues to allow raw analytic counterexample syntax, descriptive coordinates, explicit-formula support, finite evidence, density estimates, surrogate bridge programs, and ordinary theoremhood. The packet concerns only endpoint-standing use on the fixed RH-EF carrier.\par

\textbf{Closure route.} If the objection changes endpoint standing, reportability, branch exclusion, or admissible continuation, it is endpoint-admissibility-relevant. It must then survive AMetricity, ATS locking, UEAP report preservation, and no-independent-discriminator closure.\par

\newpage

\subsection{Denial Packet 7: UEAP overreporting exception}

\textbf{Required objection.} The critic must show that the RH-EF endpoint permits reporting lower support status as endpoint result.\par

\textbf{Why a generic objection fails.} A demand for an analytic estimate, a different proof class, or an expression of unfamiliarity with AASC vocabulary does not engage this denial packet. The packet requires a same-mode failure in target, carrier, role, projection, or endpoint-use governance.\par

\textbf{Manuscript lock.} The relevant lock is present in the formal theorem statement, repeated in the audit capsule, and indexed in the theorem ladder. If the objection attacks a theorem, it must identify the exact premise whose proof fails, not merely assert that the conclusion is surprising.\par

\textbf{Allowed mathematical content.} The manuscript continues to allow raw analytic counterexample syntax, descriptive coordinates, explicit-formula support, finite evidence, density estimates, surrogate bridge programs, and ordinary theoremhood. The packet concerns only endpoint-standing use on the fixed RH-EF carrier.\par

\textbf{Closure route.} If the objection changes endpoint standing, reportability, branch exclusion, or admissible continuation, it is endpoint-admissibility-relevant. It must then survive AMetricity, ATS locking, UEAP report preservation, and no-independent-discriminator closure.\par

\newpage

\subsection{Denial Packet 8: Independent discriminator denial}

\textbf{Required objection.} The critic must show that noncritical amplitude status work has a remaining lawful role after trichotomy exhaustion.\par

\textbf{Why a generic objection fails.} A demand for an analytic estimate, a different proof class, or an expression of unfamiliarity with AASC vocabulary does not engage this denial packet. The packet requires a same-mode failure in target, carrier, role, projection, or endpoint-use governance.\par

\textbf{Manuscript lock.} The relevant lock is present in the formal theorem statement, repeated in the audit capsule, and indexed in the theorem ladder. If the objection attacks a theorem, it must identify the exact premise whose proof fails, not merely assert that the conclusion is surprising.\par

\textbf{Allowed mathematical content.} The manuscript continues to allow raw analytic counterexample syntax, descriptive coordinates, explicit-formula support, finite evidence, density estimates, surrogate bridge programs, and ordinary theoremhood. The packet concerns only endpoint-standing use on the fixed RH-EF carrier.\par

\textbf{Closure route.} If the objection changes endpoint standing, reportability, branch exclusion, or admissible continuation, it is endpoint-admissibility-relevant. It must then survive AMetricity, ATS locking, UEAP report preservation, and no-independent-discriminator closure.\par

\newpage

\subsection{Denial Packet 9: A+ consequence denial}

\textbf{Required objection.} The critic must show that endpoint use on a fixed carrier does not trigger no-independent-discriminator closure.\par

\textbf{Why a generic objection fails.} A demand for an analytic estimate, a different proof class, or an expression of unfamiliarity with AASC vocabulary does not engage this denial packet. The packet requires a same-mode failure in target, carrier, role, projection, or endpoint-use governance.\par

\textbf{Manuscript lock.} The relevant lock is present in the formal theorem statement, repeated in the audit capsule, and indexed in the theorem ladder. If the objection attacks a theorem, it must identify the exact premise whose proof fails, not merely assert that the conclusion is surprising.\par

\textbf{Allowed mathematical content.} The manuscript continues to allow raw analytic counterexample syntax, descriptive coordinates, explicit-formula support, finite evidence, density estimates, surrogate bridge programs, and ordinary theoremhood. The packet concerns only endpoint-standing use on the fixed RH-EF carrier.\par

\textbf{Closure route.} If the objection changes endpoint standing, reportability, branch exclusion, or admissible continuation, it is endpoint-admissibility-relevant. It must then survive AMetricity, ATS locking, UEAP report preservation, and no-independent-discriminator closure.\par

\newpage

\subsection{Denial Packet 10: Raw negation promotion}

\textbf{Required objection.} The critic must show that the manuscript promotes bare negation before projection-role content and counterforce.\par

\textbf{Why a generic objection fails.} A demand for an analytic estimate, a different proof class, or an expression of unfamiliarity with AASC vocabulary does not engage this denial packet. The packet requires a same-mode failure in target, carrier, role, projection, or endpoint-use governance.\par

\textbf{Manuscript lock.} The relevant lock is present in the formal theorem statement, repeated in the audit capsule, and indexed in the theorem ladder. If the objection attacks a theorem, it must identify the exact premise whose proof fails, not merely assert that the conclusion is surprising.\par

\textbf{Allowed mathematical content.} The manuscript continues to allow raw analytic counterexample syntax, descriptive coordinates, explicit-formula support, finite evidence, density estimates, surrogate bridge programs, and ordinary theoremhood. The packet concerns only endpoint-standing use on the fixed RH-EF carrier.\par

\textbf{Closure route.} If the objection changes endpoint standing, reportability, branch exclusion, or admissible continuation, it is endpoint-admissibility-relevant. It must then survive AMetricity, ATS locking, UEAP report preservation, and no-independent-discriminator closure.\par

\newpage

\subsection{Denial Packet 11: Annotation premise objection}

\textbf{Required objection.} The critic must show that the reductio annotation functions as an object-level premise.\par

\textbf{Why a generic objection fails.} A demand for an analytic estimate, a different proof class, or an expression of unfamiliarity with AASC vocabulary does not engage this denial packet. The packet requires a same-mode failure in target, carrier, role, projection, or endpoint-use governance.\par

\textbf{Manuscript lock.} The relevant lock is present in the formal theorem statement, repeated in the audit capsule, and indexed in the theorem ladder. If the objection attacks a theorem, it must identify the exact premise whose proof fails, not merely assert that the conclusion is surprising.\par

\textbf{Allowed mathematical content.} The manuscript continues to allow raw analytic counterexample syntax, descriptive coordinates, explicit-formula support, finite evidence, density estimates, surrogate bridge programs, and ordinary theoremhood. The packet concerns only endpoint-standing use on the fixed RH-EF carrier.\par

\textbf{Closure route.} If the objection changes endpoint standing, reportability, branch exclusion, or admissible continuation, it is endpoint-admissibility-relevant. It must then survive AMetricity, ATS locking, UEAP report preservation, and no-independent-discriminator closure.\par

\newpage

\subsection{Denial Packet 12: Ordinary theoremhood shield}

\textbf{Required objection.} The critic must show that official endpoint use can retain endpoint force while escaping endpoint-role discipline.\par

\textbf{Why a generic objection fails.} A demand for an analytic estimate, a different proof class, or an expression of unfamiliarity with AASC vocabulary does not engage this denial packet. The packet requires a same-mode failure in target, carrier, role, projection, or endpoint-use governance.\par

\textbf{Manuscript lock.} The relevant lock is present in the formal theorem statement, repeated in the audit capsule, and indexed in the theorem ladder. If the objection attacks a theorem, it must identify the exact premise whose proof fails, not merely assert that the conclusion is surprising.\par

\textbf{Allowed mathematical content.} The manuscript continues to allow raw analytic counterexample syntax, descriptive coordinates, explicit-formula support, finite evidence, density estimates, surrogate bridge programs, and ordinary theoremhood. The packet concerns only endpoint-standing use on the fixed RH-EF carrier.\par

\textbf{Closure route.} If the objection changes endpoint standing, reportability, branch exclusion, or admissible continuation, it is endpoint-admissibility-relevant. It must then survive AMetricity, ATS locking, UEAP report preservation, and no-independent-discriminator closure.\par

\newpage

\subsection{Denial Packet 13: Incompleteness rescue}

\textbf{Required objection.} The critic must show that independence or true-but-unprovable status changes endpoint standing without becoming gate work.\par

\textbf{Why a generic objection fails.} A demand for an analytic estimate, a different proof class, or an expression of unfamiliarity with AASC vocabulary does not engage this denial packet. The packet requires a same-mode failure in target, carrier, role, projection, or endpoint-use governance.\par

\textbf{Manuscript lock.} The relevant lock is present in the formal theorem statement, repeated in the audit capsule, and indexed in the theorem ladder. If the objection attacks a theorem, it must identify the exact premise whose proof fails, not merely assert that the conclusion is surprising.\par

\textbf{Allowed mathematical content.} The manuscript continues to allow raw analytic counterexample syntax, descriptive coordinates, explicit-formula support, finite evidence, density estimates, surrogate bridge programs, and ordinary theoremhood. The packet concerns only endpoint-standing use on the fixed RH-EF carrier.\par

\textbf{Closure route.} If the objection changes endpoint standing, reportability, branch exclusion, or admissible continuation, it is endpoint-admissibility-relevant. It must then survive AMetricity, ATS locking, UEAP report preservation, and no-independent-discriminator closure.\par

\newpage

\subsection{Denial Packet 14: Carrier shift rescue}

\textbf{Required objection.} The critic must show that a surrogate carrier remains the same RH-EF endpoint carrier without a bridge.\par

\textbf{Why a generic objection fails.} A demand for an analytic estimate, a different proof class, or an expression of unfamiliarity with AASC vocabulary does not engage this denial packet. The packet requires a same-mode failure in target, carrier, role, projection, or endpoint-use governance.\par

\textbf{Manuscript lock.} The relevant lock is present in the formal theorem statement, repeated in the audit capsule, and indexed in the theorem ladder. If the objection attacks a theorem, it must identify the exact premise whose proof fails, not merely assert that the conclusion is surprising.\par

\textbf{Allowed mathematical content.} The manuscript continues to allow raw analytic counterexample syntax, descriptive coordinates, explicit-formula support, finite evidence, density estimates, surrogate bridge programs, and ordinary theoremhood. The packet concerns only endpoint-standing use on the fixed RH-EF carrier.\par

\textbf{Closure route.} If the objection changes endpoint standing, reportability, branch exclusion, or admissible continuation, it is endpoint-admissibility-relevant. It must then survive AMetricity, ATS locking, UEAP report preservation, and no-independent-discriminator closure.\par

\newpage

\subsection{Denial Packet 15: Critical normalization smuggling}

\textbf{Required objection.} The critic must show that critical normalization asserts universal neutral occupation rather than defining the positive endpoint image.\par

\textbf{Why a generic objection fails.} A demand for an analytic estimate, a different proof class, or an expression of unfamiliarity with AASC vocabulary does not engage this denial packet. The packet requires a same-mode failure in target, carrier, role, projection, or endpoint-use governance.\par

\textbf{Manuscript lock.} The relevant lock is present in the formal theorem statement, repeated in the audit capsule, and indexed in the theorem ladder. If the objection attacks a theorem, it must identify the exact premise whose proof fails, not merely assert that the conclusion is surprising.\par

\textbf{Allowed mathematical content.} The manuscript continues to allow raw analytic counterexample syntax, descriptive coordinates, explicit-formula support, finite evidence, density estimates, surrogate bridge programs, and ordinary theoremhood. The packet concerns only endpoint-standing use on the fixed RH-EF carrier.\par

\textbf{Closure route.} If the objection changes endpoint standing, reportability, branch exclusion, or admissible continuation, it is endpoint-admissibility-relevant. It must then survive AMetricity, ATS locking, UEAP report preservation, and no-independent-discriminator closure.\par

\newpage

\subsection{Denial Packet 16: Final universalization failure}

\textbf{Required objection.} The critic must show why the arbitrary-zerohood theorem cannot be universally generalized.\par

\textbf{Why a generic objection fails.} A demand for an analytic estimate, a different proof class, or an expression of unfamiliarity with AASC vocabulary does not engage this denial packet. The packet requires a same-mode failure in target, carrier, role, projection, or endpoint-use governance.\par

\textbf{Manuscript lock.} The relevant lock is present in the formal theorem statement, repeated in the audit capsule, and indexed in the theorem ladder. If the objection attacks a theorem, it must identify the exact premise whose proof fails, not merely assert that the conclusion is surprising.\par

\textbf{Allowed mathematical content.} The manuscript continues to allow raw analytic counterexample syntax, descriptive coordinates, explicit-formula support, finite evidence, density estimates, surrogate bridge programs, and ordinary theoremhood. The packet concerns only endpoint-standing use on the fixed RH-EF carrier.\par

\textbf{Closure route.} If the objection changes endpoint standing, reportability, branch exclusion, or admissible continuation, it is endpoint-admissibility-relevant. It must then survive AMetricity, ATS locking, UEAP report preservation, and no-independent-discriminator closure.\par

\newpage

\subsection{Denial Packet 17: Hidden fifth role}

\textbf{Required objection.} The critic must exhibit endpoint-resolving non-neutral amplitude that is neither support, bookkeeping, carrier shift, repair, coequal occupation, nor independent discriminator.\par

\textbf{Why a generic objection fails.} A demand for an analytic estimate, a different proof class, or an expression of unfamiliarity with AASC vocabulary does not engage this denial packet. The packet requires a same-mode failure in target, carrier, role, projection, or endpoint-use governance.\par

\textbf{Manuscript lock.} The relevant lock is present in the formal theorem statement, repeated in the audit capsule, and indexed in the theorem ladder. If the objection attacks a theorem, it must identify the exact premise whose proof fails, not merely assert that the conclusion is surprising.\par

\textbf{Allowed mathematical content.} The manuscript continues to allow raw analytic counterexample syntax, descriptive coordinates, explicit-formula support, finite evidence, density estimates, surrogate bridge programs, and ordinary theoremhood. The packet concerns only endpoint-standing use on the fixed RH-EF carrier.\par

\textbf{Closure route.} If the objection changes endpoint standing, reportability, branch exclusion, or admissible continuation, it is endpoint-admissibility-relevant. It must then survive AMetricity, ATS locking, UEAP report preservation, and no-independent-discriminator closure.\par

\newpage

\subsection{Denial Packet 18: Same-domain weakening}

\textbf{Required objection.} The critic must define a weaker same-carrier regime that preserves target adequacy while admitting amplitude separator governance.\par

\textbf{Why a generic objection fails.} A demand for an analytic estimate, a different proof class, or an expression of unfamiliarity with AASC vocabulary does not engage this denial packet. The packet requires a same-mode failure in target, carrier, role, projection, or endpoint-use governance.\par

\textbf{Manuscript lock.} The relevant lock is present in the formal theorem statement, repeated in the audit capsule, and indexed in the theorem ladder. If the objection attacks a theorem, it must identify the exact premise whose proof fails, not merely assert that the conclusion is surprising.\par

\textbf{Allowed mathematical content.} The manuscript continues to allow raw analytic counterexample syntax, descriptive coordinates, explicit-formula support, finite evidence, density estimates, surrogate bridge programs, and ordinary theoremhood. The packet concerns only endpoint-standing use on the fixed RH-EF carrier.\par

\textbf{Closure route.} If the objection changes endpoint standing, reportability, branch exclusion, or admissible continuation, it is endpoint-admissibility-relevant. It must then survive AMetricity, ATS locking, UEAP report preservation, and no-independent-discriminator closure.\par

\newpage
\appendix
\section{Theorem Ladder}\label{app:ladder}

\begin{enumerate}
\item Bare counterexample content is separated from endpoint use.
\item Ordinary counterexample theoremhood is granted as kernel-internal theoremhood, not rejected for negativity.
\item Official counterexample force is proved to be endpoint use.
\item No autonomous counterexample endpoint role is proved: official endpoint-defeating force cannot stand outside endpoint use.
\item Autonomous ordinary-negative endpoint roles collapse to raw content, support, bookkeeping, ordinary theoremhood without endpoint force, local syntax without endpoint force, or carrier/domain shift.
\item Official counterexample force is shown to satisfy target adequacy and therefore to trigger the kernel and A+ consequence layer.
\item Local exact-complement counterexample force is proved to be local endpoint use without becoming global negative endpoint outcome.
\item RH off-line counterexample force projects to explicit-formula non-neutral order and enters native-counterforce exhaustion.
\end{enumerate}


\begin{enumerate}
\item The official RH endpoint is fixed as \(\forall\rho(\Az(\rho)\Rightarrow\Lcrit(\rho))\).
\item The RH-EF context fixes zeta carrier, prime-trace carrier, explicit-formula route, branch slots, and amplitude endpoint image without asserting RH.
\item Official endpoint resolution through the explicit-formula route satisfies target adequacy.
\item Target adequacy forces the kernel; endpoint use on a fixed domain forces A+.
\item ATS locks anchor, tensor, and skin roles; UEAP locks report levels.
\item The context-bound explicit-formula model is adequate for zero-to-amplitude projection.
\item A zero \(\rho=\beta+i\gamma\) projects to critical-normalized amplitude exponent \(\beta-1/2\).
\item The structural-to-EF bridge identifies \(\Lcrit(\rho)\) with \(\ordhalf(\rho)=0\) and \(\neg\Lcrit(\rho)\) with \(\ordhalf(\rho)\ne0\).
\item \(\BzetaEF\) instantiates the structural carrier-preservation bridge without containing line support by definition.
\item The bridge fork gives: \(\BridgeCompleteEF\), \(\DomainShiftEF\), or \(\Damp\).
\item Critical-line readout projects to neutral amplitude occupation.
\item Off-line readout projects to non-neutral amplitude non-occupation.
\item Projection-role content binds local exact-complement countercase use to non-neutral amplitude content under the official route.
\item Non-neutral amplitude alone is lawful descriptive/support content.
\item Endpoint-standing non-neutral amplitude is same-domain noncritical amplitude status work.
\item Noncritical amplitude status work induces \(\Damp(\rho)\).
\item Endpoint use of local amplitude separator induces \(\Namp(\rho)\), hence \(\neg\Damp(\rho)\).
\item Local amplitude separator occupation is impossible.
\item Exact-complement local countercase induces endpoint counterforce and local separator occupation.
\item The local countercase contradiction discharges \([\neg\Lcrit(\rho)]_i\).
\item Arbitrary zerohood yields \(\Lcrit(\rho)\).
\item Universal generalization gives RH.
\end{enumerate}

\section{Anti-Circularity Audit}\label{app:anticirc}

\begin{longtable}{L{0.32\textwidth}L{0.58\textwidth}}
\toprule
Potential misreading & Audit response\\
\midrule
The paper treats every counterexample as forbidden. & No.  Bare counterexample content is not endpoint use.  Only official counterexample force or local exact-complement endpoint counterforce enters endpoint-use discipline.\\
Objects are not acts. & Correct.  The object \(\rho\) is raw analytic content; offering it as a same-carrier endpoint-defeating counterexample is an act.\\
There is an autonomous ordinary-negative endpoint role. & No.  Corollary~\ref{cor:no-autonomous-counterexample-endpoint-role} states that official endpoint-defeating force entails endpoint use; without endpoint use the branch is raw content, support, bookkeeping, ordinary theoremhood without endpoint force, local syntax without endpoint force, or carrier/domain shift.\\
The paper bans analytic falsification. & No.  It classifies endpoint-standing falsification as endpoint use and then applies the same endpoint discipline that governs official resolution.\\
The proof proves arbitrary propositions. & No.  Arbitrary negation lacks official route projection, endpoint counterforce, and endpoint-status work.\\
\bottomrule
\end{longtable}


\begin{longtable}{L{0.31\textwidth}L{0.61\textwidth}}
\toprule
\textbf{Potential circle or misreading} & \textbf{Audit response}\\
\midrule
Context smuggles RH & \(\CRHEF\) fixes carrier and route only.  It does not assert line support or neutral occupation.\\
Structural EF bridge smuggles RH & \(\BzetaEF\) fixes same-zero, route, report-level, and trajectory conditions. It does not include \(\Lcrit\), \(\ordhalf=0\), or RH as a premise.\\
Bridge fork bans counterexamples by fiat & The fork allows bridge-complete neutralization and domain shift. Only bridge-failing same-carrier endpoint counterforce yields \(\Damp\).\\
Order class is just relabeled coordinate & \(\ordhalf\) is a recoverable asymptotic order invariant of the critical-normalized EF contribution. It is lawful as data and becomes AASC-relevant only under endpoint standing.\\
Explicit formula smuggles RH & The explicit formula gives zero-to-prime-trace projection; it does not assert \(\beta=1/2\).\\
Neutral amplitude assumes positivity & Neutral amplitude is the positive endpoint image, not a universal occupation claim.\\
Off-line coordinates are illegal & Coordinate reports are lawful.  Only endpoint-standing separator use is excluded.\\
Off-line amplitude is illegal & Non-neutral amplitude is lawful descriptive/support content.  Endpoint-standing use is the excluded role.\\
Pairing ignored & Pairing is lawful analytic structure.  It does not erase scale-class typing without endpoint-neutral bridge.\\
Arbitrary negation becomes separator & Blocked by the anti-trivialization theorem.\\
Reductio annotation is an added premise & The annotation is proof-act role only; it is not an object-level conjunct.\\
A+ does analytic estimation & A+ excludes independent endpoint-status work after explicit-formula projection identifies the role.\\
Ordinary negative theoremhood banned & Negative theoremhood is allowed.  Official endpoint use of it is role-classified.\\
Open or independent status is third branch & Open/independent status is epistemic unless it changes endpoint standing; then it is gate work.\\
\bottomrule
\end{longtable}

\section{Notation Ledger}\label{app:notation}

\begin{longtable}{L{0.26\textwidth}L{0.64\textwidth}}
\toprule
\textbf{Notation} & \textbf{Meaning}\\
\midrule
\(\Az(\rho)\) & Nontrivial zeta zerohood: \(\zeta(\rho)=0\) and \(0<\Ree(\rho)<1\).\\
\(\Lcrit(\rho)\) & Critical-line readout: \(\Ree(\rho)=1/2\).\\
\(\rho=\beta+i\gamma\) & Zero decomposed into real and imaginary parts.\\
\(\psi(x)\) & Chebyshev prime trace \(\sum_{n\le x}\Lambda(n)\).\\
\(\psierr(x)\) & Prime-trace error \(\psi(x)-x\).\\
\(\psinorm(x)\) & Critical-normalized trace \((\psi(x)-x)/x^{1/2}\).\\
\(\ordhalf(\rho)\) & Critical-normalized amplitude order \(\beta-1/2\).\\
\(\Dord(\rho)\) & Nonzero order-class discriminator \(\ordhalf(\rho)\ne0\).\\
\(\BzetaEF(\rho)\) & Explicit-formula instantiation of the structural bridge \(B_\zeta\), without built-in line support.\\
\(\BridgeCompleteEF(\rho)\) & Same-zero EF bridge completion yielding endpoint-neutral order.\\
\(\DomainShiftEF(\rho)\) & Shift of carrier, route, zero identity, normalization, or endpoint image.\\
\(\CRHEF\) & Fixed RH explicit-formula endpoint context.\\
\(\ModelEF\) & Context-bound explicit-formula endpoint model.\\
\(\OfficialEFRoute\) & Official evaluation of RH through the explicit-formula zero-to-prime-trace route.\\
\(\CritAmpOcc\) & Critical-neutral amplitude occupation.\\
\(\OffAmpExcl\) & Non-neutral amplitude non-occupation.\\
\(\AmpStand\) & Endpoint-standing non-neutral amplitude.\\
\(\AmpSepLoc\) & Local amplitude separator occupation inside the reductio subproof.\\
\(\Damp\) & Independent amplitude endpoint-status discriminator.\\
\(\Namp\) & No-independent amplitude discriminator closure.\\
\(\ATS\) & Anchor--tensor--skin structural locking discipline.\\
\(\UEAP\) & Report-preservation and laundering-exclusion discipline.\\
\bottomrule
\end{longtable}


\section{Corpus Support Ledger}\label{app:corpus}

This appendix records the selected support rows extracted from the uploaded AASC Corpus Control Matrix.  The local use of these rows is structural: kernel forcing, UEAP report discipline, ATS decomposition, AMetric no-selector discipline, no repair, no carrier transfer, silent-redescription control, role-occupancy closure, boundary-trace fixation, non-factorizability, and same-domain incompleteness control.  None of these rows supplies an analytic zero-location estimate or a new explicit formula.

\begin{longtable}{L{0.16\textwidth} L{0.20\textwidth} L{0.12\textwidth} L{0.21\textwidth} L{0.21\textwidth}}
\toprule
Local use & Source id & Type & Object & Manuscript use \\
\midrule
\endhead
AMetric boundary & ERK20260527-New\_erk\_files\_May\_41605D-THM-023 & Theorem & Admissibility constraint on Operator undefinedness & prevents boundary operators from carrying unauthorized endpoint status \\
AMetric no selector & ERK20260527-New\_erk\_files\_May\_41605D-THM-002 & Theorem & Admissibility constraint on no-selector/no-parameter boundary authority & no endpoint-bearing selector or metric at boundary \\
ATS identity lock & ERK20260527-New\_erk\_files\_May\_6DC183-LEM-003 & Lemma & (Congruence; identity persistence is forced by stable & same-zero/same-route identity under redescription \\
ATS redescription lock & ERK20260527-New\_erk\_files\_May\_6DC183-THM-028 & Theorem & i & blocks treating skin changes as tensor content \\
Bivalence & ERK20260527-New\_erk\_files\_May\_ACEB9B-PRO-012 & Proposition & Admissibility and bivalence in the typed domain & fail-closed endpoint status on fixed domain \\
Boundary no-overpromotion & AASC-EX6-THM-10-1 & Theorem & Master endpoint non-promotion theorem & endpoint non-promotion \\
Boundary trace & ERK20260527-New\_erk\_files\_May\_045695-EXM-001 & Example & Two-token environmental witness & declared trace requirement for continuation loci \\
Calibration/post-selection & AASC-EX5-COR-11-2 & Corollary & Post-selection exclusion & post-selection and after-the-fact exclusion \\
Endpoint boundary exactness & AASC-EX6-DEF-9-1 & Definition & Endpoint-fixed & endpoint/boundary exactness \\
Governed construction & ERK20260527-New\_erk\_files\_May\_35B594-DEF-004 & Definition & Governed construction and non-degenerate construction & endpoint use as governed construction \\
Kernel raw sequence & ERK20260527-New\_erk\_files\_May\_35B594-DEF-002 & Definition & Kernel-neutral target adequacy & neutral adequacy precondition for non-degenerate endpoint acts \\
Kernel target adequacy & ERK20260527-New\_erk\_files\_May\_35B594-LEM-001 & Lemma & Target adequacy forces the kernel roles & kernel forcing for endpoint use \\
No carriers & ERK20260527-New\_erk\_files\_May\_35B594-COR-012 & Corollary & No independent carriers beneath standing & blocks independent carrier of endpoint standing \\
No generators & ERK20260527-New\_erk\_files\_May\_EA512E-THM-010 & Theorem & no generators & blocks standing creation from failed or lower-status acts \\
No hidden transport & NOAR-LEM-001 & Lemma & No hidden flavor transport & hidden transport analogy for no hidden amplitude route \\
No repair & ERK20260527-New\_erk\_files\_May\_311F0C-THM-009 & Theorem & No post-selection repair & blocks later proof repairing earlier unbridged endpoint act \\
No same-domain totalization & ERK20260527-New\_erk\_files\_May\_3D011A-THM-056 & Theorem & No internal totalization & same-carrier totalization block \\
Quotient identity & ERK20260527-New\_erk\_files\_May\_4C9E73-THM-021 & Theorem & (Quotient-only presentation is insufficient for fixed-scope & redescription orbit vs invariant identity \\
Rescue route exhaustion & ERK20260527-New\_erk\_files\_May\_3D011A-THM-083 & Theorem & Rescue-route exhaustion & blocks semantic/reflection/richer-extension rescue \\
Role occupancy & ERK20260527-New\_erk\_files\_May\_DCCC53-THM-003 & Theorem & Forced Role Set & licensed realization of role occupation \\
Same-domain incompleteness & ERK20260527-New\_erk\_files\_May\_FBEDF2-THM-007 & Theorem & Exhaustion of same-domain incompleteness objection routes & Godel/independence firewall \\
Silent redescription & ERK20260527-New\_erk\_files\_May\_35B594-COR-007 & Corollary & No silent redescription & blocks surrogate/normalization replacement of target \\
Standing conservation & ERK20260527-New\_erk\_files\_May\_8B37FD-PRO-005 & Proposition & ; & standing can be preserved or destroyed, not created \\
Transport closure & ERK20260527-New\_erk\_files\_May\_12778C-PRO-001 & Proposition & Explicit mismatch under tensor omission & same-target transport closure \\
UEAP carrier adequacy & ERK20260527-New\_erk\_files\_May\_AF0C6B-CRT-001 & Criterion & Carrier Adequacy Test & report-status and carrier adequacy gate \\
UEAP laundering & ERK20260527-New\_erk\_files\_May\_AF0C6B-DEF-004 & Definition & Carrier & blocks report promotion from support to endpoint standing \\
\bottomrule
\end{longtable}


\begin{audit}[Corpus-use boundary]
The support ledger is not an authority import for RH.  It is a provenance record for structural constraints replayed locally in the manuscript.  A same-mode objection must attack the local use of a support family, not merely observe that the support family comes from the AASC corpus.
\end{audit}



\newpage
\section{Extended Hostile-Referee Defense Capsules}\label{app:extendedcapsules}

This section intentionally repeats the proof spine at high resolution.  Each capsule states a load-bearing claim, gives the local proof reason, records the hostile misreading it blocks, and names the denial burden.  The point is defensive redundancy: no transition used in the endpoint proof is allowed to remain implicit.

\subsection{Capsule 1: Context non-smuggling hard lock}
\begin{theorem}[Context non-smuggling hard lock]
The context object fixes the zeta/prime-trace carrier but does not assert RH, critical-line readout, no off-line branch, or no-independent amplitude closure.
\end{theorem}
\begin{proof}
The carrier fields name the analytic object, trace object, route, and role slots.  A slot is not an occupant.  If the context asserted universal critical-line support, it would be a theorem object rather than a context object.  Endpoint closure is derived only after projection-role content, endpoint counterforce, amplitude separator occupation, and no-independent closure are in force.
\end{proof}
\begin{audit}[Context non-smuggling hard lock anti-misreading]
Do not read the positive endpoint image as occupied merely because it is named.  The target must be fixed before it can be proved.
\end{audit}
\textbf{Denial burden.} Show a line in the context definition where universal critical-line support or absence of off-line zeros is asserted.

\subsection{Capsule 2: Official route is not truth}
\begin{theorem}[Official route is not truth]
The official explicit-formula route fixes endpoint-role projection; it does not prove either RH or its negation.
\end{theorem}
\begin{proof}
A route tells the proof where branch content is read.  In this manuscript the branch content is read through critical-normalized prime-trace amplitude.  The route makes non-neutral amplitude the projected role content of the negative branch, but it does not assert that such content exists or is impossible.
\end{proof}
\begin{audit}[Official route is not truth anti-misreading]
Projection is not proof.  The route makes the later classification meaningful; it is not the classification itself.
\end{audit}
\textbf{Denial burden.} Produce an inference in which OfficialEFRoute alone yields Lcrit or neg Lcrit.

\subsection{Capsule 3: OSEval is closure evaluation, not endpoint completion}
\begin{theorem}[OSEval is closure evaluation, not endpoint completion]
The endpoint-closure evaluation predicate fixes local countercase discipline without assuming an endpoint outcome.
\end{theorem}
\begin{proof}
Endpoint closure needs a fixed target and a rule for how a countercase can oppose it.  It does not need a prior declaration that one branch has already been proved.  The local reductio opens the complement branch and tests whether it can stand as counterforce.
\end{proof}
\begin{audit}[OSEval is closure evaluation, not endpoint completion anti-misreading]
This blocks the objection that the manuscript assumes endpoint-complete bivalence at report level.  Raw bivalence is separate.
\end{audit}
\textbf{Denial burden.} Show that OSEval contains PosEndpointOutcome or NegEndpointOutcome as a premise.

\subsection{Capsule 4: Raw bivalence remains bare bivalence}
\begin{theorem}[Raw bivalence remains bare bivalence]
The split Lcrit(rho) or not Lcrit(rho) is used only at the raw local proposition layer.
\end{theorem}
\begin{proof}
Classical logic provides a local complement split for a fixed proposition.  The proof does not infer endpoint standing from that split.  The assumption becomes endpoint-relevant only after it is annotated as the live exact-complement countercase and routed through official EF projection.
\end{proof}
\begin{audit}[Raw bivalence remains bare bivalence anti-misreading]
A raw assumption is not official endpoint resolution.  This is the anti-promotion rule that protects the proof from trivialization.
\end{audit}
\textbf{Denial burden.} Show a step where raw negation alone implies AmpSepLoc.

\subsection{Capsule 5: Endpoint image versus symmetric outcome carrier}
\begin{theorem}[Endpoint image versus symmetric outcome carrier]
The broad carrier {RH, not RH} is not the same as the critical-neutral amplitude endpoint image.
\end{theorem}
\begin{proof}
The broad carrier is a two-outcome problem presentation.  The endpoint image is the positive amplitude interface induced by the official EF route.  Off-line amplitude can be a negative broad outcome, but on the critical-neutral image it is non-occupation, not coequal occupation.
\end{proof}
\begin{audit}[Endpoint image versus symmetric outcome carrier anti-misreading]
This distinction does not deny the legitimacy of the negative problem branch.  It identifies its role on the chosen endpoint image.
\end{audit}
\textbf{Denial burden.} Give a same-image occupant that is non-neutral and nevertheless occupies the critical-neutral slot.

\subsection{Capsule 6: Kernel forcing is not optional framework import}
\begin{theorem}[Kernel forcing is not optional framework import]
The endpoint-use act already requires reference, standing, admissibility, and irreversibility.
\end{theorem}
\begin{proof}
An official endpoint resolution must fix what is being resolved, what counts as a valid step, whether a later repair is the same act, and how lawful redescription preserves identity.  Those are exactly the kernel roles.  The proof uses the role work, not the vocabulary as a premise.
\end{proof}
\begin{audit}[Kernel forcing is not optional framework import anti-misreading]
Rejecting AASC terminology does not remove target, step, finality, and fidelity requirements.
\end{audit}
\textbf{Denial burden.} Supply an ordinary official endpoint act preserving those features while lacking one kernel role.

\subsection{Capsule 7: Aplus is consequence not source}
\begin{theorem}[Aplus is consequence not source]
A+ does not create endpoint standing; it records fixed-domain consequences of endpoint use.
\end{theorem}
\begin{proof}
Endpoint use first instantiates the kernel.  The AMetric boundary, unique interior, no-carrier, no-repair, and no-independent-discriminator consequences then follow from fixed-domain closure.  The RH-EF specialization names the carrier on which these consequences apply.
\end{proof}
\begin{audit}[Aplus is consequence not source anti-misreading]
A+ is not a hidden premise inserted to prove RH.  It is activated only after endpoint use is present.
\end{audit}
\textbf{Denial burden.} Identify where A+ is used before EndpointUse or FixedDomain.

\subsection{Capsule 8: AMetric no-selector applies only at endpoint boundary}
\begin{theorem}[AMetric no-selector applies only at endpoint boundary]
Coordinates, exponents, cutoffs, and normalizations are not banned; their endpoint-gate use is controlled.
\end{theorem}
\begin{proof}
The AMetric boundary forbids selector authority at the admissibility boundary.  A coordinate may be described; an exponent may be calculated; a smoothing may be chosen.  They become prohibited only if they decide endpoint standing without a bridge.
\end{proof}
\begin{audit}[AMetric no-selector applies only at endpoint boundary anti-misreading]
This prevents the proof from banning analytic tools.
\end{audit}
\textbf{Denial burden.} Exhibit a passage where a descriptive coordinate or support-level exponent is declared invalid merely as data.

\subsection{Capsule 9: UEAP blocks report laundering}
\begin{theorem}[UEAP blocks report laundering]
A support-level amplitude report cannot be promoted to endpoint result without endpoint-counterforce and bridge discipline.
\end{theorem}
\begin{proof}
UEAP records claim target, carrier, support, and report status.  Non-neutral amplitude is support-level unless it is used as endpoint counterforce in the official route.  Reporting it as endpoint-standing without that use would exceed its licensed status.
\end{proof}
\begin{audit}[UEAP blocks report laundering anti-misreading]
UEAP is a no-overreporting discipline, not a zero-location theorem.
\end{audit}
\textbf{Denial burden.} Show a support-level report that the manuscript treats as endpoint result without the counterforce premise.

\subsection{Capsule 10: ATS locks anchor and tensor}
\begin{theorem}[ATS locks anchor and tensor]
Same-zero identity, the zeta carrier, the prime-trace carrier, and normalized amplitude exponent cannot be replaced by skin.
\end{theorem}
\begin{proof}
ATS distinguishes anchor, tensor, and skin.  The same zero and route are anchors; the normalized amplitude exponent is tensor content; notation, smoothing, and presentation are skin when they preserve that tensor.  Skin cannot carry endpoint force.
\end{proof}
\begin{audit}[ATS locks anchor and tensor anti-misreading]
This is why spectral analogies, random-matrix statistics, and finite computations cannot silently replace the EF amplitude role.
\end{audit}
\textbf{Denial burden.} Show a skin-level presentation change that the manuscript uses as tensor-level endpoint force.

\subsection{Capsule 11: No carrier transfer}
\begin{theorem}[No carrier transfer]
Spectral, numerical, positivity, density, and random-matrix carriers cannot transport endpoint standing into the zeta/prime-trace carrier without a bridge.
\end{theorem}
\begin{proof}
Each auxiliary carrier may be mathematically meaningful, but endpoint standing is carrier-bound.  Without an explicit bridge to the same zerohood and same amplitude endpoint image, the auxiliary carrier supplies support only.
\end{proof}
\begin{audit}[No carrier transfer anti-misreading]
The theorem protects analytic programs by classifying them as support or bridge programs rather than rejecting them.
\end{audit}
\textbf{Denial burden.} Produce a lawful unbridged carrier transfer that preserves endpoint standing on the same RH-EF carrier.

\subsection{Capsule 12: No post-selection repair}
\begin{theorem}[No post-selection repair]
Later analytic information cannot make an earlier unbridged endpoint report valid as the same act.
\end{theorem}
\begin{proof}
If a later proof supplies a missing bridge, the result is a fresh bridge-complete act.  It may vindicate content, but it does not retroactively give standing to the earlier unbridged act.
\end{proof}
\begin{audit}[No post-selection repair anti-misreading]
This is act-time finality, not hostility to later proof.
\end{audit}
\textbf{Denial burden.} Give a same-act repair that changes endpoint standing while preserving act identity.

\subsection{Capsule 13: Explicit-formula model adequacy is not RH}
\begin{theorem}[Explicit-formula model adequacy is not RH]
The model records zero coverage, contribution soundness, same-zero identity, amplitude exactness, and no-smuggling; it does not assert all zeros have beta=1/2.
\end{theorem}
\begin{proof}
The model makes both neutral and non-neutral amplitude roles meaningful.  A model that only allowed neutral contributions would already encode RH.  This manuscript instead allows OffLineAmpExcl as support-level content.
\end{proof}
\begin{audit}[Explicit-formula model adequacy is not RH anti-misreading]
The model is an adequacy packet, not a conclusion.
\end{audit}
\textbf{Denial burden.} Show that ModelEF includes a universal neutrality clause.

\subsection{Capsule 14: Amplitude exactness is algebraic role content}
\begin{theorem}[Amplitude exactness is algebraic role content]
For rho=beta+i gamma, critical normalization sends the contribution scale to \(x^{\beta-1/2}\).
\end{theorem}
\begin{proof}
The contribution scale \(x^\rho\) has magnitude scale \(x^\beta\).  Dividing by \(x^{1/2}\) leaves \(x^{\beta-1/2}\).  This establishes the tensor role used for projection.
\end{proof}
\begin{audit}[Amplitude exactness is algebraic role content anti-misreading]
This is not a dominance estimate or global error term.
\end{audit}
\textbf{Denial burden.} Attack must show the scale role is not correct for the explicit-formula contribution.

\subsection{Capsule 15: Term isolation does not assert dominance}
\begin{theorem}[Term isolation does not assert dominance]
An individual zero contribution may be assigned endpoint-relevant role content without claiming it dominates the prime trace.
\end{theorem}
\begin{proof}
Term isolation fixes same-zero identity, route, symmetry partners, and smoothing skin.  It lets the proof discuss role content without claiming global analytic dominance.
\end{proof}
\begin{audit}[Term isolation does not assert dominance anti-misreading]
This directly answers the objection that explicit formulae are sums with cancellations.
\end{audit}
\textbf{Denial burden.} Show that the proof requires dominance rather than role identification.

\subsection{Capsule 16: Pairing does not erase endpoint role}
\begin{theorem}[Pairing does not erase endpoint role]
Symmetry-paired off-line terms retain non-neutral scale classes unless a bridge proves endpoint-neutral occupation.
\end{theorem}
\begin{proof}
Functional-equation symmetry may pair beta with 1-beta.  Relative to critical normalization, those terms carry reciprocal non-neutral scale classes.  Pairing is lawful but not automatically neutral endpoint occupation.
\end{proof}
\begin{audit}[Pairing does not erase endpoint role anti-misreading]
Cancellation or pairing can be a bridge if proved; it is not a hidden exemption.
\end{audit}
\textbf{Denial burden.} Produce a same-carrier theorem that pairing alone yields critical-neutral endpoint occupation.

\subsection{Capsule 17: Smoothing and truncation are skin unless they change tensor}
\begin{theorem}[Smoothing and truncation are skin unless they change tensor]
Smoothing and truncation choices do not change endpoint role unless they change the normalized exponent content.
\end{theorem}
\begin{proof}
Equivalent formula presentations can manage convergence or boundary terms.  When they leave the same-zero exponent class invariant, they are skin.  When they alter the role, they are a different route requiring a bridge.
\end{proof}
\begin{audit}[Smoothing and truncation are skin unless they change tensor anti-misreading]
This preserves analytic rigor without letting presentation choices carry endpoint force.
\end{audit}
\textbf{Denial burden.} Identify a smoothing skin that the manuscript treats as an endpoint tensor.

\subsection{Capsule 18: Critical-neutral interface is target image}
\begin{theorem}[Critical-neutral interface is target image]
Neutral amplitude \(x^0\) is the positive endpoint image because the official endpoint asks beta=1/2.
\end{theorem}
\begin{proof}
Under the EF route, beta=1/2 corresponds to \(x^{\beta-1/2}=x^0\).  Thus neutral amplitude is the image of the positive endpoint.  This fixes the interface without asserting universal occupation.
\end{proof}
\begin{audit}[Critical-neutral interface is target image anti-misreading]
Target image fixation is not endpoint proof.
\end{audit}
\textbf{Denial burden.} Show where universal neutral occupation is asserted before the contradiction.

\subsection{Capsule 19: Off-line amplitude is non-occupation not contradiction}
\begin{theorem}[Off-line amplitude is non-occupation not contradiction]
OffLineAmpExcl is non-occupation of the critical-neutral image and is lawful until used as endpoint-standing separator work.
\end{theorem}
\begin{proof}
Non-neutral amplitude is the complement of neutral amplitude on the fixed endpoint image.  Complementhood alone is not contradiction.  Endpoint-standing use supplies the additional governance role.
\end{proof}
\begin{audit}[Off-line amplitude is non-occupation not contradiction anti-misreading]
This is the central non-explosiveness statement.
\end{audit}
\textbf{Denial burden.} Show an inference OffLineAmpExcl implies contradiction without endpoint counterforce.

\subsection{Capsule 20: Projection-role content is mediator}
\begin{theorem}[Projection-role content is mediator]
Projection-role content binds the countercase to its amplitude role but does not itself produce the forbidden discriminator.
\end{theorem}
\begin{proof}
The projection theorem says what the branch means on the official EF route.  The discriminator theorem requires endpoint counterforce and exclusion of support, bookkeeping, carrier shift, repair, and coequal occupation.
\end{proof}
\begin{audit}[Projection-role content is mediator anti-misreading]
This prevents one-step retyping.
\end{audit}
\textbf{Denial burden.} Show that ProjRoleEF alone implies Damp.

\subsection{Capsule 21: Descriptive amplitude remains lawful}
\begin{theorem}[Descriptive amplitude remains lawful]
The relation \(x^{\beta-1/2}\ne x^0\) may be written, analyzed, and used as support without contradiction.
\end{theorem}
\begin{proof}
Descriptive amplitude is below endpoint-standing level.  UEAP allows it as a report of scale; ATS locks it as tensor content.  It becomes dangerous only if used to defeat endpoint closure.
\end{proof}
\begin{audit}[Descriptive amplitude remains lawful anti-misreading]
This protects the paper from banning ordinary explicit-formula analysis.
\end{audit}
\textbf{Denial burden.} Find a theorem that forbids descriptive amplitude displacement.

\subsection{Capsule 22: Support amplitude remains lawful}
\begin{theorem}[Support amplitude remains lawful]
Support-level non-neutral amplitude is not the excluded object.
\end{theorem}
\begin{proof}
A hypothetical off-line zero may be analyzed through its explicit-formula contribution.  Such analysis can be support, counterexample syntax, or route diagnosis.  It is excluded only when it claims endpoint-standing authority on the same endpoint image.
\end{proof}
\begin{audit}[Support amplitude remains lawful anti-misreading]
This is why the proof does not reject all negative analytic methods.
\end{audit}
\textbf{Denial burden.} Show that support-level amplitude is identified with AmpSepLoc without counterforce.

\subsection{Capsule 23: Endpoint-standing non-neutral amplitude is status work}
\begin{theorem}[Endpoint-standing non-neutral amplitude is status work]
If non-neutral amplitude defeats endpoint closure on the same carrier, it performs noncritical amplitude status work.
\end{theorem}
\begin{proof}
Defeating endpoint closure is a status assignment: the local branch is treated as blocking the positive endpoint.  Since it is not support, bookkeeping, repair, carrier shift, or coequal image occupation, the act assigns same-domain noncritical status.
\end{proof}
\begin{audit}[Endpoint-standing non-neutral amplitude is status work anti-misreading]
This is a role classification, not an analytic estimate.
\end{audit}
\textbf{Denial burden.} Give an endpoint-defeating non-neutral amplitude act that changes no endpoint status.

\subsection{Capsule 24: Native data can be non-native as authority}
\begin{theorem}[Native data can be non-native as authority]
A datum may be native to analysis but non-native as endpoint authority over a fixed endpoint image.
\end{theorem}
\begin{proof}
The real part and amplitude exponent are native analytic data.  Endpoint authority is a different role.  When native data is used as a second classifier of endpoint standing, it becomes non-native relative to the fixed endpoint interface.
\end{proof}
\begin{audit}[Native data can be non-native as authority anti-misreading]
The native/non-native distinction concerns use, not origin.
\end{audit}
\textbf{Denial burden.} Show that native analytic origin automatically licenses endpoint-standing authority.

\subsection{Capsule 25: No RH-equivalent error bound is assumed}
\begin{theorem}[No RH-equivalent error bound is assumed]
The non-neutral scale standing theorem does not assert a prime-trace error bound equivalent to RH.
\end{theorem}
\begin{proof}
It does not say non-neutral contributions are absent.  It says endpoint-standing use of them is a same-domain discriminator.  Absence follows later only after the discriminator contradiction in the local reductio.
\end{proof}
\begin{audit}[No RH-equivalent error bound is assumed anti-misreading]
This blocks the charge that the proof smuggles the classical error-term equivalent of RH.
\end{audit}
\textbf{Denial burden.} Identify an assumed bound on psi(x)-x strong enough to imply RH.

\subsection{Capsule 26: Amplitude separator is not support}
\begin{theorem}[Amplitude separator is not support]
AmpSepLoc is non-neutral amplitude used as live endpoint-defeating occupation, not support-level amplitude.
\end{theorem}
\begin{proof}
The definition requires counterforce, fixed route, and exclusion of support/bookkeeping/shift/repair.  Therefore it starts where support ends.
\end{proof}
\begin{audit}[Amplitude separator is not support anti-misreading]
This prevents reading AmpSep as just \(x^{\beta-1/2}\ne x^0\).
\end{audit}
\textbf{Denial burden.} Find a use of AmpSepLoc without endpoint counterforce.

\subsection{Capsule 27: Endpoint-discriminator trichotomy is exhaustive}
\begin{theorem}[Endpoint-discriminator trichotomy is exhaustive]
An endpoint-status classifier is bookkeeping, governance-equivalent, independent same-domain work, or domain shift.
\end{theorem}
\begin{proof}
The classifier either changes endpoint governance or not; if not, it is bookkeeping.  If yes, it either preserves the same admissibility interface, leaves the domain, or performs independent same-domain work.  These cases exhaust the options.
\end{proof}
\begin{audit}[Endpoint-discriminator trichotomy is exhaustive anti-misreading]
No hidden fifth role remains by naming a classifier ordinary or analytic.
\end{audit}
\textbf{Denial burden.} Produce a fifth endpoint role that changes endpoint standing without falling into the trichotomy.

\subsection{Capsule 28: Noncritical amplitude work induces Damp}
\begin{theorem}[Noncritical amplitude work induces Damp]
Noncritical amplitude status work is the independent amplitude discriminator \(D_{\mathrm{amp}}\).
\end{theorem}
\begin{proof}
By definition it is endpoint-status work on the same critical-neutral image, not support, not bookkeeping, not equivalent to neutral occupation, not repair, and not carrier shift.  The trichotomy leaves independent same-domain amplitude discrimination.
\end{proof}
\begin{audit}[Noncritical amplitude work induces Damp anti-misreading]
This is the main hinge.  It is explicitly proved and audited rather than assumed.
\end{audit}
\textbf{Denial burden.} Give noncritical same-domain endpoint-status work that is not independent and not one of the excluded roles.

\subsection{Capsule 29: Endpoint use yields Namp}
\begin{theorem}[Endpoint use yields Namp]
Endpoint use of the local amplitude separator triggers the no-independent amplitude closure.
\end{theorem}
\begin{proof}
Endpoint use instantiates the kernel and A+.  A+ on the fixed RH-EF carrier includes no independent same-domain endpoint-status discriminator.  The local separator act is endpoint use inside the reductio subproof.
\end{proof}
\begin{audit}[Endpoint use yields Namp anti-misreading]
The closure does not apply to support-level amplitude, only to endpoint use.
\end{audit}
\textbf{Denial burden.} Show that AmpSepLoc is not endpoint use despite being endpoint-defeating occupation.

\subsection{Capsule 30: Damp and Namp contradiction is local}
\begin{theorem}[Damp and Namp contradiction is local]
The contradiction is local to the annotated countercase subproof.
\end{theorem}
\begin{proof}
The local assumption induces AmpSepLoc, which yields Damp and Namp.  The contradiction occurs under the open assumption and discharges it.  The proof does not require a global negative endpoint outcome.
\end{proof}
\begin{audit}[Damp and Namp contradiction is local anti-misreading]
This prevents raw case promotion.
\end{audit}
\textbf{Denial burden.} Show that the proof asserts NegEndpointOutcome before local discharge.

\subsection{Capsule 31: Exact-complement annotation is not a premise}
\begin{theorem}[Exact-complement annotation is not a premise]
ReductioCountercase is a proof-act annotation of the sole local assumption, not a second object-level assumption.
\end{theorem}
\begin{proof}
Ordinary reductio has an assumption and a role for that assumption.  The annotation records that the assumption is the live exact complement of the target endpoint readout.  No new formula is added to the object language.
\end{proof}
\begin{audit}[Exact-complement annotation is not a premise anti-misreading]
This blocks the strengthened-assumption objection.
\end{audit}
\textbf{Denial burden.} Show that the discharged assumption is neg Lcrit conjoined with annotation.

\subsection{Capsule 32: Local countercase standing is temporary}
\begin{theorem}[Local countercase standing is temporary]
The local off-line countercase has temporary derivational standing only inside the reductio.
\end{theorem}
\begin{proof}
A local countercase is not a theorem of the negative branch.  It is a proof-act role sufficient to test endpoint closure.  Once contradiction is reached, the assumption is discharged.
\end{proof}
\begin{audit}[Local countercase standing is temporary anti-misreading]
Temporary standing is not global endpoint resolution.
\end{audit}
\textbf{Denial burden.} Show that LocalCountercase implies official negative endpoint resolution.

\subsection{Capsule 33: Endpoint counterforce is narrower than truth}
\begin{theorem}[Endpoint counterforce is narrower than truth]
A branch has endpoint counterforce only when used to defeat endpoint closure on the fixed carrier.
\end{theorem}
\begin{proof}
Raw truth, hypothetical syntax, and support content may exist without endpoint counterforce.  The exact-complement countercase has counterforce because it is used against the endpoint target.
\end{proof}
\begin{audit}[Endpoint counterforce is narrower than truth anti-misreading]
This is the local analogue of official negative endpoint use, not a truth-to-standing rule.
\end{audit}
\textbf{Denial burden.} Give a raw branch that defeats endpoint closure while lacking counterforce.

\subsection{Capsule 34: Annotation discharge returns Lcrit}
\begin{theorem}[Annotation discharge returns Lcrit]
Contradiction under the sole local assumption neg Lcrit(rho) discharges to Lcrit(rho).
\end{theorem}
\begin{proof}
Classical reductio discharges the open assumption when contradiction is derived.  Since the annotation is not an object-level premise, there is no strengthened conjunction to discharge.
\end{proof}
\begin{audit}[Annotation discharge returns Lcrit anti-misreading]
This is ordinary proof discipline with role annotation.
\end{audit}
\textbf{Denial burden.} Show that an additional object-level assumption remains undischarged.

\subsection{Capsule 35: Universalization is over arbitrary zerohood}
\begin{theorem}[Universalization is over arbitrary zerohood]
The proof fixes arbitrary rho under Az(rho), derives Lcrit(rho), and universally generalizes.
\end{theorem}
\begin{proof}
No special property of rho is used other than nontrivial zerohood and the fixed EF route.  Therefore the local result is universal over the RH carrier.
\end{proof}
\begin{audit}[Universalization is over arbitrary zerohood anti-misreading]
The proof does not infer from checked zeros or finite ranges.
\end{audit}
\textbf{Denial burden.} Identify a hidden restriction on rho beyond Az(rho).

\subsection{Capsule 36: Ordinary negative theoremhood is not banned}
\begin{theorem}[Ordinary negative theoremhood is not banned]
A theorem of not RH is not rejected merely because it is negative.
\end{theorem}
\begin{proof}
Ordinary theoremhood is proof legitimacy.  The excluded role is endpoint-standing amplitude separator occupation after official route projection.  If a negative theorem is not used as official endpoint resolution, it may remain ordinary mathematics or support.
\end{proof}
\begin{audit}[Ordinary negative theoremhood is not banned anti-misreading]
This protects negative analytic results and hypothetical counterexamples as syntax.
\end{audit}
\textbf{Denial burden.} Show that negative theoremhood is declared invalid simply for being negative.

\subsection{Capsule 37: Official negative resolution becomes endpoint use}
\begin{theorem}[Official negative resolution becomes endpoint use]
When a negative theorem is offered as official RH endpoint resolution, it enters endpoint-use discipline.
\end{theorem}
\begin{proof}
Official resolution fixes target, carrier, branch, theorem status, downstream reuse, and finality.  That is endpoint use.  Endpoint use activates the kernel and A+ consequences.
\end{proof}
\begin{audit}[Official negative resolution becomes endpoint use anti-misreading]
Ordinary theoremhood supplies legitimacy, not exemption from endpoint role.
\end{audit}
\textbf{Denial burden.} Give official endpoint force without endpoint use.

\subsection{Capsule 38: Incompleteness objections are endpoint-inert or governed}
\begin{theorem}[Incompleteness objections are endpoint-inert or governed]
Independence, true-but-unprovable, semantic, or stronger-theory distinctions matter only if they affect endpoint standing.
\end{theorem}
\begin{proof}
If they do not affect standing, reportability, reuse, continuation, or branch exclusion, they are endpoint-inert.  If they do, they are endpoint-admissibility-relevant gate work and fall under A+.
\end{proof}
\begin{audit}[Incompleteness objections are endpoint-inert or governed anti-misreading]
This does not deny classical incompleteness.
\end{audit}
\textbf{Denial burden.} Produce a same-domain incompleteness distinction that changes endpoint standing while escaping gate classification.

\subsection{Capsule 39: Open status is not endpoint outcome}
\begin{theorem}[Open status is not endpoint outcome]
Open, unresolved, or under-review status is epistemic and not a third endpoint branch.
\end{theorem}
\begin{proof}
The endpoint branches are RH and not RH.  Open status describes the state of inquiry, not the truth or endpoint occupation branch.  Treating it as a branch changes target.
\end{proof}
\begin{audit}[Open status is not endpoint outcome anti-misreading]
This prevents epistemic status from becoming hidden fifth endpoint status.
\end{audit}
\textbf{Denial burden.} Show a legitimate third endpoint truth branch in the official RH target.

\subsection{Capsule 40: Route-locus exhaustion is role-complete}
\begin{theorem}[Route-locus exhaustion is role-complete]
Every endpoint-resolving negative RH-EF route enters one of the audited loci.
\end{theorem}
\begin{proof}
A negative route must use coordinate report, EF support, amplitude separator, spectral surrogate, positivity criterion, finite verification, density estimate, zero-free region, symmetry orbit, pairing/cancellation, carrier shift, repair, or hidden fifth role.  If it enters none, it does no endpoint work.
\end{proof}
\begin{audit}[Route-locus exhaustion is role-complete anti-misreading]
The table is not decoration; it is the exhaustion proof in operational form.
\end{audit}
\textbf{Denial burden.} Produce a same-carrier endpoint-resolving negative route outside the loci.

\subsection{Capsule 41: Weakening resistance}
\begin{theorem}[Weakening resistance]
A weaker same-carrier regime either preserves the endpoint behavior or reintroduces forbidden amplitude governance.
\end{theorem}
\begin{proof}
If the weakening changes no endpoint verdict, it is bookkeeping.  If it permits endpoint-standing non-neutral amplitude, it permits independent discriminator work.  If it shifts carrier, it is not same endpoint.
\end{proof}
\begin{audit}[Weakening resistance anti-misreading]
This prevents weakening the packet just enough to let the separator through.
\end{audit}
\textbf{Denial burden.} Define a weaker fixed-carrier regime preserving all endpoint roles while allowing AmpSepLoc without Damp.

\subsection{Capsule 42: Faithful counterexample burden}
\begin{theorem}[Faithful counterexample burden]
A valid objection must produce a concrete faithful counterexample to one named theorem.
\end{theorem}
\begin{proof}
General discomfort with AASC, lack of estimates, or unfamiliar vocabulary does not attack the theorem chain.  A mathematical objection must break context adequacy, EF model adequacy, projection, scale standing, trichotomy, A+ closure, local discharge, or universalization.
\end{proof}
\begin{audit}[Faithful counterexample burden anti-misreading]
This is not burden shifting; it follows from the proof structure.
\end{audit}
\textbf{Denial burden.} Supply one named theorem failure with a same-domain counterexample.


\section{Transition-by-Transition Audit}

The following audit isolates the exact proof transition being used at each stage of the final local reductio.  Each transition must be attacked directly if a hostile referee wishes to break the proof.

\begin{longtable}{L{0.10\textwidth} L{0.27\textwidth} L{0.38\textwidth} L{0.15\textwidth}}
\toprule
Code & Transition & Authorized by & Main forbidden misreading \\
\midrule
\endhead
T1 & \(\Az(\rho)\Rightarrow \Lcrit(\rho)\vee\neg\Lcrit(\rho)\) & Raw classical bivalence under fixed zerohood & Endpoint standing from raw truth \\
T2 & \(\neg\Lcrit(\rho)\Rightarrow \beta\ne1/2\) & Definition of critical-line readout & Off-line coordinate declared illicit \\
T3 & \(\beta\ne1/2\Rightarrow x^{\beta-1/2}\ne x^0\) & EF amplitude exactness & Contradiction from amplitude alone \\
T4 & Off-line amplitude to projection-role content & OfficialEFRoute and ModelEF & One-step separator typing \\
T5 & Projection-role content plus counterforce to endpoint-standing non-neutral amplitude & Endpoint-standing criterion & Support promoted to endpoint result \\
T6 & Endpoint-standing non-neutral amplitude to noncritical status work & Non-neutral scale standing theorem & Native analytic data banned as data \\
T7 & Noncritical status work to \(\Damp\) & Endpoint-discriminator trichotomy & Hidden fifth endpoint role \\
T8 & \(\AmpSepLoc\Rightarrow\EndpointUse\) & Definition of local separator occupation & Support-level use called endpoint use \\
T9 & Endpoint use to \(\Namp\) & Kernel and A+ closure & A+ applied before endpoint use \\
T10 & \(\Namp\Rightarrow\neg\Damp\) & No-independent amplitude closure & Descriptive labels banned \\
T11 & \(\Damp\wedge\neg\Damp\Rightarrow\bot\) & Classical contradiction & None \\
T12 & Contradiction discharges \([\neg\Lcrit]_i\) & Annotation-discharge rule & Strengthened-assumption discharge \\
T13 & Arbitrary zerohood to universal RH & Universal generalization & Finite or local-only proof \\
\bottomrule
\end{longtable}

\section{Expanded Route-Locus Exhaustion}

\begin{longtable}{L{0.20\textwidth} L{0.34\textwidth} L{0.34\textwidth}}
\toprule
Locus & RH-EF form & Disposition \\
\midrule
\endhead
Coordinate-only report & \(\Ree(\rho)\ne1/2\) as coordinate data & Lawful analytic data; not endpoint standing. \\
Explicit-formula support & \(x^{\beta-1/2}\ne x^0\) as contribution scale & Lawful support; not contradiction. \\
Endpoint amplitude separator & Non-neutral amplitude used as endpoint counterforce & Routes to \(\Damp\) and is excluded by \(\Namp\). \\
Symmetry orbit & \(\rho,1-\rho,\bar\rho,1-\bar\rho\) & Lawful carrier structure; not same-zero neutral occupation without bridge. \\
Pairing/cancellation & Paired contributions cancel or oscillate & Bridge if proved; otherwise not endpoint-neutral occupation. \\
Spectral surrogate & Hilbert--Polya or model spectrum & Support or bridge program; not same endpoint without carrier identification. \\
Positivity criterion & Equivalent or proposed positivity condition & Bridge if fully proved; surrogate if not. \\
Numerical verification & Finite-height zero checking & Finite support; not universal endpoint. \\
Density estimate & Off-line zeros rare or bounded & Support; absence must be proved. \\
Zero-free region & Region excludes zeros & Support or bridge piece; not universal endpoint unless covers all off-line region. \\
Carrier shift & Replacement by another domain & Not same RH-EF endpoint. \\
Repair & Later proof validates earlier report as same act & Violates act-time finality. \\
Hidden fifth role & Endpoint-resolving but non-governing off-line amplitude & Impossible by trichotomy. \\
\bottomrule
\end{longtable}

\section{Expanded Right-Mode Objection Targets}

A right-mode objection must attack at least one of the following named points: context adequacy, context non-smuggling, official route fixation, EF model adequacy, amplitude exactness, term isolation, pairing closure, projection-role content, single amplitude interface, non-explosiveness, endpoint-standing criterion, non-neutral scale standing, non-native amplitude theorem, endpoint-discriminator trichotomy, A+ no-independent closure, exact-complement annotation discipline, annotation discharge, or universalization.

An objection that demands a zero-free region, density estimate, or spectral operator is wrong-mode unless it also explains why this proof class is not a legitimate endpoint-exclusion class.  An objection that says ``off-line zeros are meaningful'' is already granted.  An objection that says ``non-neutral amplitude is native analytic data'' is also granted.  The proof excludes only endpoint-standing use of that data as independent amplitude separator authority on the fixed endpoint image.

\section{Formalization Map}\label{app:formalization}

A future Lean-facing support layer would use objects with the following shape:
\begin{verbatim}
def NontrivialZerohood (rho : Complex) : Prop := zeta rho = 0 /\ 0 < re rho /\ re rho < 1

def CriticalLineReadout (rho : Complex) : Prop := re rho = (1/2 : Real)

def PrimeTraceCarrier := ...
def OfficialEFRoute := ...
def EFContributionScale (rho : Complex) := x ^ (re rho - 1/2)

def CritAmpOcc (rho : Complex) : Prop := EFContributionScale rho = x ^ 0

def OffLineAmpExcl (rho : Complex) : Prop := EFContributionScale rho <> x ^ 0

def ProjectionRoleContentEF (B : Prop) (rho : Complex) : Prop := ...
def LocalAmpSeparatorOccupation (rho : Complex) : Prop := ...
def AmpEndpointDiscriminator (rho : Complex) : Prop := ...
def NoIndependentAmpDiscriminator (rho : Complex) : Prop := ...
\end{verbatim}

Key theorem surfaces:
\begin{verbatim}
ef_branch_projection
single_amplitude_interface
descriptive_non_neutral_amplitude_lawful
nonneutral_scale_standing
noncritical_amp_work_induces_discriminator
endpoint_use_forces_no_independent_amp
local_amp_separator_impossible
annotation_discharge
critical_line_readout_of_zerohood
riemann_hypothesis_endpoint
\end{verbatim}

\clearpage
\section{Detailed Load-Bearing Proof Capsules}\label{app:detailed-capsules}
This appendix expands the theorem ladder into individual proof capsules.  Each capsule states what the manuscript uses, why the use is not a hidden RH assumption, how the hostile reading is blocked, and what a successful same-mode objection must supply.  This section is intentionally redundant: a Clay-facing endpoint-exclusion proof cannot rely on a compressed proof spine.

\subsection{Capsule 1: Target fixation without conclusion}
\textbf{Load-bearing role.} This capsule controls the official endpoint is fixed as a target role while the positive branch is not asserted.  It is used because the proof would otherwise be vulnerable to the hostile reading that an analytic datum, route choice, or proof annotation is silently doing endpoint work before its status has been fixed.

\textbf{Dependency chain.} The local dependency is context non-smuggling and target adequacy.  No conclusion about the universal location of zeros is available at this stage.  The capsule only fixes the status of the object under its declared carrier and route.  If the object remains descriptive or support-level, it remains lawful.  If it becomes endpoint-standing, it enters the endpoint-status discipline.

\textbf{Anti-smuggling guard.} The capsule must not be read as assuming RH, as denying the explicit formula, or as banning ordinary analytic counterexample syntax.  Its only exclusion is the role-confused use of an object at a stronger report level than its carrier and bridge license.  This is why the proof repeatedly separates coordinate, support, endpoint-counterforce, and endpoint-outcome roles.

\textbf{Hostile-referee stress test.} A hostile referee will try to collapse this capsule into either an analytic estimate claim or a tautological role claim.  The response is that the capsule makes no estimate claim and no tautological zero-location claim.  It fixes the role boundary needed before the later discriminator theorem can be applied.

\textbf{Denial burden.} To defeat this capsule, the objection must identify a premise that asserts critical-line occupation before the local countercase is discharged.  A mere statement that the method is not analytic-native does not touch the capsule, because the proof class has already been locked as endpoint-rigidity rather than analytic estimation.

\subsection{Capsule 2: Official EF route versus analytic proof}
\textbf{Load-bearing role.} This capsule controls the explicit formula is an endpoint route rather than a new estimate or zero-free region.  It is used because the proof would otherwise be vulnerable to the hostile reading that an analytic datum, route choice, or proof annotation is silently doing endpoint work before its status has been fixed.

\textbf{Dependency chain.} The local dependency is official route fixation and EF model adequacy.  No conclusion about the universal location of zeros is available at this stage.  The capsule only fixes the status of the object under its declared carrier and route.  If the object remains descriptive or support-level, it remains lawful.  If it becomes endpoint-standing, it enters the endpoint-status discipline.

\textbf{Anti-smuggling guard.} The capsule must not be read as assuming RH, as denying the explicit formula, or as banning ordinary analytic counterexample syntax.  Its only exclusion is the role-confused use of an object at a stronger report level than its carrier and bridge license.  This is why the proof repeatedly separates coordinate, support, endpoint-counterforce, and endpoint-outcome roles.

\textbf{Hostile-referee stress test.} A hostile referee will try to collapse this capsule into either an analytic estimate claim or a tautological role claim.  The response is that the capsule makes no estimate claim and no tautological zero-location claim.  It fixes the role boundary needed before the later discriminator theorem can be applied.

\textbf{Denial burden.} To defeat this capsule, the objection must show that using the explicit formula as a projection route already assumes RH.  A mere statement that the method is not analytic-native does not touch the capsule, because the proof class has already been locked as endpoint-rigidity rather than analytic estimation.

\subsection{Capsule 3: Same-zero identity}
\textbf{Load-bearing role.} This capsule controls a contribution attributed to rho remains the contribution of that zerohood instance.  It is used because the proof would otherwise be vulnerable to the hostile reading that an analytic datum, route choice, or proof annotation is silently doing endpoint work before its status has been fixed.

\textbf{Dependency chain.} The local dependency is ATS anchor preservation and EF model same-zero clause.  No conclusion about the universal location of zeros is available at this stage.  The capsule only fixes the status of the object under its declared carrier and route.  If the object remains descriptive or support-level, it remains lawful.  If it becomes endpoint-standing, it enters the endpoint-status discipline.

\textbf{Anti-smuggling guard.} The capsule must not be read as assuming RH, as denying the explicit formula, or as banning ordinary analytic counterexample syntax.  Its only exclusion is the role-confused use of an object at a stronger report level than its carrier and bridge license.  This is why the proof repeatedly separates coordinate, support, endpoint-counterforce, and endpoint-outcome roles.

\textbf{Hostile-referee stress test.} A hostile referee will try to collapse this capsule into either an analytic estimate claim or a tautological role claim.  The response is that the capsule makes no estimate claim and no tautological zero-location claim.  It fixes the role boundary needed before the later discriminator theorem can be applied.

\textbf{Denial burden.} To defeat this capsule, the objection must produce a lawful route where a symmetry partner or averaged class replaces rho without carrier change.  A mere statement that the method is not analytic-native does not touch the capsule, because the proof class has already been locked as endpoint-rigidity rather than analytic estimation.

\subsection{Capsule 4: Smoothing and truncation skin}
\textbf{Load-bearing role.} This capsule controls formula presentation choices do not carry endpoint authority.  It is used because the proof would otherwise be vulnerable to the hostile reading that an analytic datum, route choice, or proof annotation is silently doing endpoint work before its status has been fixed.

\textbf{Dependency chain.} The local dependency is ATS skin/tensor separation.  No conclusion about the universal location of zeros is available at this stage.  The capsule only fixes the status of the object under its declared carrier and route.  If the object remains descriptive or support-level, it remains lawful.  If it becomes endpoint-standing, it enters the endpoint-status discipline.

\textbf{Anti-smuggling guard.} The capsule must not be read as assuming RH, as denying the explicit formula, or as banning ordinary analytic counterexample syntax.  Its only exclusion is the role-confused use of an object at a stronger report level than its carrier and bridge license.  This is why the proof repeatedly separates coordinate, support, endpoint-counterforce, and endpoint-outcome roles.

\textbf{Hostile-referee stress test.} A hostile referee will try to collapse this capsule into either an analytic estimate claim or a tautological role claim.  The response is that the capsule makes no estimate claim and no tautological zero-location claim.  It fixes the role boundary needed before the later discriminator theorem can be applied.

\textbf{Denial burden.} To defeat this capsule, the objection must show a smoothing convention that changes endpoint standing while remaining skin.  A mere statement that the method is not analytic-native does not touch the capsule, because the proof class has already been locked as endpoint-rigidity rather than analytic estimation.

\subsection{Capsule 5: Critical normalization}
\textbf{Load-bearing role.} This capsule controls normalization by \(x^{1/2}\) names the critical-line endpoint scale without asserting universal occupancy.  It is used because the proof would otherwise be vulnerable to the hostile reading that an analytic datum, route choice, or proof annotation is silently doing endpoint work before its status has been fixed.

\textbf{Dependency chain.} The local dependency is single amplitude interface and context non-smuggling.  No conclusion about the universal location of zeros is available at this stage.  The capsule only fixes the status of the object under its declared carrier and route.  If the object remains descriptive or support-level, it remains lawful.  If it becomes endpoint-standing, it enters the endpoint-status discipline.

\textbf{Anti-smuggling guard.} The capsule must not be read as assuming RH, as denying the explicit formula, or as banning ordinary analytic counterexample syntax.  Its only exclusion is the role-confused use of an object at a stronger report level than its carrier and bridge license.  This is why the proof repeatedly separates coordinate, support, endpoint-counterforce, and endpoint-outcome roles.

\textbf{Hostile-referee stress test.} A hostile referee will try to collapse this capsule into either an analytic estimate claim or a tautological role claim.  The response is that the capsule makes no estimate claim and no tautological zero-location claim.  It fixes the role boundary needed before the later discriminator theorem can be applied.

\textbf{Denial burden.} To defeat this capsule, the objection must show that naming the positive endpoint image is equivalent to assuming it is universally occupied.  A mere statement that the method is not analytic-native does not touch the capsule, because the proof class has already been locked as endpoint-rigidity rather than analytic estimation.

\subsection{Capsule 6: Amplitude exactness}
\textbf{Load-bearing role.} This capsule controls \(\rho=\beta+i\gamma\) produces critical-normalized scale class \(x^{\beta-1/2}\).  It is used because the proof would otherwise be vulnerable to the hostile reading that an analytic datum, route choice, or proof annotation is silently doing endpoint work before its status has been fixed.

\textbf{Dependency chain.} The local dependency is explicit-formula contribution adequacy.  No conclusion about the universal location of zeros is available at this stage.  The capsule only fixes the status of the object under its declared carrier and route.  If the object remains descriptive or support-level, it remains lawful.  If it becomes endpoint-standing, it enters the endpoint-status discipline.

\textbf{Anti-smuggling guard.} The capsule must not be read as assuming RH, as denying the explicit formula, or as banning ordinary analytic counterexample syntax.  Its only exclusion is the role-confused use of an object at a stronger report level than its carrier and bridge license.  This is why the proof repeatedly separates coordinate, support, endpoint-counterforce, and endpoint-outcome roles.

\textbf{Hostile-referee stress test.} A hostile referee will try to collapse this capsule into either an analytic estimate claim or a tautological role claim.  The response is that the capsule makes no estimate claim and no tautological zero-location claim.  It fixes the role boundary needed before the later discriminator theorem can be applied.

\textbf{Denial burden.} To defeat this capsule, the objection must attack the contribution-scale bridge rather than the AASC closure.  A mere statement that the method is not analytic-native does not touch the capsule, because the proof class has already been locked as endpoint-rigidity rather than analytic estimation.

\subsection{Capsule 7: Positive projection}
\textbf{Load-bearing role.} This capsule controls critical-line readout is neutral amplitude occupation.  It is used because the proof would otherwise be vulnerable to the hostile reading that an analytic datum, route choice, or proof annotation is silently doing endpoint work before its status has been fixed.

\textbf{Dependency chain.} The local dependency is amplitude exactness and EF route.  No conclusion about the universal location of zeros is available at this stage.  The capsule only fixes the status of the object under its declared carrier and route.  If the object remains descriptive or support-level, it remains lawful.  If it becomes endpoint-standing, it enters the endpoint-status discipline.

\textbf{Anti-smuggling guard.} The capsule must not be read as assuming RH, as denying the explicit formula, or as banning ordinary analytic counterexample syntax.  Its only exclusion is the role-confused use of an object at a stronger report level than its carrier and bridge license.  This is why the proof repeatedly separates coordinate, support, endpoint-counterforce, and endpoint-outcome roles.

\textbf{Hostile-referee stress test.} A hostile referee will try to collapse this capsule into either an analytic estimate claim or a tautological role claim.  The response is that the capsule makes no estimate claim and no tautological zero-location claim.  It fixes the role boundary needed before the later discriminator theorem can be applied.

\textbf{Denial burden.} To defeat this capsule, the objection must produce a critical-line zero whose contribution is non-neutral under the fixed normalization.  A mere statement that the method is not analytic-native does not touch the capsule, because the proof class has already been locked as endpoint-rigidity rather than analytic estimation.

\subsection{Capsule 8: Negative projection}
\textbf{Load-bearing role.} This capsule controls off-line readout is non-neutral amplitude exclusion.  It is used because the proof would otherwise be vulnerable to the hostile reading that an analytic datum, route choice, or proof annotation is silently doing endpoint work before its status has been fixed.

\textbf{Dependency chain.} The local dependency is amplitude exactness and branch complementarity.  No conclusion about the universal location of zeros is available at this stage.  The capsule only fixes the status of the object under its declared carrier and route.  If the object remains descriptive or support-level, it remains lawful.  If it becomes endpoint-standing, it enters the endpoint-status discipline.

\textbf{Anti-smuggling guard.} The capsule must not be read as assuming RH, as denying the explicit formula, or as banning ordinary analytic counterexample syntax.  Its only exclusion is the role-confused use of an object at a stronger report level than its carrier and bridge license.  This is why the proof repeatedly separates coordinate, support, endpoint-counterforce, and endpoint-outcome roles.

\textbf{Hostile-referee stress test.} A hostile referee will try to collapse this capsule into either an analytic estimate claim or a tautological role claim.  The response is that the capsule makes no estimate claim and no tautological zero-location claim.  It fixes the role boundary needed before the later discriminator theorem can be applied.

\textbf{Denial burden.} To defeat this capsule, the objection must produce an off-line zero whose normalized contribution still occupies \(x^0\) without a bridge.  A mere statement that the method is not analytic-native does not touch the capsule, because the proof class has already been locked as endpoint-rigidity rather than analytic estimation.

\subsection{Capsule 9: Projection-role content}
\textbf{Load-bearing role.} This capsule controls a live endpoint countercase is bound to its EF amplitude role.  It is used because the proof would otherwise be vulnerable to the hostile reading that an analytic datum, route choice, or proof annotation is silently doing endpoint work before its status has been fixed.

\textbf{Dependency chain.} The local dependency is official route plus reductio countercase annotation.  No conclusion about the universal location of zeros is available at this stage.  The capsule only fixes the status of the object under its declared carrier and route.  If the object remains descriptive or support-level, it remains lawful.  If it becomes endpoint-standing, it enters the endpoint-status discipline.

\textbf{Anti-smuggling guard.} The capsule must not be read as assuming RH, as denying the explicit formula, or as banning ordinary analytic counterexample syntax.  Its only exclusion is the role-confused use of an object at a stronger report level than its carrier and bridge license.  This is why the proof repeatedly separates coordinate, support, endpoint-counterforce, and endpoint-outcome roles.

\textbf{Hostile-referee stress test.} A hostile referee will try to collapse this capsule into either an analytic estimate claim or a tautological role claim.  The response is that the capsule makes no estimate claim and no tautological zero-location claim.  It fixes the role boundary needed before the later discriminator theorem can be applied.

\textbf{Denial burden.} To defeat this capsule, the objection must show the live countercase can defeat the endpoint while refusing its projected route content.  A mere statement that the method is not analytic-native does not touch the capsule, because the proof class has already been locked as endpoint-rigidity rather than analytic estimation.

\subsection{Capsule 10: Descriptive coordinate preservation}
\textbf{Load-bearing role.} This capsule controls Re(rho) not equal 1/2 remains lawful coordinate data.  It is used because the proof would otherwise be vulnerable to the hostile reading that an analytic datum, route choice, or proof annotation is silently doing endpoint work before its status has been fixed.

\textbf{Dependency chain.} The local dependency is UEAP report-level separation.  No conclusion about the universal location of zeros is available at this stage.  The capsule only fixes the status of the object under its declared carrier and route.  If the object remains descriptive or support-level, it remains lawful.  If it becomes endpoint-standing, it enters the endpoint-status discipline.

\textbf{Anti-smuggling guard.} The capsule must not be read as assuming RH, as denying the explicit formula, or as banning ordinary analytic counterexample syntax.  Its only exclusion is the role-confused use of an object at a stronger report level than its carrier and bridge license.  This is why the proof repeatedly separates coordinate, support, endpoint-counterforce, and endpoint-outcome roles.

\textbf{Hostile-referee stress test.} A hostile referee will try to collapse this capsule into either an analytic estimate claim or a tautological role claim.  The response is that the capsule makes no estimate claim and no tautological zero-location claim.  It fixes the role boundary needed before the later discriminator theorem can be applied.

\textbf{Denial burden.} To defeat this capsule, the objection must identify a line where descriptive coordinate data is itself declared illicit.  A mere statement that the method is not analytic-native does not touch the capsule, because the proof class has already been locked as endpoint-rigidity rather than analytic estimation.

\subsection{Capsule 11: Descriptive amplitude preservation}
\textbf{Load-bearing role.} This capsule controls \(x^{\beta-1/2}\ne x^0\) remains lawful support-level EF data.  It is used because the proof would otherwise be vulnerable to the hostile reading that an analytic datum, route choice, or proof annotation is silently doing endpoint work before its status has been fixed.

\textbf{Dependency chain.} The local dependency is UEAP report-level separation and non-explosiveness.  No conclusion about the universal location of zeros is available at this stage.  The capsule only fixes the status of the object under its declared carrier and route.  If the object remains descriptive or support-level, it remains lawful.  If it becomes endpoint-standing, it enters the endpoint-status discipline.

\textbf{Anti-smuggling guard.} The capsule must not be read as assuming RH, as denying the explicit formula, or as banning ordinary analytic counterexample syntax.  Its only exclusion is the role-confused use of an object at a stronger report level than its carrier and bridge license.  This is why the proof repeatedly separates coordinate, support, endpoint-counterforce, and endpoint-outcome roles.

\textbf{Hostile-referee stress test.} A hostile referee will try to collapse this capsule into either an analytic estimate claim or a tautological role claim.  The response is that the capsule makes no estimate claim and no tautological zero-location claim.  It fixes the role boundary needed before the later discriminator theorem can be applied.

\textbf{Denial burden.} To defeat this capsule, the objection must identify a line where support-level amplitude alone is contradicted.  A mere statement that the method is not analytic-native does not touch the capsule, because the proof class has already been locked as endpoint-rigidity rather than analytic estimation.

\subsection{Capsule 12: Endpoint-standing criterion}
\textbf{Load-bearing role.} This capsule controls non-neutral amplitude becomes excluded only as endpoint-standing counterforce.  It is used because the proof would otherwise be vulnerable to the hostile reading that an analytic datum, route choice, or proof annotation is silently doing endpoint work before its status has been fixed.

\textbf{Dependency chain.} The local dependency is counterforce definition and UEAP promotion discipline.  No conclusion about the universal location of zeros is available at this stage.  The capsule only fixes the status of the object under its declared carrier and route.  If the object remains descriptive or support-level, it remains lawful.  If it becomes endpoint-standing, it enters the endpoint-status discipline.

\textbf{Anti-smuggling guard.} The capsule must not be read as assuming RH, as denying the explicit formula, or as banning ordinary analytic counterexample syntax.  Its only exclusion is the role-confused use of an object at a stronger report level than its carrier and bridge license.  This is why the proof repeatedly separates coordinate, support, endpoint-counterforce, and endpoint-outcome roles.

\textbf{Hostile-referee stress test.} A hostile referee will try to collapse this capsule into either an analytic estimate claim or a tautological role claim.  The response is that the capsule makes no estimate claim and no tautological zero-location claim.  It fixes the role boundary needed before the later discriminator theorem can be applied.

\textbf{Denial burden.} To defeat this capsule, the objection must supply endpoint-defeating use that remains only support.  A mere statement that the method is not analytic-native does not touch the capsule, because the proof class has already been locked as endpoint-rigidity rather than analytic estimation.

\subsection{Capsule 13: ATS anchor lock}
\textbf{Load-bearing role.} This capsule controls zeta carrier, prime trace, same zero, route, and normalization are fixed anchors.  It is used because the proof would otherwise be vulnerable to the hostile reading that an analytic datum, route choice, or proof annotation is silently doing endpoint work before its status has been fixed.

\textbf{Dependency chain.} The local dependency is ATS local packet.  No conclusion about the universal location of zeros is available at this stage.  The capsule only fixes the status of the object under its declared carrier and route.  If the object remains descriptive or support-level, it remains lawful.  If it becomes endpoint-standing, it enters the endpoint-status discipline.

\textbf{Anti-smuggling guard.} The capsule must not be read as assuming RH, as denying the explicit formula, or as banning ordinary analytic counterexample syntax.  Its only exclusion is the role-confused use of an object at a stronger report level than its carrier and bridge license.  This is why the proof repeatedly separates coordinate, support, endpoint-counterforce, and endpoint-outcome roles.

\textbf{Hostile-referee stress test.} A hostile referee will try to collapse this capsule into either an analytic estimate claim or a tautological role claim.  The response is that the capsule makes no estimate claim and no tautological zero-location claim.  It fixes the role boundary needed before the later discriminator theorem can be applied.

\textbf{Denial burden.} To defeat this capsule, the objection must show an anchor change that is still same-domain.  A mere statement that the method is not analytic-native does not touch the capsule, because the proof class has already been locked as endpoint-rigidity rather than analytic estimation.

\subsection{Capsule 14: ATS tensor lock}
\textbf{Load-bearing role.} This capsule controls the normalized exponent beta-1/2 is tensor content, not skin.  It is used because the proof would otherwise be vulnerable to the hostile reading that an analytic datum, route choice, or proof annotation is silently doing endpoint work before its status has been fixed.

\textbf{Dependency chain.} The local dependency is ATS local packet and amplitude exactness.  No conclusion about the universal location of zeros is available at this stage.  The capsule only fixes the status of the object under its declared carrier and route.  If the object remains descriptive or support-level, it remains lawful.  If it becomes endpoint-standing, it enters the endpoint-status discipline.

\textbf{Anti-smuggling guard.} The capsule must not be read as assuming RH, as denying the explicit formula, or as banning ordinary analytic counterexample syntax.  Its only exclusion is the role-confused use of an object at a stronger report level than its carrier and bridge license.  This is why the proof repeatedly separates coordinate, support, endpoint-counterforce, and endpoint-outcome roles.

\textbf{Hostile-referee stress test.} A hostile referee will try to collapse this capsule into either an analytic estimate claim or a tautological role claim.  The response is that the capsule makes no estimate claim and no tautological zero-location claim.  It fixes the role boundary needed before the later discriminator theorem can be applied.

\textbf{Denial burden.} To defeat this capsule, the objection must show the exponent can be removed by notation while preserving endpoint role.  A mere statement that the method is not analytic-native does not touch the capsule, because the proof class has already been locked as endpoint-rigidity rather than analytic estimation.

\subsection{Capsule 15: ATS skin lock}
\textbf{Load-bearing role.} This capsule controls notation, smoothing, and representative formula choices cannot do endpoint work.  It is used because the proof would otherwise be vulnerable to the hostile reading that an analytic datum, route choice, or proof annotation is silently doing endpoint work before its status has been fixed.

\textbf{Dependency chain.} The local dependency is ATS no-skin-promotion theorem.  No conclusion about the universal location of zeros is available at this stage.  The capsule only fixes the status of the object under its declared carrier and route.  If the object remains descriptive or support-level, it remains lawful.  If it becomes endpoint-standing, it enters the endpoint-status discipline.

\textbf{Anti-smuggling guard.} The capsule must not be read as assuming RH, as denying the explicit formula, or as banning ordinary analytic counterexample syntax.  Its only exclusion is the role-confused use of an object at a stronger report level than its carrier and bridge license.  This is why the proof repeatedly separates coordinate, support, endpoint-counterforce, and endpoint-outcome roles.

\textbf{Hostile-referee stress test.} A hostile referee will try to collapse this capsule into either an analytic estimate claim or a tautological role claim.  The response is that the capsule makes no estimate claim and no tautological zero-location claim.  It fixes the role boundary needed before the later discriminator theorem can be applied.

\textbf{Denial burden.} To defeat this capsule, the objection must turn skin into endpoint status without becoming a surrogate.  A mere statement that the method is not analytic-native does not touch the capsule, because the proof class has already been locked as endpoint-rigidity rather than analytic estimation.

\subsection{Capsule 16: UEAP no-overreporting}
\textbf{Load-bearing role.} This capsule controls zerohood or amplitude support cannot be reported as endpoint outcome without bridge.  It is used because the proof would otherwise be vulnerable to the hostile reading that an analytic datum, route choice, or proof annotation is silently doing endpoint work before its status has been fixed.

\textbf{Dependency chain.} The local dependency is UEAP report ladder.  No conclusion about the universal location of zeros is available at this stage.  The capsule only fixes the status of the object under its declared carrier and route.  If the object remains descriptive or support-level, it remains lawful.  If it becomes endpoint-standing, it enters the endpoint-status discipline.

\textbf{Anti-smuggling guard.} The capsule must not be read as assuming RH, as denying the explicit formula, or as banning ordinary analytic counterexample syntax.  Its only exclusion is the role-confused use of an object at a stronger report level than its carrier and bridge license.  This is why the proof repeatedly separates coordinate, support, endpoint-counterforce, and endpoint-outcome roles.

\textbf{Hostile-referee stress test.} A hostile referee will try to collapse this capsule into either an analytic estimate claim or a tautological role claim.  The response is that the capsule makes no estimate claim and no tautological zero-location claim.  It fixes the role boundary needed before the later discriminator theorem can be applied.

\textbf{Denial burden.} To defeat this capsule, the objection must show report promotion that preserves report discipline.  A mere statement that the method is not analytic-native does not touch the capsule, because the proof class has already been locked as endpoint-rigidity rather than analytic estimation.

\subsection{Capsule 17: No carrier transfer}
\textbf{Load-bearing role.} This capsule controls spectral, numerical, positivity, or random matrix carriers cannot replace the zeta/prime trace route.  It is used because the proof would otherwise be vulnerable to the hostile reading that an analytic datum, route choice, or proof annotation is silently doing endpoint work before its status has been fixed.

\textbf{Dependency chain.} The local dependency is no carriers and silent-redescription closure.  No conclusion about the universal location of zeros is available at this stage.  The capsule only fixes the status of the object under its declared carrier and route.  If the object remains descriptive or support-level, it remains lawful.  If it becomes endpoint-standing, it enters the endpoint-status discipline.

\textbf{Anti-smuggling guard.} The capsule must not be read as assuming RH, as denying the explicit formula, or as banning ordinary analytic counterexample syntax.  Its only exclusion is the role-confused use of an object at a stronger report level than its carrier and bridge license.  This is why the proof repeatedly separates coordinate, support, endpoint-counterforce, and endpoint-outcome roles.

\textbf{Hostile-referee stress test.} A hostile referee will try to collapse this capsule into either an analytic estimate claim or a tautological role claim.  The response is that the capsule makes no estimate claim and no tautological zero-location claim.  It fixes the role boundary needed before the later discriminator theorem can be applied.

\textbf{Denial burden.} To defeat this capsule, the objection must supply a carrier transfer that preserves same endpoint identity.  A mere statement that the method is not analytic-native does not touch the capsule, because the proof class has already been locked as endpoint-rigidity rather than analytic estimation.

\subsection{Capsule 18: No repair}
\textbf{Load-bearing role.} This capsule controls later analytic proof does not repair an earlier unbridged endpoint act as the same act.  It is used because the proof would otherwise be vulnerable to the hostile reading that an analytic datum, route choice, or proof annotation is silently doing endpoint work before its status has been fixed.

\textbf{Dependency chain.} The local dependency is irreversibility and boundary-trace fixation.  No conclusion about the universal location of zeros is available at this stage.  The capsule only fixes the status of the object under its declared carrier and route.  If the object remains descriptive or support-level, it remains lawful.  If it becomes endpoint-standing, it enters the endpoint-status discipline.

\textbf{Anti-smuggling guard.} The capsule must not be read as assuming RH, as denying the explicit formula, or as banning ordinary analytic counterexample syntax.  Its only exclusion is the role-confused use of an object at a stronger report level than its carrier and bridge license.  This is why the proof repeatedly separates coordinate, support, endpoint-counterforce, and endpoint-outcome roles.

\textbf{Hostile-referee stress test.} A hostile referee will try to collapse this capsule into either an analytic estimate claim or a tautological role claim.  The response is that the capsule makes no estimate claim and no tautological zero-location claim.  It fixes the role boundary needed before the later discriminator theorem can be applied.

\textbf{Denial burden.} To defeat this capsule, the objection must show same-act repair compatible with act-time finality.  A mere statement that the method is not analytic-native does not touch the capsule, because the proof class has already been locked as endpoint-rigidity rather than analytic estimation.

\subsection{Capsule 19: Pairing support}
\textbf{Load-bearing role.} This capsule controls symmetry-paired off-line terms are lawful but not endpoint-neutral without a bridge.  It is used because the proof would otherwise be vulnerable to the hostile reading that an analytic datum, route choice, or proof annotation is silently doing endpoint work before its status has been fixed.

\textbf{Dependency chain.} The local dependency is pairing closure theorem.  No conclusion about the universal location of zeros is available at this stage.  The capsule only fixes the status of the object under its declared carrier and route.  If the object remains descriptive or support-level, it remains lawful.  If it becomes endpoint-standing, it enters the endpoint-status discipline.

\textbf{Anti-smuggling guard.} The capsule must not be read as assuming RH, as denying the explicit formula, or as banning ordinary analytic counterexample syntax.  Its only exclusion is the role-confused use of an object at a stronger report level than its carrier and bridge license.  This is why the proof repeatedly separates coordinate, support, endpoint-counterforce, and endpoint-outcome roles.

\textbf{Hostile-referee stress test.} A hostile referee will try to collapse this capsule into either an analytic estimate claim or a tautological role claim.  The response is that the capsule makes no estimate claim and no tautological zero-location claim.  It fixes the role boundary needed before the later discriminator theorem can be applied.

\textbf{Denial burden.} To defeat this capsule, the objection must show paired endpoint counterforce avoiding discriminator classification.  A mere statement that the method is not analytic-native does not touch the capsule, because the proof class has already been locked as endpoint-rigidity rather than analytic estimation.

\subsection{Capsule 20: Cancellation support}
\textbf{Load-bearing role.} This capsule controls aggregate cancellation is support unless certified as endpoint-neutral occupation.  It is used because the proof would otherwise be vulnerable to the hostile reading that an analytic datum, route choice, or proof annotation is silently doing endpoint work before its status has been fixed.

\textbf{Dependency chain.} The local dependency is term isolation and role occupation.  No conclusion about the universal location of zeros is available at this stage.  The capsule only fixes the status of the object under its declared carrier and route.  If the object remains descriptive or support-level, it remains lawful.  If it becomes endpoint-standing, it enters the endpoint-status discipline.

\textbf{Anti-smuggling guard.} The capsule must not be read as assuming RH, as denying the explicit formula, or as banning ordinary analytic counterexample syntax.  Its only exclusion is the role-confused use of an object at a stronger report level than its carrier and bridge license.  This is why the proof repeatedly separates coordinate, support, endpoint-counterforce, and endpoint-outcome roles.

\textbf{Hostile-referee stress test.} A hostile referee will try to collapse this capsule into either an analytic estimate claim or a tautological role claim.  The response is that the capsule makes no estimate claim and no tautological zero-location claim.  It fixes the role boundary needed before the later discriminator theorem can be applied.

\textbf{Denial burden.} To defeat this capsule, the objection must show cancellation that is endpoint-standing but not a bridge and not discriminator work.  A mere statement that the method is not analytic-native does not touch the capsule, because the proof class has already been locked as endpoint-rigidity rather than analytic estimation.

\subsection{Capsule 21: Single amplitude interface}
\textbf{Load-bearing role.} This capsule controls the positive endpoint image has one occupation interface: neutral amplitude.  It is used because the proof would otherwise be vulnerable to the hostile reading that an analytic datum, route choice, or proof annotation is silently doing endpoint work before its status has been fixed.

\textbf{Dependency chain.} The local dependency is official endpoint identity and amplitude exactness.  No conclusion about the universal location of zeros is available at this stage.  The capsule only fixes the status of the object under its declared carrier and route.  If the object remains descriptive or support-level, it remains lawful.  If it becomes endpoint-standing, it enters the endpoint-status discipline.

\textbf{Anti-smuggling guard.} The capsule must not be read as assuming RH, as denying the explicit formula, or as banning ordinary analytic counterexample syntax.  Its only exclusion is the role-confused use of an object at a stronger report level than its carrier and bridge license.  This is why the proof repeatedly separates coordinate, support, endpoint-counterforce, and endpoint-outcome roles.

\textbf{Hostile-referee stress test.} A hostile referee will try to collapse this capsule into either an analytic estimate claim or a tautological role claim.  The response is that the capsule makes no estimate claim and no tautological zero-location claim.  It fixes the role boundary needed before the later discriminator theorem can be applied.

\textbf{Denial burden.} To defeat this capsule, the objection must show lawful non-neutral coequal occupation inside the same image.  A mere statement that the method is not analytic-native does not touch the capsule, because the proof class has already been locked as endpoint-rigidity rather than analytic estimation.

\subsection{Capsule 22: Coequal carrier separation}
\textbf{Load-bearing role.} This capsule controls broad RH versus not-RH outcome is not the same as endpoint-image occupation.  It is used because the proof would otherwise be vulnerable to the hostile reading that an analytic datum, route choice, or proof annotation is silently doing endpoint work before its status has been fixed.

\textbf{Dependency chain.} The local dependency is outcome carrier distinction.  No conclusion about the universal location of zeros is available at this stage.  The capsule only fixes the status of the object under its declared carrier and route.  If the object remains descriptive or support-level, it remains lawful.  If it becomes endpoint-standing, it enters the endpoint-status discipline.

\textbf{Anti-smuggling guard.} The capsule must not be read as assuming RH, as denying the explicit formula, or as banning ordinary analytic counterexample syntax.  Its only exclusion is the role-confused use of an object at a stronger report level than its carrier and bridge license.  This is why the proof repeatedly separates coordinate, support, endpoint-counterforce, and endpoint-outcome roles.

\textbf{Hostile-referee stress test.} A hostile referee will try to collapse this capsule into either an analytic estimate claim or a tautological role claim.  The response is that the capsule makes no estimate claim and no tautological zero-location claim.  It fixes the role boundary needed before the later discriminator theorem can be applied.

\textbf{Denial burden.} To defeat this capsule, the objection must identify a negative branch that is coequal occupation of the neutral amplitude image.  A mere statement that the method is not analytic-native does not touch the capsule, because the proof class has already been locked as endpoint-rigidity rather than analytic estimation.

\subsection{Capsule 23: Non-native amplitude theorem}
\textbf{Load-bearing role.} This capsule controls endpoint-standing non-neutral amplitude is independent endpoint-status work after harmless roles are removed.  It is used because the proof would otherwise be vulnerable to the hostile reading that an analytic datum, route choice, or proof annotation is silently doing endpoint work before its status has been fixed.

\textbf{Dependency chain.} The local dependency is endpoint discriminator trichotomy.  No conclusion about the universal location of zeros is available at this stage.  The capsule only fixes the status of the object under its declared carrier and route.  If the object remains descriptive or support-level, it remains lawful.  If it becomes endpoint-standing, it enters the endpoint-status discipline.

\textbf{Anti-smuggling guard.} The capsule must not be read as assuming RH, as denying the explicit formula, or as banning ordinary analytic counterexample syntax.  Its only exclusion is the role-confused use of an object at a stronger report level than its carrier and bridge license.  This is why the proof repeatedly separates coordinate, support, endpoint-counterforce, and endpoint-outcome roles.

\textbf{Hostile-referee stress test.} A hostile referee will try to collapse this capsule into either an analytic estimate claim or a tautological role claim.  The response is that the capsule makes no estimate claim and no tautological zero-location claim.  It fixes the role boundary needed before the later discriminator theorem can be applied.

\textbf{Denial burden.} To defeat this capsule, the objection must supply the hidden fifth role.  A mere statement that the method is not analytic-native does not touch the capsule, because the proof class has already been locked as endpoint-rigidity rather than analytic estimation.

\subsection{Capsule 24: Amplitude discriminator induction}
\textbf{Load-bearing role.} This capsule controls local separator occupation induces \(D_{\mathrm{amp}}\).  It is used because the proof would otherwise be vulnerable to the hostile reading that an analytic datum, route choice, or proof annotation is silently doing endpoint work before its status has been fixed.

\textbf{Dependency chain.} The local dependency is noncritical status work and trichotomy.  No conclusion about the universal location of zeros is available at this stage.  The capsule only fixes the status of the object under its declared carrier and route.  If the object remains descriptive or support-level, it remains lawful.  If it becomes endpoint-standing, it enters the endpoint-status discipline.

\textbf{Anti-smuggling guard.} The capsule must not be read as assuming RH, as denying the explicit formula, or as banning ordinary analytic counterexample syntax.  Its only exclusion is the role-confused use of an object at a stronger report level than its carrier and bridge license.  This is why the proof repeatedly separates coordinate, support, endpoint-counterforce, and endpoint-outcome roles.

\textbf{Hostile-referee stress test.} A hostile referee will try to collapse this capsule into either an analytic estimate claim or a tautological role claim.  The response is that the capsule makes no estimate claim and no tautological zero-location claim.  It fixes the role boundary needed before the later discriminator theorem can be applied.

\textbf{Denial burden.} To defeat this capsule, the objection must show separator occupation that changes endpoint status without discrimination.  A mere statement that the method is not analytic-native does not touch the capsule, because the proof class has already been locked as endpoint-rigidity rather than analytic estimation.

\subsection{Capsule 25: No-independent closure}
\textbf{Load-bearing role.} This capsule controls endpoint use gives \(N_{\mathrm{amp}}\) and hence not \(D_{\mathrm{amp}}\).  It is used because the proof would otherwise be vulnerable to the hostile reading that an analytic datum, route choice, or proof annotation is silently doing endpoint work before its status has been fixed.

\textbf{Dependency chain.} The local dependency is kernel-forced A+ consequence layer.  No conclusion about the universal location of zeros is available at this stage.  The capsule only fixes the status of the object under its declared carrier and route.  If the object remains descriptive or support-level, it remains lawful.  If it becomes endpoint-standing, it enters the endpoint-status discipline.

\textbf{Anti-smuggling guard.} The capsule must not be read as assuming RH, as denying the explicit formula, or as banning ordinary analytic counterexample syntax.  Its only exclusion is the role-confused use of an object at a stronger report level than its carrier and bridge license.  This is why the proof repeatedly separates coordinate, support, endpoint-counterforce, and endpoint-outcome roles.

\textbf{Hostile-referee stress test.} A hostile referee will try to collapse this capsule into either an analytic estimate claim or a tautological role claim.  The response is that the capsule makes no estimate claim and no tautological zero-location claim.  It fixes the role boundary needed before the later discriminator theorem can be applied.

\textbf{Denial burden.} To defeat this capsule, the objection must break endpoint use or no-independent closure.  A mere statement that the method is not analytic-native does not touch the capsule, because the proof class has already been locked as endpoint-rigidity rather than analytic estimation.

\subsection{Capsule 26: Local reductio role}
\textbf{Load-bearing role.} This capsule controls the exact-complement assumption has local countercase standing only as proof role.  It is used because the proof would otherwise be vulnerable to the hostile reading that an analytic datum, route choice, or proof annotation is silently doing endpoint work before its status has been fixed.

\textbf{Dependency chain.} The local dependency is annotation rule and endpoint closure evaluation.  No conclusion about the universal location of zeros is available at this stage.  The capsule only fixes the status of the object under its declared carrier and route.  If the object remains descriptive or support-level, it remains lawful.  If it becomes endpoint-standing, it enters the endpoint-status discipline.

\textbf{Anti-smuggling guard.} The capsule must not be read as assuming RH, as denying the explicit formula, or as banning ordinary analytic counterexample syntax.  Its only exclusion is the role-confused use of an object at a stronger report level than its carrier and bridge license.  This is why the proof repeatedly separates coordinate, support, endpoint-counterforce, and endpoint-outcome roles.

\textbf{Hostile-referee stress test.} A hostile referee will try to collapse this capsule into either an analytic estimate claim or a tautological role claim.  The response is that the capsule makes no estimate claim and no tautological zero-location claim.  It fixes the role boundary needed before the later discriminator theorem can be applied.

\textbf{Denial burden.} To defeat this capsule, the objection must show the annotation is an object-level premise.  A mere statement that the method is not analytic-native does not touch the capsule, because the proof class has already been locked as endpoint-rigidity rather than analytic estimation.

\subsection{Capsule 27: Annotation discharge}
\textbf{Load-bearing role.} This capsule controls contradiction discharges the original assumption.  It is used because the proof would otherwise be vulnerable to the hostile reading that an analytic datum, route choice, or proof annotation is silently doing endpoint work before its status has been fixed.

\textbf{Dependency chain.} The local dependency is ordinary reductio discipline plus annotation non-premise theorem.  No conclusion about the universal location of zeros is available at this stage.  The capsule only fixes the status of the object under its declared carrier and route.  If the object remains descriptive or support-level, it remains lawful.  If it becomes endpoint-standing, it enters the endpoint-status discipline.

\textbf{Anti-smuggling guard.} The capsule must not be read as assuming RH, as denying the explicit formula, or as banning ordinary analytic counterexample syntax.  Its only exclusion is the role-confused use of an object at a stronger report level than its carrier and bridge license.  This is why the proof repeatedly separates coordinate, support, endpoint-counterforce, and endpoint-outcome roles.

\textbf{Hostile-referee stress test.} A hostile referee will try to collapse this capsule into either an analytic estimate claim or a tautological role claim.  The response is that the capsule makes no estimate claim and no tautological zero-location claim.  It fixes the role boundary needed before the later discriminator theorem can be applied.

\textbf{Denial burden.} To defeat this capsule, the objection must show a strengthened premise was discharged.  A mere statement that the method is not analytic-native does not touch the capsule, because the proof class has already been locked as endpoint-rigidity rather than analytic estimation.

\subsection{Capsule 28: Anti-trivialization}
\textbf{Load-bearing role.} This capsule controls arbitrary negation is not separator occupation.  It is used because the proof would otherwise be vulnerable to the hostile reading that an analytic datum, route choice, or proof annotation is silently doing endpoint work before its status has been fixed.

\textbf{Dependency chain.} The local dependency is route, projection, counterforce, and role-exhaustion prerequisites.  No conclusion about the universal location of zeros is available at this stage.  The capsule only fixes the status of the object under its declared carrier and route.  If the object remains descriptive or support-level, it remains lawful.  If it becomes endpoint-standing, it enters the endpoint-status discipline.

\textbf{Anti-smuggling guard.} The capsule must not be read as assuming RH, as denying the explicit formula, or as banning ordinary analytic counterexample syntax.  Its only exclusion is the role-confused use of an object at a stronger report level than its carrier and bridge license.  This is why the proof repeatedly separates coordinate, support, endpoint-counterforce, and endpoint-outcome roles.

\textbf{Hostile-referee stress test.} A hostile referee will try to collapse this capsule into either an analytic estimate claim or a tautological role claim.  The response is that the capsule makes no estimate claim and no tautological zero-location claim.  It fixes the role boundary needed before the later discriminator theorem can be applied.

\textbf{Denial burden.} To defeat this capsule, the objection must find a theorem that infers separator status from negation alone.  A mere statement that the method is not analytic-native does not touch the capsule, because the proof class has already been locked as endpoint-rigidity rather than analytic estimation.

\subsection{Capsule 29: Ordinary negative theoremhood}
\textbf{Load-bearing role.} This capsule controls negative theoremhood is not rejected for polarity.  It is used because the proof would otherwise be vulnerable to the hostile reading that an analytic datum, route choice, or proof annotation is silently doing endpoint work before its status has been fixed.

\textbf{Dependency chain.} The local dependency is ordinary theoremhood versus endpoint use distinction.  No conclusion about the universal location of zeros is available at this stage.  The capsule only fixes the status of the object under its declared carrier and route.  If the object remains descriptive or support-level, it remains lawful.  If it becomes endpoint-standing, it enters the endpoint-status discipline.

\textbf{Anti-smuggling guard.} The capsule must not be read as assuming RH, as denying the explicit formula, or as banning ordinary analytic counterexample syntax.  Its only exclusion is the role-confused use of an object at a stronger report level than its carrier and bridge license.  This is why the proof repeatedly separates coordinate, support, endpoint-counterforce, and endpoint-outcome roles.

\textbf{Hostile-referee stress test.} A hostile referee will try to collapse this capsule into either an analytic estimate claim or a tautological role claim.  The response is that the capsule makes no estimate claim and no tautological zero-location claim.  It fixes the role boundary needed before the later discriminator theorem can be applied.

\textbf{Denial burden.} To defeat this capsule, the objection must show the paper bans negative mathematics as such.  A mere statement that the method is not analytic-native does not touch the capsule, because the proof class has already been locked as endpoint-rigidity rather than analytic estimation.

\subsection{Capsule 30: Incompleteness firewall}
\textbf{Load-bearing role.} This capsule controls independence and true-but-unprovable distinctions are endpoint-inert unless they change endpoint standing.  It is used because the proof would otherwise be vulnerable to the hostile reading that an analytic datum, route choice, or proof annotation is silently doing endpoint work before its status has been fixed.

\textbf{Dependency chain.} The local dependency is same-domain incompleteness route exhaustion.  No conclusion about the universal location of zeros is available at this stage.  The capsule only fixes the status of the object under its declared carrier and route.  If the object remains descriptive or support-level, it remains lawful.  If it becomes endpoint-standing, it enters the endpoint-status discipline.

\textbf{Anti-smuggling guard.} The capsule must not be read as assuming RH, as denying the explicit formula, or as banning ordinary analytic counterexample syntax.  Its only exclusion is the role-confused use of an object at a stronger report level than its carrier and bridge license.  This is why the proof repeatedly separates coordinate, support, endpoint-counterforce, and endpoint-outcome roles.

\textbf{Hostile-referee stress test.} A hostile referee will try to collapse this capsule into either an analytic estimate claim or a tautological role claim.  The response is that the capsule makes no estimate claim and no tautological zero-location claim.  It fixes the role boundary needed before the later discriminator theorem can be applied.

\textbf{Denial burden.} To defeat this capsule, the objection must produce a third endpoint status that changes standing without becoming gate work.  A mere statement that the method is not analytic-native does not touch the capsule, because the proof class has already been locked as endpoint-rigidity rather than analytic estimation.

\subsection{Capsule 31: Open status separation}
\textbf{Load-bearing role.} This capsule controls open or unresolved status is epistemic and not a third endpoint branch.  It is used because the proof would otherwise be vulnerable to the hostile reading that an analytic datum, route choice, or proof annotation is silently doing endpoint work before its status has been fixed.

\textbf{Dependency chain.} The local dependency is endpoint outcome definition.  No conclusion about the universal location of zeros is available at this stage.  The capsule only fixes the status of the object under its declared carrier and route.  If the object remains descriptive or support-level, it remains lawful.  If it becomes endpoint-standing, it enters the endpoint-status discipline.

\textbf{Anti-smuggling guard.} The capsule must not be read as assuming RH, as denying the explicit formula, or as banning ordinary analytic counterexample syntax.  Its only exclusion is the role-confused use of an object at a stronger report level than its carrier and bridge license.  This is why the proof repeatedly separates coordinate, support, endpoint-counterforce, and endpoint-outcome roles.

\textbf{Hostile-referee stress test.} A hostile referee will try to collapse this capsule into either an analytic estimate claim or a tautological role claim.  The response is that the capsule makes no estimate claim and no tautological zero-location claim.  It fixes the role boundary needed before the later discriminator theorem can be applied.

\textbf{Denial burden.} To defeat this capsule, the objection must show unresolved status can occupy endpoint truth.  A mere statement that the method is not analytic-native does not touch the capsule, because the proof class has already been locked as endpoint-rigidity rather than analytic estimation.

\subsection{Capsule 32: Route-locus exhaustion}
\textbf{Load-bearing role.} This capsule controls every negative route is support, bookkeeping, carrier shift, repair, surrogate, bridge, or forbidden discriminator.  It is used because the proof would otherwise be vulnerable to the hostile reading that an analytic datum, route choice, or proof annotation is silently doing endpoint work before its status has been fixed.

\textbf{Dependency chain.} The local dependency is route-locus theorem.  No conclusion about the universal location of zeros is available at this stage.  The capsule only fixes the status of the object under its declared carrier and route.  If the object remains descriptive or support-level, it remains lawful.  If it becomes endpoint-standing, it enters the endpoint-status discipline.

\textbf{Anti-smuggling guard.} The capsule must not be read as assuming RH, as denying the explicit formula, or as banning ordinary analytic counterexample syntax.  Its only exclusion is the role-confused use of an object at a stronger report level than its carrier and bridge license.  This is why the proof repeatedly separates coordinate, support, endpoint-counterforce, and endpoint-outcome roles.

\textbf{Hostile-referee stress test.} A hostile referee will try to collapse this capsule into either an analytic estimate claim or a tautological role claim.  The response is that the capsule makes no estimate claim and no tautological zero-location claim.  It fixes the role boundary needed before the later discriminator theorem can be applied.

\textbf{Denial burden.} To defeat this capsule, the objection must inhabit a route outside the finite list.  A mere statement that the method is not analytic-native does not touch the capsule, because the proof class has already been locked as endpoint-rigidity rather than analytic estimation.

\subsection{Capsule 33: Weakening resistance}
\textbf{Load-bearing role.} This capsule controls strict weakening either preserves endpoint behavior or permits forbidden governance.  It is used because the proof would otherwise be vulnerable to the hostile reading that an analytic datum, route choice, or proof annotation is silently doing endpoint work before its status has been fixed.

\textbf{Dependency chain.} The local dependency is weakening resistance theorem.  No conclusion about the universal location of zeros is available at this stage.  The capsule only fixes the status of the object under its declared carrier and route.  If the object remains descriptive or support-level, it remains lawful.  If it becomes endpoint-standing, it enters the endpoint-status discipline.

\textbf{Anti-smuggling guard.} The capsule must not be read as assuming RH, as denying the explicit formula, or as banning ordinary analytic counterexample syntax.  Its only exclusion is the role-confused use of an object at a stronger report level than its carrier and bridge license.  This is why the proof repeatedly separates coordinate, support, endpoint-counterforce, and endpoint-outcome roles.

\textbf{Hostile-referee stress test.} A hostile referee will try to collapse this capsule into either an analytic estimate claim or a tautological role claim.  The response is that the capsule makes no estimate claim and no tautological zero-location claim.  It fixes the role boundary needed before the later discriminator theorem can be applied.

\textbf{Denial burden.} To defeat this capsule, the objection must define a weaker same-carrier regime with endpoint separator governance but no second authority.  A mere statement that the method is not analytic-native does not touch the capsule, because the proof class has already been locked as endpoint-rigidity rather than analytic estimation.

\subsection{Capsule 34: Corpus support boundary}
\textbf{Load-bearing role.} This capsule controls corpus rows supply structural constraints, not analytic zero-location facts.  It is used because the proof would otherwise be vulnerable to the hostile reading that an analytic datum, route choice, or proof annotation is silently doing endpoint work before its status has been fixed.

\textbf{Dependency chain.} The local dependency is source-use audit.  No conclusion about the universal location of zeros is available at this stage.  The capsule only fixes the status of the object under its declared carrier and route.  If the object remains descriptive or support-level, it remains lawful.  If it becomes endpoint-standing, it enters the endpoint-status discipline.

\textbf{Anti-smuggling guard.} The capsule must not be read as assuming RH, as denying the explicit formula, or as banning ordinary analytic counterexample syntax.  Its only exclusion is the role-confused use of an object at a stronger report level than its carrier and bridge license.  This is why the proof repeatedly separates coordinate, support, endpoint-counterforce, and endpoint-outcome roles.

\textbf{Hostile-referee stress test.} A hostile referee will try to collapse this capsule into either an analytic estimate claim or a tautological role claim.  The response is that the capsule makes no estimate claim and no tautological zero-location claim.  It fixes the role boundary needed before the later discriminator theorem can be applied.

\textbf{Denial burden.} To defeat this capsule, the objection must show an AASC citation smuggling analytic content.  A mere statement that the method is not analytic-native does not touch the capsule, because the proof class has already been locked as endpoint-rigidity rather than analytic estimation.

\subsection{Capsule 35: Reference integrity}
\textbf{Load-bearing role.} This capsule controls classical references fix analytic objects while AASC references fix structural constraints.  It is used because the proof would otherwise be vulnerable to the hostile reading that an analytic datum, route choice, or proof annotation is silently doing endpoint work before its status has been fixed.

\textbf{Dependency chain.} The local dependency is reference integrity appendix.  No conclusion about the universal location of zeros is available at this stage.  The capsule only fixes the status of the object under its declared carrier and route.  If the object remains descriptive or support-level, it remains lawful.  If it becomes endpoint-standing, it enters the endpoint-status discipline.

\textbf{Anti-smuggling guard.} The capsule must not be read as assuming RH, as denying the explicit formula, or as banning ordinary analytic counterexample syntax.  Its only exclusion is the role-confused use of an object at a stronger report level than its carrier and bridge license.  This is why the proof repeatedly separates coordinate, support, endpoint-counterforce, and endpoint-outcome roles.

\textbf{Hostile-referee stress test.} A hostile referee will try to collapse this capsule into either an analytic estimate claim or a tautological role claim.  The response is that the capsule makes no estimate claim and no tautological zero-location claim.  It fixes the role boundary needed before the later discriminator theorem can be applied.

\textbf{Denial burden.} To defeat this capsule, the objection must show cross-source role confusion.  A mere statement that the method is not analytic-native does not touch the capsule, because the proof class has already been locked as endpoint-rigidity rather than analytic estimation.

\clearpage
\section{Hostile-Referee Simulation and Replies}\label{app:hostile-sim}
The following simulations are phrased in the strongest hostile form.  The replies do not ask the referee to accept vocabulary by preference; each reply locates the exact role distinction the objection must break.

\subsection{Objection 1: This still does not prove a zero-free region}
\textbf{Hostile form.} This still does not prove a zero-free region.

\textbf{Reply.} The proof class is not zero-free-region construction.  The explicit formula is used as endpoint route.  The route does not supply contradiction until non-neutral amplitude is used as endpoint-standing counterforce.

\textbf{Audit location.} The relevant manuscript controls are the proof-class lock, context non-smuggling theorem, EF model adequacy packet, projection-role-content theorem, non-explosiveness theorem, non-native amplitude theorem, local reductio annotation rule, and route-locus exhaustion table.

\textbf{Remaining denial burden.} The objection succeeds only if it produces a same-carrier endpoint role that avoids support, bookkeeping, carrier shift, repair, surrogate substitution, bridge completion, coequal image occupation, and independent discriminator classification.

\subsection{Objection 2: Off-line zeros are ordinary analytic objects}
\textbf{Hostile form.} Off-line zeros are ordinary analytic objects.

\textbf{Reply.} They are preserved as raw syntax and descriptive coordinate reports.  The excluded object is official endpoint-standing use of their non-neutral amplitude projection.

\textbf{Audit location.} The relevant manuscript controls are the proof-class lock, context non-smuggling theorem, EF model adequacy packet, projection-role-content theorem, non-explosiveness theorem, non-native amplitude theorem, local reductio annotation rule, and route-locus exhaustion table.

\textbf{Remaining denial burden.} The objection succeeds only if it produces a same-carrier endpoint role that avoids support, bookkeeping, carrier shift, repair, surrogate substitution, bridge completion, coequal image occupation, and independent discriminator classification.

\subsection{Objection 3: The amplitude exponent is native analytic data}
\textbf{Hostile form.} The amplitude exponent is native analytic data.

\textbf{Reply.} Yes at support level.  It becomes AASC-relevant only when used as boundary authority over the critical-neutral endpoint image.

\textbf{Audit location.} The relevant manuscript controls are the proof-class lock, context non-smuggling theorem, EF model adequacy packet, projection-role-content theorem, non-explosiveness theorem, non-native amplitude theorem, local reductio annotation rule, and route-locus exhaustion table.

\textbf{Remaining denial burden.} The objection succeeds only if it produces a same-carrier endpoint role that avoids support, bookkeeping, carrier shift, repair, surrogate substitution, bridge completion, coequal image occupation, and independent discriminator classification.

\subsection{Objection 4: The proof assumes the critical scale}
\textbf{Hostile form.} The proof assumes the critical scale.

\textbf{Reply.} The critical scale names the endpoint image induced by the official RH statement.  Naming a target image is not asserting universal occupation.

\textbf{Audit location.} The relevant manuscript controls are the proof-class lock, context non-smuggling theorem, EF model adequacy packet, projection-role-content theorem, non-explosiveness theorem, non-native amplitude theorem, local reductio annotation rule, and route-locus exhaustion table.

\textbf{Remaining denial burden.} The objection succeeds only if it produces a same-carrier endpoint role that avoids support, bookkeeping, carrier shift, repair, surrogate substitution, bridge completion, coequal image occupation, and independent discriminator classification.

\subsection{Objection 5: Pairing could cancel the term}
\textbf{Hostile form.} Pairing could cancel the term.

\textbf{Reply.} Pairing and cancellation are support-level unless bridged to endpoint-neutral occupation.  Without such a bridge, they do not supply a hidden fifth endpoint role.

\textbf{Audit location.} The relevant manuscript controls are the proof-class lock, context non-smuggling theorem, EF model adequacy packet, projection-role-content theorem, non-explosiveness theorem, non-native amplitude theorem, local reductio annotation rule, and route-locus exhaustion table.

\textbf{Remaining denial burden.} The objection succeeds only if it produces a same-carrier endpoint role that avoids support, bookkeeping, carrier shift, repair, surrogate substitution, bridge completion, coequal image occupation, and independent discriminator classification.

\subsection{Objection 6: The EF route is just an equivalent RH reformulation}
\textbf{Hostile form.} The EF route is just an equivalent RH reformulation.

\textbf{Reply.} The route is a projection discipline.  The contradiction is not obtained from equivalence alone but from endpoint-standing non-neutral amplitude plus A+ closure.

\textbf{Audit location.} The relevant manuscript controls are the proof-class lock, context non-smuggling theorem, EF model adequacy packet, projection-role-content theorem, non-explosiveness theorem, non-native amplitude theorem, local reductio annotation rule, and route-locus exhaustion table.

\textbf{Remaining denial burden.} The objection succeeds only if it produces a same-carrier endpoint role that avoids support, bookkeeping, carrier shift, repair, surrogate substitution, bridge completion, coequal image occupation, and independent discriminator classification.

\subsection{Objection 7: A+ is doing analytic work}
\textbf{Hostile form.} A+ is doing analytic work.

\textbf{Reply.} A+ only excludes the independent discriminator.  The EF route supplies the amplitude projection; A+ does not compute beta.

\textbf{Audit location.} The relevant manuscript controls are the proof-class lock, context non-smuggling theorem, EF model adequacy packet, projection-role-content theorem, non-explosiveness theorem, non-native amplitude theorem, local reductio annotation rule, and route-locus exhaustion table.

\textbf{Remaining denial burden.} The objection succeeds only if it produces a same-carrier endpoint role that avoids support, bookkeeping, carrier shift, repair, surrogate substitution, bridge completion, coequal image occupation, and independent discriminator classification.

\subsection{Objection 8: The argument proves any proposition}
\textbf{Hostile form.} The argument proves any proposition.

\textbf{Reply.} The anti-trivialization theorem blocks this.  Arbitrary negation lacks EF route, amplitude projection, endpoint counterforce, and role exhaustion.

\textbf{Audit location.} The relevant manuscript controls are the proof-class lock, context non-smuggling theorem, EF model adequacy packet, projection-role-content theorem, non-explosiveness theorem, non-native amplitude theorem, local reductio annotation rule, and route-locus exhaustion table.

\textbf{Remaining denial burden.} The objection succeeds only if it produces a same-carrier endpoint role that avoids support, bookkeeping, carrier shift, repair, surrogate substitution, bridge completion, coequal image occupation, and independent discriminator classification.

\subsection{Objection 9: Ordinary negative theoremhood is forbidden}
\textbf{Hostile form.} Ordinary negative theoremhood is forbidden.

\textbf{Reply.} The manuscript grants ordinary negative theoremhood.  Official endpoint resolution is a different role and is subject to endpoint discipline.

\textbf{Audit location.} The relevant manuscript controls are the proof-class lock, context non-smuggling theorem, EF model adequacy packet, projection-role-content theorem, non-explosiveness theorem, non-native amplitude theorem, local reductio annotation rule, and route-locus exhaustion table.

\textbf{Remaining denial burden.} The objection succeeds only if it produces a same-carrier endpoint role that avoids support, bookkeeping, carrier shift, repair, surrogate substitution, bridge completion, coequal image occupation, and independent discriminator classification.

\subsection{Objection 10: Independence could create a third status}
\textbf{Hostile form.} Independence could create a third status.

\textbf{Reply.} Independence is epistemic unless it changes endpoint standing.  If it does, it is gate work and falls under A+.

\textbf{Audit location.} The relevant manuscript controls are the proof-class lock, context non-smuggling theorem, EF model adequacy packet, projection-role-content theorem, non-explosiveness theorem, non-native amplitude theorem, local reductio annotation rule, and route-locus exhaustion table.

\textbf{Remaining denial burden.} The objection succeeds only if it produces a same-carrier endpoint role that avoids support, bookkeeping, carrier shift, repair, surrogate substitution, bridge completion, coequal image occupation, and independent discriminator classification.

\subsection{Objection 11: Numerical verification supports RH}
\textbf{Hostile form.} Numerical verification supports RH.

\textbf{Reply.} Correct, as finite-domain support.  It does not occupy the universal endpoint without a bridge.

\textbf{Audit location.} The relevant manuscript controls are the proof-class lock, context non-smuggling theorem, EF model adequacy packet, projection-role-content theorem, non-explosiveness theorem, non-native amplitude theorem, local reductio annotation rule, and route-locus exhaustion table.

\textbf{Remaining denial burden.} The objection succeeds only if it produces a same-carrier endpoint role that avoids support, bookkeeping, carrier shift, repair, surrogate substitution, bridge completion, coequal image occupation, and independent discriminator classification.

\subsection{Objection 12: A spectral operator may solve RH}
\textbf{Hostile form.} A spectral operator may solve RH.

\textbf{Reply.} Correct, if it supplies a bridge to the same zeta zerohood and line-readout carrier.  Without that bridge it is a surrogate.

\textbf{Audit location.} The relevant manuscript controls are the proof-class lock, context non-smuggling theorem, EF model adequacy packet, projection-role-content theorem, non-explosiveness theorem, non-native amplitude theorem, local reductio annotation rule, and route-locus exhaustion table.

\textbf{Remaining denial burden.} The objection succeeds only if it produces a same-carrier endpoint role that avoids support, bookkeeping, carrier shift, repair, surrogate substitution, bridge completion, coequal image occupation, and independent discriminator classification.

\subsection{Objection 13: A positivity criterion could solve RH}
\textbf{Hostile form.} A positivity criterion could solve RH.

\textbf{Reply.} Correct, if the positivity theorem is established and bridged.  The manuscript rejects only unproved carrier transfer.

\textbf{Audit location.} The relevant manuscript controls are the proof-class lock, context non-smuggling theorem, EF model adequacy packet, projection-role-content theorem, non-explosiveness theorem, non-native amplitude theorem, local reductio annotation rule, and route-locus exhaustion table.

\textbf{Remaining denial burden.} The objection succeeds only if it produces a same-carrier endpoint role that avoids support, bookkeeping, carrier shift, repair, surrogate substitution, bridge completion, coequal image occupation, and independent discriminator classification.

\subsection{Objection 14: Density estimates almost prove RH}
\textbf{Hostile form.} Density estimates almost prove RH.

\textbf{Reply.} Density estimates are support unless the exceptional set is eliminated by a bridge.

\textbf{Audit location.} The relevant manuscript controls are the proof-class lock, context non-smuggling theorem, EF model adequacy packet, projection-role-content theorem, non-explosiveness theorem, non-native amplitude theorem, local reductio annotation rule, and route-locus exhaustion table.

\textbf{Remaining denial burden.} The objection succeeds only if it produces a same-carrier endpoint role that avoids support, bookkeeping, carrier shift, repair, surrogate substitution, bridge completion, coequal image occupation, and independent discriminator classification.

\subsection{Objection 15: The proof is too formal}
\textbf{Hostile form.} The proof is too formal.

\textbf{Reply.} Endpoint-use claims are formal claims about standing, reportability, and fixed-domain reuse.  A same-mode objection must attack those links, not demand a different genre.

\textbf{Audit location.} The relevant manuscript controls are the proof-class lock, context non-smuggling theorem, EF model adequacy packet, projection-role-content theorem, non-explosiveness theorem, non-native amplitude theorem, local reductio annotation rule, and route-locus exhaustion table.

\textbf{Remaining denial burden.} The objection succeeds only if it produces a same-carrier endpoint role that avoids support, bookkeeping, carrier shift, repair, surrogate substitution, bridge completion, coequal image occupation, and independent discriminator classification.

\subsection{Objection 16: The context object contains too much}
\textbf{Hostile form.} The context object contains too much.

\textbf{Reply.} The context non-smuggling theorem lists exactly what is not contained: RH, line support, no-independent closure, and branch exclusion.

\textbf{Audit location.} The relevant manuscript controls are the proof-class lock, context non-smuggling theorem, EF model adequacy packet, projection-role-content theorem, non-explosiveness theorem, non-native amplitude theorem, local reductio annotation rule, and route-locus exhaustion table.

\textbf{Remaining denial burden.} The objection succeeds only if it produces a same-carrier endpoint role that avoids support, bookkeeping, carrier shift, repair, surrogate substitution, bridge completion, coequal image occupation, and independent discriminator classification.

\subsection{Objection 17: Projection-role content is artificial}
\textbf{Hostile form.} Projection-role content is artificial.

\textbf{Reply.} Projection-role content is the formal name for reading a branch through the official route.  Refusing it while accepting the route is carrier drift.

\textbf{Audit location.} The relevant manuscript controls are the proof-class lock, context non-smuggling theorem, EF model adequacy packet, projection-role-content theorem, non-explosiveness theorem, non-native amplitude theorem, local reductio annotation rule, and route-locus exhaustion table.

\textbf{Remaining denial burden.} The objection succeeds only if it produces a same-carrier endpoint role that avoids support, bookkeeping, carrier shift, repair, surrogate substitution, bridge completion, coequal image occupation, and independent discriminator classification.

\subsection{Objection 18: The local annotation adds a premise}
\textbf{Hostile form.} The local annotation adds a premise.

\textbf{Reply.} The annotation is proof-role metadata; the object-level assumption remains only not Lcrit.

\textbf{Audit location.} The relevant manuscript controls are the proof-class lock, context non-smuggling theorem, EF model adequacy packet, projection-role-content theorem, non-explosiveness theorem, non-native amplitude theorem, local reductio annotation rule, and route-locus exhaustion table.

\textbf{Remaining denial burden.} The objection succeeds only if it produces a same-carrier endpoint role that avoids support, bookkeeping, carrier shift, repair, surrogate substitution, bridge completion, coequal image occupation, and independent discriminator classification.

\subsection{Objection 19: Neutral amplitude is arbitrary}
\textbf{Hostile form.} Neutral amplitude is arbitrary.

\textbf{Reply.} It is forced by critical normalization of the official RH endpoint.  Other normalizations are skin unless bridged to the same endpoint image.

\textbf{Audit location.} The relevant manuscript controls are the proof-class lock, context non-smuggling theorem, EF model adequacy packet, projection-role-content theorem, non-explosiveness theorem, non-native amplitude theorem, local reductio annotation rule, and route-locus exhaustion table.

\textbf{Remaining denial burden.} The objection succeeds only if it produces a same-carrier endpoint role that avoids support, bookkeeping, carrier shift, repair, surrogate substitution, bridge completion, coequal image occupation, and independent discriminator classification.

\subsection{Objection 20: The corpus matrix is not an analytic proof}
\textbf{Hostile form.} The corpus matrix is not an analytic proof.

\textbf{Reply.} Correct.  It supplies structural locks only.  Analytic objects are supplied by classical references and definitions.

\textbf{Audit location.} The relevant manuscript controls are the proof-class lock, context non-smuggling theorem, EF model adequacy packet, projection-role-content theorem, non-explosiveness theorem, non-native amplitude theorem, local reductio annotation rule, and route-locus exhaustion table.

\textbf{Remaining denial burden.} The objection succeeds only if it produces a same-carrier endpoint role that avoids support, bookkeeping, carrier shift, repair, surrogate substitution, bridge completion, coequal image occupation, and independent discriminator classification.

\clearpage
\section{Detailed Route-Certificate Audit}\label{app:route-certificates}
Each certificate identifies the exact route by which an unlicensed negative branch would attempt to enter the endpoint.  The certificates are operational versions of the route-locus exhaustion theorem.

\subsection{Certificate 1: Status coordinate certificate}
\textbf{Failure pattern.} The route occurs when zerohood or off-line amplitude is reported at a stronger status than its carrier permits.

\textbf{Control theorem.} UEAP blocks the promotion unless endpoint counterforce and bridge conditions are present.

\textbf{RH-EF application.} In the amplitude manuscript the same certificate asks whether a statement about \(\rho\), \(\beta\), \(x^{\beta-1/2}\), a symmetry family, a finite height interval, or an auxiliary carrier is being used as endpoint-standing content.  If it is not, it remains lawful support.  If it is, the route must pass through the endpoint-discriminator trichotomy.

\textbf{Reviewer test.} A reviewer can test the route by asking: does this object change endpoint standing, reportability, branch exclusion, or downstream reuse on the same carrier?  If no, it is endpoint-inert support.  If yes, it is endpoint-admissibility-relevant and enters the no-independent-discriminator closure.

\subsection{Certificate 2: Readout-slot certificate}
\textbf{Failure pattern.} The route occurs when the target readout is replaced by spectral, numerical, density, or positivity output.

\textbf{Control theorem.} silent redescription blocks the substitution unless a bridge identifies the output with same-zero line readout.

\textbf{RH-EF application.} In the amplitude manuscript the same certificate asks whether a statement about \(\rho\), \(\beta\), \(x^{\beta-1/2}\), a symmetry family, a finite height interval, or an auxiliary carrier is being used as endpoint-standing content.  If it is not, it remains lawful support.  If it is, the route must pass through the endpoint-discriminator trichotomy.

\textbf{Reviewer test.} A reviewer can test the route by asking: does this object change endpoint standing, reportability, branch exclusion, or downstream reuse on the same carrier?  If no, it is endpoint-inert support.  If yes, it is endpoint-admissibility-relevant and enters the no-independent-discriminator closure.

\subsection{Certificate 3: Occupation certificate}
\textbf{Failure pattern.} The route occurs when neutral amplitude is declared occupied without local countercase exclusion or bridge completion.

\textbf{Control theorem.} no-generator discipline blocks standing creation by declaration.

\textbf{RH-EF application.} In the amplitude manuscript the same certificate asks whether a statement about \(\rho\), \(\beta\), \(x^{\beta-1/2}\), a symmetry family, a finite height interval, or an auxiliary carrier is being used as endpoint-standing content.  If it is not, it remains lawful support.  If it is, the route must pass through the endpoint-discriminator trichotomy.

\textbf{Reviewer test.} A reviewer can test the route by asking: does this object change endpoint standing, reportability, branch exclusion, or downstream reuse on the same carrier?  If no, it is endpoint-inert support.  If yes, it is endpoint-admissibility-relevant and enters the no-independent-discriminator closure.

\subsection{Certificate 4: Carrier-transfer certificate}
\textbf{Failure pattern.} The route occurs when auxiliary carrier standing is imported into the zeta/prime-trace carrier.

\textbf{Control theorem.} no-carrier discipline blocks the transfer unless a lawful bridge is supplied.

\textbf{RH-EF application.} In the amplitude manuscript the same certificate asks whether a statement about \(\rho\), \(\beta\), \(x^{\beta-1/2}\), a symmetry family, a finite height interval, or an auxiliary carrier is being used as endpoint-standing content.  If it is not, it remains lawful support.  If it is, the route must pass through the endpoint-discriminator trichotomy.

\textbf{Reviewer test.} A reviewer can test the route by asking: does this object change endpoint standing, reportability, branch exclusion, or downstream reuse on the same carrier?  If no, it is endpoint-inert support.  If yes, it is endpoint-admissibility-relevant and enters the no-independent-discriminator closure.

\subsection{Certificate 5: Temporal certificate}
\textbf{Failure pattern.} The route occurs when later analytic work is used to rescue an earlier unbridged endpoint act as the same act.

\textbf{Control theorem.} irreversibility and no post-selection repair force a new act.

\textbf{RH-EF application.} In the amplitude manuscript the same certificate asks whether a statement about \(\rho\), \(\beta\), \(x^{\beta-1/2}\), a symmetry family, a finite height interval, or an auxiliary carrier is being used as endpoint-standing content.  If it is not, it remains lawful support.  If it is, the route must pass through the endpoint-discriminator trichotomy.

\textbf{Reviewer test.} A reviewer can test the route by asking: does this object change endpoint standing, reportability, branch exclusion, or downstream reuse on the same carrier?  If no, it is endpoint-inert support.  If yes, it is endpoint-admissibility-relevant and enters the no-independent-discriminator closure.

\subsection{Certificate 6: Locality certificate}
\textbf{Failure pattern.} The route occurs when finite or local facts are treated as universal endpoint standing.

\textbf{Control theorem.} nonfactorizability demands a global compatibility certificate.

\textbf{RH-EF application.} In the amplitude manuscript the same certificate asks whether a statement about \(\rho\), \(\beta\), \(x^{\beta-1/2}\), a symmetry family, a finite height interval, or an auxiliary carrier is being used as endpoint-standing content.  If it is not, it remains lawful support.  If it is, the route must pass through the endpoint-discriminator trichotomy.

\textbf{Reviewer test.} A reviewer can test the route by asking: does this object change endpoint standing, reportability, branch exclusion, or downstream reuse on the same carrier?  If no, it is endpoint-inert support.  If yes, it is endpoint-admissibility-relevant and enters the no-independent-discriminator closure.

\subsection{Certificate 7: Selection certificate}
\textbf{Failure pattern.} The route occurs when a scale, coordinate, or representative is used as boundary authority.

\textbf{Control theorem.} AMetric no-selector discipline blocks endpoint-standing selector use.

\textbf{RH-EF application.} In the amplitude manuscript the same certificate asks whether a statement about \(\rho\), \(\beta\), \(x^{\beta-1/2}\), a symmetry family, a finite height interval, or an auxiliary carrier is being used as endpoint-standing content.  If it is not, it remains lawful support.  If it is, the route must pass through the endpoint-discriminator trichotomy.

\textbf{Reviewer test.} A reviewer can test the route by asking: does this object change endpoint standing, reportability, branch exclusion, or downstream reuse on the same carrier?  If no, it is endpoint-inert support.  If yes, it is endpoint-admissibility-relevant and enters the no-independent-discriminator closure.

\subsection{Certificate 8: Pairing certificate}
\textbf{Failure pattern.} The route occurs when symmetry pair or cancellation is used as endpoint-neutral without proof of neutral occupation.

\textbf{Control theorem.} pairing closure classifies it as support unless bridge-complete.

\textbf{RH-EF application.} In the amplitude manuscript the same certificate asks whether a statement about \(\rho\), \(\beta\), \(x^{\beta-1/2}\), a symmetry family, a finite height interval, or an auxiliary carrier is being used as endpoint-standing content.  If it is not, it remains lawful support.  If it is, the route must pass through the endpoint-discriminator trichotomy.

\textbf{Reviewer test.} A reviewer can test the route by asking: does this object change endpoint standing, reportability, branch exclusion, or downstream reuse on the same carrier?  If no, it is endpoint-inert support.  If yes, it is endpoint-admissibility-relevant and enters the no-independent-discriminator closure.

\subsection{Certificate 9: Projection certificate}
\textbf{Failure pattern.} The route occurs when the branch is used through the EF route but its projected role content is denied.

\textbf{Control theorem.} official route discipline blocks accepting endpoint force while refusing projection.

\textbf{RH-EF application.} In the amplitude manuscript the same certificate asks whether a statement about \(\rho\), \(\beta\), \(x^{\beta-1/2}\), a symmetry family, a finite height interval, or an auxiliary carrier is being used as endpoint-standing content.  If it is not, it remains lawful support.  If it is, the route must pass through the endpoint-discriminator trichotomy.

\textbf{Reviewer test.} A reviewer can test the route by asking: does this object change endpoint standing, reportability, branch exclusion, or downstream reuse on the same carrier?  If no, it is endpoint-inert support.  If yes, it is endpoint-admissibility-relevant and enters the no-independent-discriminator closure.

\subsection{Certificate 10: Endpoint-use certificate}
\textbf{Failure pattern.} The route occurs when ordinary theoremhood is used as official endpoint resolution while claiming exemption from A+.

\textbf{Control theorem.} ordinary official resolution is endpoint use and triggers the kernel consequence layer.

\textbf{RH-EF application.} In the amplitude manuscript the same certificate asks whether a statement about \(\rho\), \(\beta\), \(x^{\beta-1/2}\), a symmetry family, a finite height interval, or an auxiliary carrier is being used as endpoint-standing content.  If it is not, it remains lawful support.  If it is, the route must pass through the endpoint-discriminator trichotomy.

\textbf{Reviewer test.} A reviewer can test the route by asking: does this object change endpoint standing, reportability, branch exclusion, or downstream reuse on the same carrier?  If no, it is endpoint-inert support.  If yes, it is endpoint-admissibility-relevant and enters the no-independent-discriminator closure.


\section{Conclusion and Final Same-Mode Denial Fork}

This manuscript has one nonclassical dependency: official endpoint-defeating counterexample force is theorem-bearing endpoint use.  Once that is granted, the remaining proof is target instantiation plus route exhaustion.  The conclusion therefore does not hide the decisive commitment.  A reader who accepts the endpoint-use status of official counterexample force need only check the RH-EF instantiation, the structural-to-explicit-formula equivalence bridge, the native-counterforce exhaustion, the local reductio discipline, and the A+ no-independent-discriminator closure.  A reader who rejects that dependency should attack the no-autonomous-counterexample endpoint-role theorem directly.

The manuscript fixes the official Riemann Hypothesis endpoint, the zeta/prime-trace explicit-formula carrier, the critical-normalized amplitude endpoint image, and the live exact-complement local countercase discipline.  It preserves raw off-line zerohood, descriptive coordinate information, ordinary negative theoremhood, support-level explicit-formula amplitude displacement, symmetry pairing, cancellation analysis, finite numerical evidence, density estimates, positivity programs, and spectral programs in their lawful roles.  None of those is excluded merely as content.

The excluded object is narrower: same-carrier endpoint-standing non-neutral amplitude counterforce that is not bridge-neutralized, not support-only, not bookkeeping, not carrier/domain shift, and not lawful coequal occupation of the critical-neutral endpoint image.  The structural-to-explicit-formula bridge turns the live off-line countercase into non-neutral critical-normalized order.  Native-counterforce exhaustion closes the proposed ordinary-native-falsification fifth role.  The no-autonomous-counterexample theorem closes the remaining attempt to keep official counterexample force outside endpoint-use discipline.

Thus the final same-mode denial fork is exact.  A critic must break one of the following links: official counterexample force is endpoint use; endpoint use forces the kernel and A+ consequence layer; the structural-to-EF bridge maps line readout to neutral order and off-line readout to non-neutral order; native counterforce has no fifth endpoint role; fixed-domain use excludes carrier shift; or A+ excludes the independent amplitude discriminator.  If those links remain standing, the local off-line countercase induces \(\Damp(\rho)\) and the same endpoint-use act yields \(\neg\Damp(\rho)\).  The local assumption discharges, arbitrary zerohood gives critical-line readout, and universal generalization gives RH.

\section{Reference Integrity}

Classical analytic references fix the zeta object, the Riemann Hypothesis statement, and the explicit-formula background.  They do not supply AASC endpoint closure.  AASC corpus references fix endpoint-use governance, ATS locking, UEAP report preservation, AMetric boundary discipline, no repair, no carrier transfer, no silent redescription, nonfactorizability, role occupation, and same-domain incompleteness control.  They do not supply analytic estimates for \(\zeta\).  The proof combines the two only at the declared endpoint-route interface.

\section{References}
\begin{enumerate}
\item B. Riemann, \emph{Ueber die Anzahl der Primzahlen unter einer gegebenen Grosse}, Monatsberichte der Berliner Akademie, 1859.
\item H. M. Edwards, \emph{Riemann's Zeta Function}, Academic Press, 1974.
\item E. C. Titchmarsh, \emph{The Theory of the Riemann Zeta-Function}, 2nd ed., revised by D. R. Heath-Brown, Oxford University Press, 1986.
\item A. Ivic, \emph{The Riemann Zeta-Function}, Dover, 2003.
\item E. Bombieri, \emph{Problems of the Millennium: The Riemann Hypothesis}, Clay Mathematics Institute, 2000.
\item A. J. Maley, \emph{A Structural Solution to the Riemann Hypothesis as a Carrier-Preservation Role-Compression Problem}, 2026.
\item A. J. Maley, \emph{P = NP by Exclusion of the SAT Separator Endpoint}, 2026.
\item A. J. Maley, AASC Corpus Control Matrix, uploaded audit matrix, 2026.
\end{enumerate}

\end{document}
